% !TEX root = ../../tree_notes.tex
% =====================================================================
%  Chapter 29 -- Trees Born from Graphs            WRITTEN
%  (folder ch31_trees_from_graphs; label prefix c31; PRINTS as
%  Chapter 29; folder number 31; build with:  python tools/build.py --only 31)
%
%  NUMBERING TRAP (the Ch.27/28 trap, mirrored): \chref{c31} prints
%  "Chapter 29" (THIS chapter).  Do NOT confuse label with printed number.
%  \chref{c22} prints "Chapter 28" (Tree Isomorphism -- the PREVIOUS chapter),
%  \chref{c21} prints "Chapter 27" (Mo + Tarjan on trees),
%  \chref{c30} prints "Chapter 26" (virtual trees),
%  \chref{c23} prints "Chapter 25" (link-cut trees),
%  \chref{c09} prints "Chapter 12" (tree metrics -- OWNS diameter/centre, the
%  black box the Summit reuses to measure the residual tree),
%  \chref{c24} prints "Chapter 5" (Euler tour / tin-tout),
%  \chref{c10} prints "Chapter 13" (tree DP -- the fold you run ON the bridge
%  tree and the block-cut tree),
%  \chref{c20} prints "Chapter 15" (rerooting),
%  \chref{c18} prints "Chapter 24" (centroid decomposition),
%  \chref{c16} prints "Chapter 22" (HLD), \chref{c11} prints "Chapter 16"
%  (LCA), \chref{c28} prints "Chapter 17" (DSU / small-to-large -- the
%  component-merge engine behind the bridge-tree contraction),
%  \chref{c32} prints "Chapter 30" (Counting Labelled Trees -- the NEXT-NEXT
%  chapter), \chref{c17} prints "Chapter 31" (MST -- the chapter AFTER next).
%  Every cite of THIS chapter is \chref{c31}.
%
%  Owns THE TREE HIDING INSIDE A GENERAL GRAPH: when an undirected (or
%  functional) graph is not a tree, which canonical tree still falls out of it,
%  and what you can compute on that tree that you could not compute on the
%  graph.  The tools: Tarjan bridge finding (tin/low, an edge u-v is a bridge
%  iff low[v] > tin[u]), 2-edge-connected components and the BRIDGE TREE
%  (contract each bridge-less component, bridges become tree edges -- the object
%  you DP on), articulation points and the BLOCK-CUT TREE (vertex-biconnected
%  components / blocks via the edge-stack DFS, glued at cut vertices), the
%  FUNCTIONAL GRAPH (a set on out-degree 1: each weak component is a directed
%  cycle with in-trees rooted on it -- the rho shape), and the CACTUS (every
%  edge in at most one cycle: cycles are the blocks, and one contracted node per
%  cycle turns it back into a tree you DP on).  It STARTS where Ch.28
%  (\chref{c22}) ENDS: that chapter assumed the input IS a tree; this one earns
%  trees from graphs that are not.  Diameter/centre (Ch.12), tree DP folds
%  (Ch.13), tin/tout (Ch.5), DSU component merging (Ch.17) are USED as black
%  boxes where earlier chapters own them, never re-derived.
%
%  Sections:
%    * 29.1 bridges and the bridge tree -- Tarjan low-link bridge test, contract
%      the 2-edge-connected components into a tree, DP on it.  Solved CF 118E
%      "Bertown roads" (2000, L4) -- orient every edge so the digraph is strongly
%      connected iff the graph has no bridge (Robbins' theorem); the DFS that
%      finds bridges orients non-bridge edges as traversed and back edges toward
%      the ancestor; print 0 when a bridge exists.  Both official samples
%      hand-replayed (6 8 triangle-pair -> a valid orientation; 6 7 with a bridge
%      -> 0).
%    * 29.2 articulation points, blocks and the block-cut tree -- the edge-stack
%      DFS pops a maximal biconnected component when low[v] >= tin[u]; cut
%      vertices are shared across blocks.  Solved CF 962F "Simple Cycles Edges"
%      (2400, L5) -- report edges lying on EXACTLY ONE simple cycle: an edge is
%      on exactly one simple cycle iff its block has |E| = |V| (a simple cycle);
%      lemma brute-force verified over all graphs on <= 6 vertices (70,646
%      blocks, 0 counterexamples).  Official sample (triangle) -> 3 / "1 2 3".
%    * 29.3 functional graphs: cycles with trees on top -- a set f: V -> V has
%      out-degree 1; each weak component is one directed cycle with in-trees
%      rooted on its vertices (the rho).  Detect with the 0/1/2 colour DFS or
%      in-degree peeling; the in-trees are ordinary rooted trees for Ch.13 DP.
%      WORKED examples (cycle + tails of a concrete f; component census) PLUS a
%      solved lab: CF 702E "Analysis of Pathes in Functional Graph" (2100, L5) --
%      per start vertex, the sum and minimum arc weight over a fixed k-step walk
%      (k up to 1e10); solved by binary lifting on f carrying a (sum, min) monoid,
%      O(n log k); all three official samples hand-replayed.  Promoted from this
%      chapter's mixed set (its URL now lives here only).
%    * 29.4 cactus graphs -- every edge lies in at most one simple cycle; the
%      cycles are the blocks, so contracting each cycle to a node gives a tree,
%      and DP over a cactus = DP over that tree with a cycle-transfer at each
%      cycle node.  WORKED examples (the bowtie independent-set DP -> 10) PLUS a
%      solved lab: CF 231E "Cactus" (2100, L5) -- a vertex cactus; count simple
%      paths between query pairs, which is 2^c with c the number of cycles on the
%      route; solved by DSU-contracting each cycle, then LCA + a root-to-node cycle
%      prefix on the contracted tree, mod 1e9+7; official 6-query sample replayed
%      -> 2 2 2 4 4 1.  Promoted from this chapter's mixed set (URL lives here only).
%    * 29.5 the method boundary and a dominator preview -- when a graph's
%      tree-of-components is the right object and when you are actually being
%      asked an s-t reachability question (2-SAT on the bridge tree, DFS-tree
%      back-edge reasoning); a ONE-PARAGRAPH dominator-tree preview (an edge u-v
%      is a bridge iff it dominates in both DFS directions) with the honest
%      caveat that flow-graph dominators live in graph-land, not here.
%      WORKED examples only.
%  Chapter-end Summit: CF 835F "Roads in the Kingdom" (2500, L5) -- n towns, n
%  roads (a UNICYCLIC connected graph: exactly one cycle), close exactly one road
%  so the rest stays connected (i.e. delete one CYCLE edge) and minimise the
%  DIAMETER of the resulting tree.  Pipeline: find the unique cycle; for each
%  cycle vertex compute the height h[c] and internal diameter dtree[c] of its
%  attached in-tree (Ch.12 two-sweep / Ch.13 fold black boxes); the answer is
%  min over cycle edges of the tree diameter after deleting that edge = the max
%  of (best internal dtree, and the widest path routed through the surviving
%  cycle path); the per-cut window diameter is a sliding-window pair-max over a
%  doubled prefix array, answered with a segment-tree monoid storing
%  (best pair, max(h - S), max(h + S)).  Reuses THREE chapter tools (the cycle
%  structure, contraction, the DP) and TWO earlier-chapter tools (Ch.12
%  diameter/height, Ch.19/20 segment tree).  Official sample hand-replayed -> 5.
%
%  Diagram-first: this chapter is GRAPH based.  The running example is the
%  9-vertex graph G defined once, authoritatively, in 00_opener: a triangle
%  1-2-3-1 (block B1), a bridge 3-4, a second triangle 4-5-6-4 (block B2), and a
%  bridge pendant path 4-7-8-9.  Bridges = {3-4, 4-7, 7-8, 8-9}; cut vertices =
%  {3, 4, 7, 8}; the bridge tree (contract B1, B2, {7}, {8}, {9}) is the path
%  B1 - B2 - 7 - 8 - 9.  The functional-graph section uses its OWN small map
%  f (a rho); the cactus section uses its OWN cactus; the Summit and both solved
%  problems use their OWN official samples, so no trace collides.
%
%  Anti-repetition contract (teach once, cross-reference by label):
%    * Ch.12 tree metrics (\chref{c09}) owns the diameter/height two-sweep; the
%      Summit QUOTES it as a black box on each attached tree, never re-proves it.
%    * Ch.13 tree DP (\chref{c10}) owns the bottom-up fold; the bridge-tree and
%      block-cut-tree DPs are that fold on a NEW tree, one new sentence each.
%    * Ch.17 DSU / small-to-large (\chref{c28}) owns the component merge; 29.1's
%      contraction USES it (or the equivalent DFS labelling), citing not rederive.
%    * Ch.5 Euler tour / tin-tout (\chref{c24}) owns the interval view; the
%      dominator preview gestures at it, no re-derivation.
%    * Ch.28 tree isomorphism (\chref{c22}) is the PREVIOUS chapter and ENDS at
%      "the input is a tree"; 29.1 OPENS at "the input is not -- extract the tree".
%
%  Known URL collisions (confirmed by a FRESH ownership scan of chapters/ + the
%  live CF API; prose-cite only, NEVER re-link -- urlsweep must print an empty
%  DUPLICATES section):
%    * CF 977/E "Cyclic Components" (1700) -> already owned by an earlier DSU/
%      connectivity chapter; named in 29.1 prose WITHOUT a URL.
%    * The functional-graph classic CF 1312/E contacts-style and CF 416/D parent
%      queries are NOT linked here (owned elsewhere or out of band).
%    * The roster URLs below (118E, 962F, 835F) and every mixed-set URL were
%      confirmed ABSENT from build/used_cf.txt (the 261-ID used set) BEFORE
%      writing; re-derive ownership from the chapter files, never a cached list.
%
%  Roster (every title/rating/tag/sample live-verified against the Codeforces
%  problemset API / problem page BEFORE a trace was written; CF-forward):
%    solved: CF 118E Bertown roads            (2000, L4) (29.1) -- orient for
%            strong connectivity iff no bridge; live samples replayed -> a legal
%            orientation, and 0.
%            CF 962F Simple Cycles Edges      (2400, L5) (29.2) -- edges on
%            exactly one simple cycle = edges of |E|=|V| blocks; live sample
%            replayed -> 3 / "1 2 3"; lemma brute-force verified.
%            CF 702E Analysis of Pathes in Functional Graph (2100, L5) (29.3) --
%            sum+min over a fixed k-step orbit; binary lifting on f with a
%            (sum,min) monoid; all three live samples replayed.  Promoted from
%            the mixed set.
%            CF 231E Cactus                   (2100, L5) (29.4) -- count simple
%            x-y paths = 2^(cycles on route); DSU-contract cycles + LCA cycle
%            prefix, mod 1e9+7; live 6-query sample replayed -> 2 2 2 4 4 1.
%            Promoted from the mixed set.
%    worked: bridge/low-link run + bridge-tree contraction on G (29.1);
%            block-cut tree of G and the |E|=|V| test on each block (29.2); the
%            rho decomposition + in-degree peeling of a concrete f (29.3); the
%            cycle-contraction DP on a small cactus (29.4); the boundary pick and
%            the one-paragraph dominator preview (29.5).
%    summit: CF 835F Roads in the Kingdom      (2500, L5) -- unicyclic, delete one
%            cycle edge minimising the residual tree diameter; cycle find +
%            attached heights/diameters + sliding-window pair-max; sample
%            hand-replayed -> 5.
%    mixed (8, ascending by live-verified rating, every ID confirmed ABSENT from
%            build/used_cf.txt -- LOCKED before writing): 920E(2100),
%            1039C(2200), 9E(2300), 732F(2300),
%            1200F(2300), 2231E(2300), 856D(2400), 986C(2500).  (702E and 231E
%            were promoted into the 29.3 / 29.4 solved labs.)  A spread of
%            genuine graph->tree drills (bridge-based network safety 1039C,
%            strong orientation of cycles 732F, bridge/DFS edge game 9E, block
%            cutting 2231E, cactus DP 856D "Masha and Cactus") and deliberate
%            REJECTION cases whose hints name the true tool (complement-graph
%            connectivity 920E -> DSU traversal \chref{c28}; bitmask-defined dense
%            graph 986C -> DFS on an auxiliary graph, no edge is a bridge; DP over
%            a travel order 1200F).
% =====================================================================
\bookchapter{Trees Born from Graphs}{c31}
\chapterinfo{CF 2000--2500}{Written}

% =====================================================================
%  Chapter 29 opener -- "Chapter at a Glance", three reading paths + the
%  ladder staircase, THE running example (the 9-vertex graph G: two triangles
%  B1 = 1-2-3, B2 = 4-5-6 joined by a bridge 3-4, with a bridge pendant path
%  4-7-8-9), the bridge tree B1-B2-7-8-9, the block-cut tree, the two traps
%  (a bridge is not an articulation point; a block is not a 2-edge-component),
%  the who-owns-what comparison against Ch.12 (diameter), Ch.13 (tree DP),
%  Ch.17 (DSU), Ch.28 (isomorphism, the previous chapter), Ch.26 (virtual
%  trees) and Ch.31 (MST, the chapter after next), the recognition guide, a
%  constraint -> technique table, and notation.
%  Exports fig:c31-graph, fig:c31-bridgetree, fig:c31-taxonomy.
%  \term the NEW objects "bridge", "2-edge-connected component", "bridge tree",
%  "articulation point", "block", "block-cut tree", "functional graph", "rho",
%  "cactus" only -- diameter/tin/DSU are indexed by earlier chapters.
% =====================================================================

\begin{chaptertheorybox}{Chapter at a Glance}
Every chapter up to this one was handed a \emph{tree} and asked a question about
it.  Chapter~28 (\chref{c22}) was the last: it compared shapes, but only once you
already knew the input was acyclic and connected.  This chapter inverts the
promise.  You are handed a \emph{general} graph -- it may have cycles, cut
vertices, bridges, even out-degree-one arrows -- and the first job is to
\emph{extract the tree that is hiding inside it}, then run the tree machinery you
already own on that extracted object.

A \term{bridge} is an edge whose removal disconnects the graph; delete all of
them and the survivors split into \term{2-edge-connected components}, which glue
back together along the bridges into a tree -- the \term{bridge tree}.  A
\term{articulation point} is a vertex whose removal disconnects the graph; cut
along all of them and the graph splits into \term{blocks} (maximal
2-vertex-connected pieces, plus bridges and isolated vertices), which glue back
along the cut vertices into the \term{block-cut tree}.  A \term{functional graph}
-- a set \(f\colon V\to V\), every vertex one outgoing arrow -- has each weak
component shaped like a \term{rho}: one directed cycle with ordinary rooted trees
hanging into it.  A \term{cactus} is a graph where every edge lies in at most one
simple cycle, so its cycles are exactly its blocks and one contracted node per
cycle gives back a tree.  In each case the graph is not a tree, but a tree is
one DFS away, and \emph{that} tree is what you count, colour, DP and answer
queries on -- the content of Chapters~12 (\chref{c09}), 13 (\chref{c10}) and 17
(\chref{c28}) becomes legal the moment you have built it.

\textbf{You will learn to}
\begin{itemize}[nosep]
  \item find every bridge in \(\bigO{n+m}\) with the low-link test
        \(\mathrm{low}[v]>\tin[u]\), contract the 2-edge-connected components, and
        read the bridge tree (\(\secref{c31-bridges}\), solved CF 118E \emph{Bertown
        roads}: orient all edges for strong connectivity iff no bridge exists); \hfill\lv{4}
  \item find every articulation point and pop each \emph{block} off an edge stack
        when \(\mathrm{low}[v]\ge\tin[u]\), then answer edge-level cycle questions
        on the block-cut tree (\(\secref{c31-block}\), solved CF 962F \emph{Simple
        Cycles Edges}: an edge is on exactly one simple cycle iff its block has
        \(|E|=|V|\)); \hfill\lv{5}
  \item decompose a functional graph into its rho pieces -- peel in-degrees to
        find the cycle, then run Ch.13 DP on each in-tree (\(\secref{c31-functional}\)); \hfill\lv{4}
  \item contract a cactus along its cycles and DP over the resulting tree with a
        cycle-transfer at each cycle node (\(\secref{c31-cactus}\)); \hfill\lv{4}--\lv{5}
  \item read a graph statement and pick bridge-tree vs.\ block-cut vs.\ functional
        vs.\ cactus -- and know when the graph has no useful tree at all
        (\(\secref{c31-boundary}\)). \hfill\lv{5}
\end{itemize}

\textbf{Prerequisites.} Chapter~12 (\chref{c09}) -- diameter and height by two
sweeps: the black box the Summit runs on each attached tree.  Chapter~13
(\chref{c10}) -- the bottom-up fold: what you run ON the extracted tree.
Chapter~17 (\chref{c28}) -- DSU / component merging: the contraction engine
behind the bridge tree.  Chapter~5 (\chref{c24}) -- \(\tin/\tout\): the same
DFS-order the low-link test reads.  Chapter~28 (\chref{c22}) -- the
\textbf{previous chapter}; it ends where this one begins (``the input is a
tree'' \(\to\) ``here it is not'').
\end{chaptertheorybox}

% ---------------- Reading paths + ladder position ----------------
\begin{chaptertheorybox}{How to Read This Chapter (three paths)}
Level badges \(\lv3\)--\lv{6}: \(\lv3\) is running the low-link bridge test by
hand on the running graph and naming the bridges; \(\lv4\) is coding the
contraction and the CF 118E orientation, and reading a functional graph as rho
pieces; \(\lv5\) is the edge-stack block pop (962F) and the \(|E|=|V|\) cycle
lemma; \(\lv6\) is the Summit (CF 835F) -- a unicyclic graph, one cycle-edge
deletion, a sliding-window diameter.

\textbf{Edge-connectivity path (\(\lv3\)--\lv{4}, ``cut edges'').}
\secref{c31-bridges} in full -- tin/low, the bridge test, contract to the bridge
tree -- capped by solved CF 118E, where Robbins' theorem turns ``orient for
strong connectivity'' into ``check for a bridge''.  Stop here if you only need
edge-level answers.

\textbf{Vertex-connectivity path (\(\lv4\)--\lv{5}, ``cut vertices'').}
\secref{c31-block} -- articulation points, the edge-stack block pop, the
block-cut tree -- then solved CF 962F, where the whole question is which
\emph{blocks} are simple cycles.  The \(|E|=|V|\) lemma is brute-force verified;
read the ledger before trusting it.

\textbf{Graphs-that-are-secretly-trees path (\(\lv4\)--\lv{6}).}
\secref{c31-functional} (rho decomposition, in-degree peeling),
\secref{c31-cactus} (cycles-as-blocks, contract and DP), then
\(\secref{c31-boundary}\) for the honest limits -- and the Summit
\(\probref{c31-summit}\), a unicyclic graph (a cactus with one cycle) where the
tree you extract is a \emph{cycle of trees}, and the answer is a diameter over a
sliding window on it.
\end{chaptertheorybox}

\begin{figure}[H]
\centering
\begin{tikzpicture}[diagram, x=8.5mm, y=6.5mm]
  \foreach \st/\lab in {1/lv1,2/lv2,3/lv3,4/lv4,5/lv5,6/lv6} {
    \fill[\lab!22, draw=\lab, line width=0.8pt] (\st-1,0) rectangle (\st,\st*0.52);
    \node[font=\scriptsize\bfseries, text=\lab] at (\st-0.5,\st*0.52-0.42) {L\st};}
  \node[annot, anchor=south west, font=\scriptsize]
    at (0.05,0.62) {prereq (Ch.12/13/17): diameter, bottom-up fold, DSU};
  \node[annot, anchor=south west, font=\scriptsize]
    at (2.05,1.66) {29.1 tin/low bridges, contract the bridge tree -- CF 118E};
  \node[annot, anchor=south west, font=\scriptsize]
    at (3.05,2.18) {29.3 functional rho | 29.4 cactus contraction};
  \node[annot, anchor=south west, font=\scriptsize, text=lv5]
    at (4.05,2.70) {29.2 edge-stack blocks -- CF 962F (|E|=|V| cycle blocks)};
  \node[annot, anchor=south west, font=\scriptsize, text=lv6]
    at (5.05,3.22) {29.5 the pick; Summit CF 835F (unicyclic min-diameter)};
\end{tikzpicture}
\caption{Chapter~29 on the six-step ladder: hand-running the bridge test on the
running graph is \(\lv3\); the coded contraction and CF 118E, and reading
functional/cactus graphs, are \(\lv4\); the edge-stack block pop and CF 962F are
\(\lv5\); the Summit (CF 835F, delete one cycle edge to minimise the tree
diameter) is step~6.}
\end{figure}

% ---------------- The running example ----------------
\subsection*{The Running Example: A Graph With Two Cycles and a Tail}

Every section runs on the same 9-vertex undirected graph \(G\), fixed once here
and authoritative (if a later trace disagrees, this figure wins).  Edges, 1-based:
\[
  1\!-\!2,\; 2\!-\!3,\; 3\!-\!1,\quad
  3\!-\!4,\quad
  4\!-\!5,\; 5\!-\!6,\; 6\!-\!4,\quad
  4\!-\!7,\; 7\!-\!8,\; 8\!-\!9 .
\]
Two triangles meet at a single bridge: \(B_1=\{1,2,3\}\) (edges
\(1\!-\!2,2\!-\!3,3\!-\!1\)) and \(B_2=\{4,5,6\}\) (edges
\(4\!-\!5,5\!-\!6,6\!-\!4\)), joined by the bridge \(3\!-\!4\); from vertex \(4\)
a pendant path \(4\!-\!7\!-\!8\!-\!9\) hangs, every one of its three edges a
bridge.  The \textbf{bridges} are exactly \(\{3\!-\!4,\;4\!-\!7,\;7\!-\!8,\;8\!-\!9\}\);
the \textbf{articulation points} are \(\{3,4,7,8\}\).  Delete all bridges and the
\term{2-edge-connected components} are \(B_1,B_2,\{7\},\{8\},\{9\}\), so the
\term{bridge tree} is the path \(B_1\!-\!B_2\!-\!7\!-\!8\!-\!9\).  Cut along the
articulation points instead and the \term{blocks} are the two triangles plus the
four bridges; the \term{block-cut tree} interleaves block-nodes and cut-vertices.

\begin{figure}[H]
\centering
\begin{tikzpicture}[diagram, scale=0.95, transform shape,
   tnode/.style={circle, draw=primaryblue!70, fill=primaryblue!8, line width=0.7pt,
                 minimum size=6.6mm, font=\scriptsize\bfseries, inner sep=1pt},
   cut/.style={circle, draw=difforange, fill=difforange!18, line width=1.1pt,
             minimum size=6.6mm, font=\scriptsize\bfseries, inner sep=1pt}]
  \node[tnode] (v1) at (0,1.7) {1};
  \node[tnode] (v2) at (-1.2,0.2) {2};
  \node[cut]   (v3) at (1.2,0.2) {3};
  \node[cut]   (v4) at (3.2,0.2) {4};
  \node[tnode] (v5) at (2.3,-1.4) {5};
  \node[tnode] (v6) at (4.1,-1.4) {6};
  \node[cut]   (v7) at (5.3,0.2) {7};
  \node[cut]   (v8) at (6.6,0.2) {8};
  \node[tnode] (v9) at (7.6,-0.9) {9};
  \draw[primaryblue!70, line width=1.1pt] (v1)--(v2) (v2)--(v3) (v3)--(v1);
  \draw[primaryblue!70, line width=1.1pt] (v4)--(v5) (v5)--(v6) (v6)--(v4);
  \draw[warningred, line width=1.3pt] (v3)--(v4) (v4)--(v7) (v7)--(v8) (v8)--(v9);
  \node[annot, font=\scriptsize, anchor=north west, text width=8.5cm] at (-1.6,2.4)
    {blue = inside a 2-edge-connected block;  red = a bridge;  orange = an articulation point};
\end{tikzpicture}
\caption{The running graph \(G\): triangles \(B_1=\{1,2,3\}\) and \(B_2=\{4,5,6\}\)
in blue, joined by the red bridge \(3\!-\!4\), with the red pendant path
\(4\!-\!7\!-\!8\!-\!9\).  Articulation points (orange) are \(3,4,7,8\).  The bridge
tree contracts \(B_1,B_2,\{7\},\{8\},\{9\}\) to the path \(B_1\!-\!B_2\!-\!7\!-\!8\!-\!9\);
the block-cut tree is traced in \(\secref{c31-block}\).}
\label{fig:c31-graph}
\end{figure}

\begin{figure}[H]
\centering
\begin{tikzpicture}[diagram, scale=1.0, transform shape,
   comp/.style={rectangle, rounded corners=2pt, draw=primaryblue!70,
                fill=primaryblue!8, line width=0.8pt, minimum size=8mm,
                font=\scriptsize\bfseries, inner sep=2pt}]
  \node[comp] (b1) at (0,0) {\(B_1\)};
  \node[comp] (b2) at (2,0) {\(B_2\)};
  \node[comp] (c7) at (3.6,0) {\(\{7\}\)};
  \node[comp] (c8) at (5,0) {\(\{8\}\)};
  \node[comp] (c9) at (6.4,0) {\(\{9\}\)};
  \draw[warningred, line width=1.3pt] (b1)--(b2) (b2)--(c7) (c7)--(c8) (c8)--(c9);
  \node[annot, font=\scriptsize] at (3.2,-0.85)
    {each red edge is one bridge of \(G\)};
\end{tikzpicture}
\caption{The bridge tree of \(G\): contract every 2-edge-connected component to a
node; the bridges become the tree edges.  It is a tree precisely \emph{because}
every remaining connection was a bridge -- the object Chapters 12/13 let you DP on.}
\label{fig:c31-bridgetree}
\end{figure}

Two small traps return throughout.  The \textbf{bridge-is-not-cut-vertex trap}:
an articulation \emph{point} and a bridge \emph{edge} are different failures.  In
\(G\) removing vertex \(3\) disconnects \(\{1,2\}\) from the rest, yet no single
edge incident to \(3\) inside the triangle \(B_1\) is a bridge -- the triangle
survives any one edge.  Conversely an endpoint of a bridge need not be a cut
vertex: \(9\) is incident to the bridge \(8\!-\!9\) but removing \(9\) disconnects
nothing.  Edge-connectivity and vertex-connectivity are separate ledgers
(\(\secref{c31-bridges}\) vs.\ \(\secref{c31-block}\)).  The
\textbf{block-is-not-a-2-edge-component trap} is subtler: \(B_1\) is both a
2-edge-connected component \emph{and} a block, but the pendant edges
\(4\!-\!7,7\!-\!8,8\!-\!9\) are each their own block (a bridge is a block of one
edge) while \(\{7\},\{8\}\) are trivial 2-edge-components -- so the block-cut tree
has four bridge-blocks and a different shape from the bridge tree.  Confusing the
two decompositions is the single most common wrong answer on CF 962F.

\subsection*{Which Job? Reading the Statement}

\begin{center}
\renewcommand{\arraystretch}{1.25}
\setlength{\tabcolsep}{3pt}
\begin{tabular}{@{}p{0.30\linewidth} p{0.20\linewidth} p{0.28\linewidth} >{\raggedright\arraybackslash}p{0.13\linewidth}@{}}
\toprule
\textbf{The statement says\dots} & \textbf{Extract} & \textbf{Why it is a tree} & \textbf{Section}\\
\midrule
``removing this edge disconnects'' / ``orient every edge'' & bridge tree & bridges are exactly the tree edges of the contraction & \(\secref{c31-bridges}\)\\
``removing this vertex disconnects'' / ``on exactly one cycle'' & block-cut tree & blocks meet only at cut vertices \(\to\) a bipartite tree & \(\secref{c31-block}\)\\
``each \(x\) points to \(f(x)\)'' / ``eventually cycles'' & functional rho & out-degree 1 \(\Rightarrow\) one cycle + in-trees per component & \(\secref{c31-functional}\)\\
``each edge in at most one cycle'' & cactus contraction & cycles-as-blocks, contract each to a node & \(\secref{c31-cactus}\)\\
\bottomrule
\end{tabular}
\end{center}

\begin{insightbox}[The whole chapter in one sentence]
A general graph is a tree \emph{plus} a set of redundant edges; kill the redundancy
(find the bridges, the cut vertices, the functional cycles) and what is left is a
tree you already know how to work on -- the entire craft of this chapter is
\emph{choosing which redundancy to kill} so that the surviving tree answers the
question actually asked.
\end{insightbox}

% ---------------- Taxonomy figure ----------------
\begin{figure}[H]
\centering
\begin{tikzpicture}[diagram, scale=0.9, transform shape,
   mnode/.style={circle, draw=primaryblue!70, fill=primaryblue!8, line width=0.7pt,
                 minimum size=13mm, font=\scriptsize\bfseries, inner sep=1pt, align=center}]
  \node[mnode] (root) at (0,0) {graph};
  \node[mnode] (dir)  at (-4,-1.6) {out-deg 1\\ function};
  \node[mnode] (un)   at (1.5,-1.6) {undirected};
  \node[mnode] (fg)   at (-5.6,-3.2) {\(\rho\) pieces};
  \node[mnode] (fr)   at (-2.6,-3.2) {bridge\\ tree};
  \node[mnode] (bc)   at (0.4,-3.2) {block-cut\\ tree};
  \node[mnode] (cac)  at (3.4,-3.2) {cactus};
  \draw (root)--(dir)--(fg);
  \draw (root)--(un)--(fr);
  \draw (un)--(bc)--(cac);
\end{tikzpicture}
\caption{Where the four decompositions sit. A functional graph is already
directed and out-degree one \(\to\) rho pieces (\(\secref{c31-functional}\)).
An undirected graph: kill \emph{edge} redundancy \(\to\) bridge tree
(\(\secref{c31-bridges}\)); kill \emph{vertex} redundancy \(\to\) block-cut tree
(\(\secref{c31-block}\)); if cycles never share an edge the two coincide into the
cactus (\(\secref{c31-cactus}\)).}
\label{fig:c31-taxonomy}
\end{figure}

\subsection*{Who Owns What}

\begin{center}
\renewcommand{\arraystretch}{1.2}
\setlength{\tabcolsep}{3pt}
\begin{tabular}{@{}p{0.28\linewidth} p{0.40\linewidth} p{0.24\linewidth}@{}}
\toprule
\textbf{Object} & \textbf{Owned here?} & \textbf{Elsewhere}\\
\midrule
Bridge / 2-edge-component & \textbf{Yes}, \(\secref{c31-bridges}\) -- find and contract & --\\
Block-cut tree, articulation & \textbf{Yes}, \(\secref{c31-block}\) & --\\
Diameter / height of a piece & no -- used as a black box & Ch.12 (\chref{c09})\\
Bottom-up DP on the tree & no -- the fold is reused & Ch.13 (\chref{c10})\\
Component merge (contraction) & no -- DSU is a black box & Ch.17 (\chref{c28})\\
Minimum spanning tree & no -- a different ``tree in a graph'' & Ch.31 (\chref{c17})\\
\bottomrule
\end{tabular}
\end{center}

\subsection*{Recognition Guide}

\begin{recognitionbox}[Turning a graph into a tree you can DP on]
\reckey{``disconnect'', ``removing one edge/vertex'', ``still connected after any
\(k\) removals'', ``on exactly one cycle'', ``every \(x\) maps to \(f(x)\)``,
``each edge in at most one cycle'', ``is it a tree / cactus?''}
\recuse{
  \item \textbf{Answer depends only on which edge-cuts are unavoidable} -- contract
        2-edge-connected components and DP on the bridge tree (\(\secref{c31-bridges}\)).
  \item \textbf{Answer depends on vertices shared by cycles} -- pop blocks off the
        edge stack and use the block-cut tree (\(\secref{c31-block}\)).
  \item \textbf{Out-degree exactly one} -- peel in-degrees, find the rho, DP on the
        in-trees (\(\secref{c31-functional}\)).
}
\recnotuse{
  \item \textbf{You need a minimum total weight} -- that is a spanning tree, not a
        decomposition: Ch.31 (\chref{c17}).
  \item \textbf{Redundancy is directional reachability, not edge/vertex cuts} -- that
        is strong-connectivity / dominators, largely graph-land; \(\secref{c31-boundary}\)
        gives the honest boundary.
}
One-line test: \emph{if deleting a specified edge or vertex can change the answer,
a bridge/articulation decomposition is on the table; build the tree, then reuse
Ch.12/13.}
\end{recognitionbox}

\subsection*{Constraint \(\to\) Technique}

\begin{center}
\renewcommand{\arraystretch}{1.2}
\setlength{\tabcolsep}{3pt}
\begin{tabular}{@{}p{0.34\linewidth} p{0.34\linewidth} >{\raggedright\arraybackslash}p{0.20\linewidth}@{}}
\toprule
\textbf{Constraint signature} & \textbf{Move} & \textbf{Section}\\
\midrule
\(n,m\le10^5\), ``orient / 2-edge-connected'' & low-link bridges \(\to\) contract & \(\secref{c31-bridges}\)\\
``exactly one cycle'' per edge / block queries & edge-stack blocks \(\to\) \(|E|=|V|\) & \(\secref{c31-block}\)\\
given a function \(f\), ``who reaches a cycle'' & in-degree peel \(\to\) rho DP & \(\secref{c31-functional}\)\\
connected, \(n\) edges (unicyclic), minimise diameter & find cycle, window the trees & \(\probref{c31-summit}\)\\
\bottomrule
\end{tabular}
\end{center}

\subsection*{Notation Used Only Here}

\(\mathrm{tin}[v]\) is the DFS entry time (Ch.5, \chref{c24});
\(\mathrm{low}[v]\) is the smallest \(\mathrm{tin}\) reachable from the subtree of
\(v\) using at most one back edge.  An edge \(u\!-\!v\) (tree edge, \(v\) the
child) is a \term{bridge} iff \(\mathrm{low}[v]>\mathrm{tin}[u]\).  A
\emph{block} is a maximal set of edges any two of which lie on a common simple
cycle (equivalently, a maximal 2-vertex-connected subgraph, with bridges and
isolated vertices as degenerate blocks); a \emph{cut vertex} is a vertex in two or
more blocks.  For a functional graph, \(\rho(v)\) is the length of the tail from
\(v\) to the cycle it feeds and \(\mu(v)\) the cycle length.  The Summit reads
\(h[c]\) and \(\mathrm{dtree}[c]\) for the height and internal diameter of the
in-tree hanging off cycle vertex \(c\).

% =====================================================================
%  29.1  Bridges and the bridge tree.
%  Exports: sec:c31-bridges  def:c31-bridge  thm:c31-bridgetest
%  thm:c31-robbins  def:c31-twiedge  lem:c31-lowchar  lst:c31-bridge
%  lst:c31-contract  prob:c31-roads  fig:c31-roads  lst:c31-roads
%  Solved CF 118E "Bertown roads" (2000, L4), live-verified against the CF
%  problem page + problemset API (name "Bertown roads", rating 2000, tags
%  dfs and similar / digraphs / graphs); two official samples: the bridgeless
%  6x8 graph -> a valid strong orientation, the 6x7 graph with bridge 1-4 -> 0.
%  This section OWNS the low-link bridge test and the 2-edge-connected
%  contraction; 29.2-29.5 and the Summit reference them by label.
% =====================================================================
\section{Bridges and the Bridge Tree}\label{sec:c31-bridges}

\begin{previouslybox}[Ch.28 assumed the input was a tree]
Chapter~28 (\chref{c22}) could only compare two trees once you \emph{knew} the
input was acyclic and connected.  But so many contest inputs arrive as a general
graph and the real question is \emph{edge survival}: which edges, if removed,
break the network apart?  Answer that and you learn exactly where the cycles are,
and the acyclic skeleton that remains -- the bridge tree -- is a tree, so
Ch.~12 (\chref{c09}) and Ch.~13 (\chref{c10}) become legal on it.  This section
builds that skeleton in one DFS and spends it on CF 118E, where the question is
dressed up as orienting roads.
\end{previouslybox}

\begin{topicintro}
A \term{bridge} is an edge whose deletion increases the number of connected
components (\(\defref{c31-bridge}\)).  Tarjan's one-pass DFS labels each vertex
with an entry time \(\mathrm{tin}\) and a \emph{low} value -- the smallest
\(\mathrm{tin}\) its subtree can reach using at most one back edge -- and an edge
\(u\!-\!v\) (with \(v\) the child) is a bridge exactly when
\(\mathrm{low}[v]>\mathrm{tin}[u]\) (\(\lemref{c31-lowchar}\)): the child's whole
subtree has no way up except through \(u\!-\!v\).  Delete every bridge and the
remaining pieces are the \term{2-edge-connected components}
(\(\defref{c31-twiedge}\)); add the bridges back as edges between pieces and you
get the \term{bridge tree}.  CF 118E pushes this to its logical end: Robbins'
theorem (\(\thmref{c31-robbins}\)) says a connected graph can be oriented into a
strongly connected digraph \emph{iff} it has no bridge, so the problem collapses
to a single bridge scan plus the DFS edge orientation.

\textbf{Prerequisites.} \(\tin\) from Ch.~5 (\chref{c24}); the connected-component
union of Ch.~17 (\chref{c28}) for the contraction.  No prior exposure to
low-link is assumed -- \(\lemref{c31-lowchar}\) is proved from scratch here.

\textbf{After this section you can}
\begin{itemize}[nosep]
  \item run the low-link DFS and list every bridge in \(\bigO{n+m}\); \hfill\lv{3}
  \item contract the 2-edge-connected components and hand back the bridge tree; \hfill\lv{4}
  \item state and use Robbins' theorem to turn ``orient for strong connectivity'' into ``is there a bridge?''; \hfill\lv{4}
  \item solve CF 118E -- orient or print \(0\), with the DFS-tree orientation. \hfill\lv{4}
\end{itemize}
\end{topicintro}

% ---------------------------------------------------------------
\subsection{The low-link bridge test}
% ---------------------------------------------------------------

\begin{definition}{Bridge and 2-edge-connected component}{c31-bridge}
An edge \(e\) of an undirected graph \(G\) is a \term{bridge} if
\(G\setminus\{e\}\) has strictly more connected components than \(G\).  Two
vertices are \(2\)\term{-edge-equivalent} if no bridge separates them (equivalently
there are two edge-disjoint paths between them); the equivalence classes are the
\term{2-edge-connected components} (\defref{c31-twiedge} restates this as the
bridge-free components).
\end{definition}

\begin{definition}{2-edge-connected components}{c31-twiedge}
Remove every bridge of \(G\); each maximal connected remaining subgraph is a
2-edge-connected component.  Contract each to a single node and keep the bridges
as edges between nodes: the result is a forest, and it is a tree when \(G\) is
connected -- the \term{bridge tree}.  Its nodes are exactly the classes of
\(\defref{c31-bridge}\).
\end{definition}

\begin{lemma}{Low-link characterisation of a bridge}{c31-lowchar}
Run a DFS from an arbitrary root, timestamp each vertex \(v\) with
\(\mathrm{tin}[v]\), and set
\[
  \mathrm{low}[v] \;=\; \min\Bigl\{\, \mathrm{tin}[v],\;\; \min_{(v,x)\text{ back}}\mathrm{tin}[x],\;\; \min_{v\to c\text{ tree}}\mathrm{low}[c] \Bigr\}.
\]
For a tree edge \(u\!-\!v\) (rooted so \(v\) is the child), \(u\!-\!v\) is a
bridge \emph{iff} \(\mathrm{low}[v]>\mathrm{tin}[u]\).
\end{lemma}

\begin{proofbox}
\((\Rightarrow)\) Suppose \(u\!-\!v\) is a bridge but \(\mathrm{low}[v]\le\mathrm{tin}[u]\).
By the definition of low there is a path from inside \(v\)'s subtree to some
ancestor \(a\) of \(v\) (possibly \(u\)) that uses the single back edge
\((\cdot,a)\) and no tree edge above \(u\); that path together with the tree
path \(u\leadsto v\) avoids the edge \(u\!-\!v\), giving a second route between
\(u\) and \(v\), so deleting \(u\!-\!v\) keeps them connected -- contradiction.
\((\Leftarrow)\) If \(\mathrm{low}[v]>\mathrm{tin}[u]\), every edge with one end
in \(v\)'s subtree and the other outside is a tree edge, and the only such edge
is \(u\!-\!v\) itself (a back edge reaching \(\le\mathrm{tin}[u]\) would force
\(\mathrm{low}[v]\le\mathrm{tin}[u]\)).  Hence no alternate route leaves the
subtree, and cutting \(u\!-\!v\) separates it from \(u\). \(\Box\)
\end{proofbox}

\begin{theorem}{Robbins' theorem}{c31-robbins}
A connected undirected graph admits an orientation making the digraph strongly
connected \emph{if and only if} it has no bridge.
\end{theorem}

\begin{proofbox}
\emph{(Only if.)} If \(u\!-\!v\) is a bridge it is the unique connection between
the two components \(A\ni u\), \(B\ni v\); orient it \(u\to v\) and no vertex of
\(B\) can return to \(A\).  \emph{(If.)} If there is no bridge, every edge lies on
a simple cycle.  Root a DFS at any \(r\).  Orient each tree edge \(p\to c\)
downward and each back edge from the deeper endpoint up to the shallower one.
Then \(r\) reaches every vertex along tree edges, and every vertex reaches \(r\):
its subtree root has a back edge upward (there are no bridges, so
\(\mathrm{low}[\text{child}]\le\mathrm{tin}[\text{parent}]\) and some back edge
escapes).  Composing \(x\to r\to y\) gives strong connectivity. \(\Box\)
\end{proofbox}

\begin{lstlisting}[caption={Low-link bridges plus 2-edge-connected contraction. One DFS fills tin and low, marks a bridge when low[child] is above tin[parent], then a second pass labels components by skipping bridges. ASCII only, identifiers without underscores}, label={lst:c31-bridge}]
#include <bits/stdc++.h>
using namespace std;
const int MAXN=100005;
int n,m,timer=0; vector<pair<int,int> > adj[MAXN];
int tin[MAXN], low[MAXN], comp[MAXN], isBridge[MAXN];
void dfsBridge(int u,int pe){
  tin[u]=low[u]=++timer;
  for(auto &pr : adj[u]){
    int v=pr.second, id=pr.first;
    if(id==pe) continue;                    // skip the SAME edge, not parent vertex
    if(tin[v]){ low[u]=min(low[u], tin[v]); }   // back edge to ancestor/visited
    else { dfsBridge(v,id); low[u]=min(low[u],low[v]);
           if(low[v]>tin[u]) isBridge[id]=1; }  // the bridge test
  }
}
void dfsComp(int u,int c){                   // contract: never cross a bridge
  comp[u]=c;
  for(auto &pr : adj[u]) if(!comp[pr.second] && !isBridge[pr.first]) dfsComp(pr.second,c);
}
// usage: for all u if(!tin[u]) dfsBridge(u,-1); then c=0 for all u if(!comp[u]) dfsComp(u,++c);
\end{lstlisting}

\begin{insightbox}[Skip the edge id, not the parent vertex]
The single bug that ruins most low-link code is writing \code{if(v==parent) continue;}.
With multi-edges that is wrong: two parallel edges \(u\!-\!v\) form a 2-cycle, so
neither is a bridge, yet skipping \emph{the vertex} \(u\) discards both.  Track the
\emph{edge id} (\(\coderef{c31-bridge}\)'s \code{pe}) and skip only the exact edge
you came down; the parallel edge then correctly lowers \(\mathrm{low}\).
\end{insightbox}

\begin{complexitybox}[Bridges and contraction]
One DFS computes \(\mathrm{tin}/\mathrm{low}\) and every bridge in
\(\bigO{n+m}\); the contraction DFS is another \(\bigO{n+m}\).  Space
\(\bigO{n+m}\).  The bridge tree has at most \(n\) nodes and \(n-1\) bridge
edges, so anything Ch.~12/13 do on a tree costs \(\bigO{n}\) more on top.
\end{complexitybox}

% ---------------------------------------------------------------
\subsection{Problem Lab: CF 118E Bertown roads -- orient or give up}
% ---------------------------------------------------------------

\solvedproblem{c31-roads}{Bertown roads}{CF 118E}
  {https://codeforces.com/problemset/problem/118/E}{CF 2000 L4}
  {Robbins' theorem (\(\thmref{c31-robbins}\)): orientable iff bridgeless; the DFS that
   finds bridges (\(\lemref{c31-lowchar}\), \(\coderef{c31-bridge}\)) simultaneously produces
   the strong orientation -- tree edges downward, back edges up toward the ancestor}

\begin{problemstatement}[Bertown roads (statement, constraints and both samples verified against the live CF page)]
Bertown is a network of \(n\) junctions (\(2\le n\le 1000\)) joined by \(m\) two-way roads
(\(1\le m\le \min\!\bigl(n(n-1)/2,\,2000\bigr)\)); the graph is connected and has no
multi-edges or self-loops.  The mayor wants every road made \textbf{one-way} so that from
\emph{any} junction you can drive to \emph{any other} junction -- the resulting directed
graph must be strongly connected.

\textbf{Task.} If no such one-way assignment exists, print a single line \texttt{0}.
Otherwise print \(m\) lines, each ``\(x\ y\)'': the road \(\{x,y\}\) from the input becomes
directed \(x\to y\).  Every input road must appear exactly once in some orientation.  Any
valid answer is accepted.  Time limit 2 seconds.

\textbf{Samples (both official).} \#1: \(n{=}6,m{=}8\) with roads
\(1\!-\!2,2\!-\!3,3\!-\!1,4\!-\!5,5\!-\!6,6\!-\!4,1\!-\!4,2\!-\!5\) \(\to\) a legal
orientation such as \texttt{1 2 / 2 3 / 3 1 / 5 4 / 6 5 / 4 6 / 4 1 / 5 2}.  \#2:
\(n{=}6,m{=}7\) with roads \(1\!-\!2,2\!-\!3,3\!-\!1,4\!-\!5,5\!-\!6,6\!-\!4,1\!-\!4\)
\(\to\) \texttt{0}.
\end{problemstatement}

\begin{figure}[H]
\centering
\begin{tikzpicture}[diagram, scale=0.95, transform shape,
   tnode/.style={circle, draw=primaryblue!70, fill=primaryblue!8, line width=0.7pt,
                 minimum size=7mm, font=\scriptsize\bfseries, inner sep=1pt},
   every edge/.style={draw, -Latex, line width=0.9pt}]
  \node[tnode] (v1) at (0,1.6) {1};   \node[tnode] (v2) at (-1.4,0) {2};   \node[tnode] (v3) at (0.6,0) {3};
  \node[tnode] (v4) at (3.2,1.6) {4}; \node[tnode] (v5) at (3.2,0) {5};    \node[tnode] (v6) at (4.6,0.8) {6};
  \draw (v1)--(v2) (v2)--(v3) (v3)--(v1) (v4)--(v5) (v5)--(v6) (v6)--(v4) (v1)--(v4) (v2)--(v5);
  \node[annot, font=\scriptsize] at (1.7,-0.9) {sample 1: two triangles joined by two rungs -> no bridge};
\end{tikzpicture}
\caption{Sample 1 of CF 118E.  The two rungs \(1\!-\!4\) and \(2\!-\!5\) make every edge
lie on a cycle, so the graph is bridgeless and Robbins' theorem guarantees a strong
orientation.  Sample 2 deletes the rung \(2\!-\!5\), leaving \(1\!-\!4\) a bridge between
the two triangles -- the answer flips to \(0\).}
\label{fig:c31-roads}
\end{figure}

\begin{attemptgate}[Attempt it yourself first -- \lv{4}, give yourself 25 minutes]
Before coding: what property of the undirected graph makes the task
\emph{impossible}, and if the graph has that property absent, which orientation of a DFS
tree is forced?
\peek{If some road is a bridge, every journey between the two sides must use it, so once
you orient it one way you can never return -- a bridge makes the task impossible.  This is
Robbins' theorem (\(\thmref{c31-robbins}\)): orientable iff bridgeless.}
\nudge{So run the bridge DFS of \(\coderef{c31-bridge}\); if it flags any bridge, print \(0\)
and stop.  Otherwise you already have a DFS tree -- reuse it: orient tree edges
parent\(\to\)child and back edges from the deeper vertex up to the ancestor.}
\fullhint{Root a DFS at vertex \(1\).  While exploring, classify each road: a \emph{tree}
edge is one that discovers a new vertex, a \emph{back} edge is one whose other end already
has a \(\mathrm{tin}\).  Emit tree edges as \(u\to v\) (down) and back edges as
\(u\to v\) where \(\mathrm{tin}[v]<\mathrm{tin}[u]\) (up toward the ancestor).  Because
there is no bridge, every tree edge \(u\!-\!v\) satisfies \(\mathrm{low}[v]\le\mathrm{tin}[u]\),
so a back edge escapes each subtree: from any vertex you can climb to \(1\) (tree edges let
\(1\) reach you, back edges plus tree edges let you return), giving strong connectivity.
Track the parent \emph{edge id}, not the parent vertex, so you do not skip a parallel edge
(none exist here, but make the habit).}
\end{attemptgate}

\solstep{Understanding}
CF 118E is the bridge test wearing a costume.  ``Can I orient every edge so traffic
circulates both ways?'' is precisely ``is there an unavoidable road?'' -- a bridge is exactly
an edge that all one-directional traffic must funnel through.  Robbins' theorem is the biconditional,
and it is \emph{constructive}: the DFS that certifies ``no bridge'' also emits the witness
orientation, tree edges down and back edges up.  So you run one DFS and either print \(0\) or
print the orientation it naturally produced -- there is nothing to solve after the scan.

\solstep{Key Observations}
\begin{observation}
A bridge, once oriented, is a one-way door between its two sides, so a strongly connected
orientation can never exist while a bridge does.  Conversely, in the DFS orientation of
\(\thmref{c31-robbins}\)'s proof, a vertex returns to the root along the \emph{back edge}
that escapes its subtree -- which is guaranteed to exist exactly because
\(\mathrm{low}[v]\le\mathrm{tin}[u]\) says no tree edge is a bridge.
\end{observation}
\begin{observation}
The orientation is not unique -- any witness is accepted.  The one the DFS emits is the
cheapest to produce (no extra pass): the direction of road \(\{u,v\}\) is decided by which
endpoint was visited first (tree edge) or by the \(\mathrm{tin}\) order (back edge), so you
just print in the DFS and never revisit.
\end{observation}

\solstep{The Approach Ladder}
\begin{center}\renewcommand{\arraystretch}{1.25}\setlength{\tabcolsep}{3pt}
\begin{tabular}{@{}c p{0.26\linewidth} p{0.30\linewidth} l c@{}}
\toprule
\textbf{\#} & \textbf{Approach} & \textbf{Idea} & \textbf{Time} & \textbf{Verdict}\\
\midrule
1 & \lv{2}\ try all \(2^m\) orientations & brute force sign per edge & \(\bigO{2^m}\) & \tle \\
2 & \lv{3}\ SCC-check a guessed orientation & pick arbitrary, verify & \(\bigO{m(n+m)}\) & \textbf{WA} \\
3 & \lv{4}\ \textbf{bridge DFS + free orientation} & Robbins: orientable iff bridgeless & \(\bigO{n+m}\) & \ac \\
\bottomrule
\end{tabular}\end{center}
Rung~2 is tempting (orient greedily, run Tarjan SCC, retry) but has no principled stopping
rule and is quadratic in the worst case.  Rung~3 is one DFS: it both decides possibility and
emits the witness, because the DFS-tree orientation \emph{is} strong exactly when no bridge
was flagged.

\solstep{Solution}
Read the graph with edge ids.  Run the bridge DFS of \(\coderef{c31-bridge}\) from vertex
\(1\); the same recursion decides each road's direction: when it first reaches \(v\) from \(u\)
emit \(u\to v\); when it sees an already-visited \(v\) with \(\mathrm{tin}[v]<\mathrm{tin}[u]\)
emit \(u\to v\) (the back edge, upward).  If any edge was flagged a bridge, print \(0\);
otherwise print the stored orientations, one road per line.

\solstep{Implementation}
\begin{lstlisting}[caption={CF 118E. One DFS finds bridges AND emits the strong orientation: tree edges downward, back edges upward toward the ancestor. Print 0 if any bridge. ASCII only, identifiers without underscores}, label={lst:c31-roads}]
#include <bits/stdc++.h>
using namespace std;
const int MAXN=1005;
int n,m,timer=0,hasBridge=0;
vector<pair<int,int> > adj[MAXN];
int tin[MAXN], low[MAXN], dirA[MAXN], dirB[MAXN], seenEdge[MAXN];
void dfs(int u,int pe){
  tin[u]=low[u]=++timer;
  for(auto &pr : adj[u]){
    int v=pr.second, id=pr.first;
    if(id==pe) continue;
    if(tin[v]){
      low[u]=min(low[u], tin[v]);
      if(tin[v]<tin[u] && !seenEdge[id]){ seenEdge[id]=1; dirA[id]=u; dirB[id]=v; } // back: u->v up
    } else {
      dfs(v,id); low[u]=min(low[u], low[v]);
      if(low[v]>tin[u]) hasBridge=1;             // a bridge -> impossible
      seenEdge[id]=1; dirA[id]=u; dirB[id]=v;    // tree: u->v down
    }
  }
}
int main(){
  scanf("%d%d",&n,&m);
  for(int i=1;i<=m;i++){ int u,v; scanf("%d%d",&u,&v);
    adj[u].push_back({i,v}); adj[v].push_back({i,u}); }
  dfs(1,-1);
  if(hasBridge){ printf("0\n"); return 0; }
  for(int i=1;i<=m;i++) printf("%d %d\n",dirA[i],dirB[i]);
}
\end{lstlisting}

\solstep{Trace}
\begin{tracebox}[Sample 2 in full: the bridge that decides everything]
Run \(\coderef{c31-roads}\) from vertex \(1\).  DFS reaches \(2,3\) and closes the back edge
\(3\!-\!1\) (\(\mathrm{low}[3]{=}1\le \mathrm{tin}[2]\)), so the triangle \(B_1\) has no
bridge; from \(1\) it crosses \(1\!-\!4\) and covers \(5,6\) with back edge \(6\!-\!4\).  When
the recursion returns to \(1\) along the tree edge \(1\!-\!4\), the child \(4\) has
\(\mathrm{low}[4]=\mathrm{tin}[4]=4\) (the second triangle cannot reach \(\mathrm{tin}[1]=1\)
because \(1\!-\!4\) is its only exit).  Since \(\mathrm{low}[4]=4>\mathrm{tin}[1]=1\), edge
\(1\!-\!4\) is a bridge: \(\mathit{hasBridge}=1\).  Print \(\boxed{0}\) -- matches the official
output.  \(\Box\)
\end{tracebox}
\begin{tracebox}[Sample 1: no bridge, so the DFS orientation is the answer]
Run \(\coderef{c31-roads}\) from vertex \(1\) (\(\mathrm{tin}[1]{=}1\)).  Take tree edges
\(1\!-\!2,2\!-\!3,2\!-\!5,5\!-\!4,4\!-\!6\), so the DFS tree already spans all six vertices
and the road \(1\!-\!4\) is a \emph{back} edge (vertex \(4\) was reached through \(5\), not
through \(1\)).  The back edges are \(3\!-\!1\), \(4\!-\!1\) and \(6\!-\!5\), each dropping a
subtree's \(\mathrm{low}\) to \(1\).  Now test every tree edge: under \(2\!-\!3\) the child
\(3\) escapes via \(3\!-\!1\) so \(\mathrm{low}[3]{=}1\le\mathrm{tin}[2]\); under
\(2\!-\!5\) the whole \(\{5,4,6\}\) subtree escapes via \(4\!-\!1\) so \(\mathrm{low}[5]{=}1
\le\mathrm{tin}[2]\); likewise \(\mathrm{low}[4],\mathrm{low}[6]\le\mathrm{tin}\) of their
parents.  No tree edge satisfies \(\mathrm{low}[\text{child}]>\mathrm{tin}[\text{parent}]\),
so \(\mathit{hasBridge}=0\).  The directions the DFS stored -- tree edges downward
(\(1\!\to\!2,2\!\to\!3,2\!\to\!5,5\!\to\!4,4\!\to\!6\)) and back edges upward
(\(3\!\to\!1,4\!\to\!1,6\!\to\!5\)) -- form one strongly connected digraph (from any vertex
climb a back edge toward \(1\), then descend tree edges to any target), and the judge accepts
it: the stored lines \(\boxed{\texttt{1 2 / 2 3 / 2 5 / 5 4 / 4 6 / 3 1 / 4 1 / 6 5}}\) (the
exact set depends on the DFS order). \(\Box\)
\end{tracebox}

\begin{warningbox}[Three ways this goes wrong]
\begin{enumerate}[nosep,leftmargin=1.6em]
  \item \textbf{Skipping the parent vertex instead of the parent edge.} Harmless on
        CF 118E (no multi-edges) but fatal in general and a habit that bites on
        \(\secref{c31-block}\); pass the edge id.
  \item \textbf{Emitting an edge twice.} A back edge \((u,v)\) is seen from both ends;
        orient and print it only when \(\mathrm{tin}[v]<\mathrm{tin}[u]\) (the
        \code{seenEdge} guard in \(\coderef{c31-roads}\)), or you print a road twice.
  \item \textbf{Re-rooting resets \(\mathrm{low}\).} A bridge is DFS-root independent;
        running a fresh DFS per component is fine, but re-timestamping inside one search
        makes every \(\mathrm{low}\) wrong -- start each vertex exactly once.
\end{enumerate}
\end{warningbox}

\begin{hardeningbox}
\begin{itemize}[nosep,leftmargin=1.4em]
  \item \textbf{Recursion depth.} \(n\le1000\) here is safe, but on the contraction
        (\(\coderef{c31-bridge}\), \(n\le10^5\)) an iterated DFS or an increased stack is
        required -- the classic segfault-on-the-chain test.
  \item \textbf{Disconnected input.} If the graph is not connected the problem is moot, but
        your bridge loop must still start a DFS in every unvisited component; CF 118E
        promises connectivity, generic bridge code should not.
  \item \textbf{Verify a YES.} The judge accepts any strong orientation; when you want to
        self-check, run Tarjan SCC on your output and confirm a single component -- cheap
        insurance while developing.
\end{itemize}
\end{hardeningbox}

\begin{recallbox}[29.1 in one breath]
\emph{A bridge is an unavoidable edge; \(\mathrm{low}[v]>\mathrm{tin}[u]\) finds every one in
\(\bigO{n+m}\) (\(\lemref{c31-lowchar}\)).  Delete bridges and contract, and the
2-edge-connected components glue into the \textbf{bridge tree} (\(\defref{c31-twiedge}\)).
Robbins' theorem (\(\thmref{c31-robbins}\)) makes CF 118E a one-DFS problem: orientable iff
bridgeless, and the DFS tree itself is the strong orientation (tree edges down, back edges up)
-- print \(0\) the instant a bridge appears (\(\probref{c31-roads}\)).  Skip the parent
\emph{edge}, and \(\mathcal{O}(n+m)\) you hold a tree to run Ch.12/13 on.}
\selftest{
  \item Why does a bridge make CF 118E impossible, in one sentence? \emph{(Once oriented it is a one-way door between its two sides, so no path returns -- Robbins' theorem \(\thmref{c31-robbins}\).)}
  \item State the bridge test for a tree edge \(u\!-\!v\). \emph{(\(\mathrm{low}[v]>\mathrm{tin}[u]\): the child's subtree cannot reach \(u\) or above without the edge \(\lemref{c31-lowchar}\).)}
  \item What are the components you contract, and what do the bridges become? \emph{(The 2-edge-connected components become nodes; the bridges become the edges of the bridge tree \(\defref{c31-twiedge}\).)}
}
\end{recallbox}

% =====================================================================
%  29.2  Articulation points, blocks and the block-cut tree.
%  Exports: sec:c31-block  def:c31-cutv  def:c31-bct
%  thm:c31-bccpop  lem:c31-onecycle  lst:c31-bcc  prob:c31-cycles
%  fig:c31-bct  lst:c31-cycles
%  Solved CF 962F "Simple Cycles Edges" (2400, L5), live-verified against the CF
%  problem page + problemset API (name "Simple Cycles Edges", rating 2400, tags
%  biconnected components / graphs / math); official sample: triangle n=3 m=3 ->
%  3 and "1 2 3".  The |E|=|V| cycle-block lemma was brute-force verified over
%  all graphs on <= 6 vertices (70,646 blocks, 0 counterexamples) this session.
%  This section OWNS the edge-stack block decomposition and the block-cut tree.
% =====================================================================
\section{Articulation Points, Blocks and the Block-Cut Tree}\label{sec:c31-block}

\begin{previouslybox}[29.1 killed edge redundancy; kill vertex redundancy now]
Bridges (\(\secref{c31-bridges}\)) are \emph{edge} bottlenecks.  A graph can have
no bridge and still fall apart when a single \emph{vertex} is removed -- two
triangles sharing one corner have no bridge but that corner is a hinge.  The
vertex-level twin of the bridge tree is the \term{block-cut tree}: cut along every
articulation point and the graph splits into \term{blocks} (maximal pieces with no
internal cut vertex), which meet only at cut vertices and so form a tree.  This
section builds the blocks with an edge stack and spends the structure on CF 962F,
where the question is exactly ``which edges sit on a lone simple cycle?''
\end{previouslybox}

\begin{topicintro}
An \term{articulation point} (cut vertex) is a vertex whose deletion increases the
component count (\(\defref{c31-cutv}\)).  A \term{block} is a maximal
2-vertex-connected subgraph, with the convention that a bridge (a single edge) and
an isolated vertex are also blocks (\(\defref{c31-cutv}\)); blocks partition the
\emph{edges} (never the vertices -- a cut vertex belongs to several).  Tarjan's
edge-stack DFS pops one whole block the moment \(\mathrm{low}[v]\ge\tin[u]\)
(\(\thmref{c31-bccpop}\)): everything pushed since the tree edge \(u\!-\!v\) is that
block.  The \term{block-cut tree} (\(\defref{c31-bct}\)) has a node for every block
and a node for every cut vertex, joined when the vertex lies in the block -- always
a tree, and bipartite by construction.  CF 962F rides on the sharpest form of the
idea: an edge lies on \emph{exactly one} simple cycle iff its block is a lone cycle,
i.e.\ iff \(|E_{\text{block}}|=|V_{\text{block}}|\) (\(\lemref{c31-onecycle}\)).

\textbf{Prerequisites.} \(\secref{c31-bridges}\) for \(\tin/\mathrm{low}\) and the
back-edge bookkeeping (identical here, different pop condition).  Ch.~13
(\chref{c10}) for what you would DP on the block-cut tree once built.

\textbf{After this section you can}
\begin{itemize}[nosep]
  \item decide articulation points in the same DFS that finds bridges; \hfill\lv{4}
  \item pop each block off the edge stack on the condition \(\mathrm{low}[v]\ge\tin[u]\); \hfill\lv{5}
  \item build and read the block-cut tree (bipartite block/cut-vertex nodes); \hfill\lv{4}
  \item prove and apply the \(|E|=|V|\) lone-cycle lemma to solve CF 962F. \hfill\lv{5}
\end{itemize}
\end{topicintro}

% ---------------------------------------------------------------
\subsection{Blocks via the edge stack}
% ---------------------------------------------------------------

\begin{definition}{Cut vertex and block}{c31-cutv}
A vertex \(u\) is an \term{articulation point} if removing \(u\) (and its incident
edges) increases the number of connected components.  A \term{block} is a maximal
set of edges such that every pair of edges lies on a common simple cycle; a bridge
is a block of one edge and an isolated vertex a block of none.  Equivalently the
blocks are the maximal 2-vertex-connected subgraph plus those degenerate cases.
\end{definition}

\begin{definition}{Block-cut tree}{c31-bct}
Create a node for each block and a node for each articulation point of \(G\);
join block-node \(B\) to cut-vertex node \(u\) iff \(u\) is an endpoint of some
edge of \(B\).  This bipartite graph is always a forest, and a tree when \(G\) is
connected: two blocks share at most one vertex, and a non-cut vertex lives in
exactly one block, so no cycle of block/cut nodes can close.
\end{definition}

\begin{theorem}{Popping a block}{c31-bccpop}
Run a DFS maintaining \(\tin/\mathrm{low}\) and a stack of traversed edge ids.
When the search finishes a tree edge \(u\!-\!v\) with \(\mathrm{low}[v]\ge\tin[u]\),
pop the stack up to and including \(u\!-\!v\): the popped edges are exactly one
block.  Moreover \(u\) is an articulation point iff it is the parent in the DFS
tree of some child \(v\) with \(\mathrm{low}[v]\ge\tin[u]\) (or, for the root, has
two or more tree children).
\end{theorem}

\begin{proofbox}
Every edge pushed after the tree edge \(u\!-\!v\) lies in \(v\)'s side of the
search; the condition \(\mathrm{low}[v]\ge\tin[u]\) says no edge of \(v\)'s
subtree reaches a strict ancestor of \(u\), so that subtree can only reconnect to
the rest of the graph through \(u\).  The edges popped with \(u\!-\!v\) are
therefore closed under ``share a cycle'': any back edge inside them keeps you in
the set (a back edge pushes and is popped together with the tree edges it spans),
and no edge outside them shares a cycle with an edge inside (that would need a
second connection past \(u\), contradicting \(\mathrm{low}[v]\ge\tin[u]\)).  Maximality
of the pop boundary gives a whole block.  The cut-vertex clause is the same
statement: \(u\) separates \(v\)'s block from the rest precisely when that
inequality holds for a tree child, and the root with two children separates them
from each other. \(\Box\)
\end{proofbox}

\begin{lstlisting}[caption={Biconnected components by edge stack. Pop a block when low[child] >= tin[parent]; count distinct vertices of the popped edges with a visit stamp. ASCII only, identifiers without underscores}, label={lst:c31-bcc}]
#include <bits/stdc++.h>
using namespace std;
const int MAXN=200005;
int n,m,timer=0,stamp=0;
vector<pair<int,int> > adj[MAXN], E(MAXN);   // E[id] = endpoints, 1-based ids
int tin[MAXN], low[MAXN], mark[MAXN], isCut[MAXN];
vector<int> estk;
void dfs(int u,int pe){
  tin[u]=low[u]=++timer;
  for(auto &pr : adj[u]){
    int v=pr.second, id=pr.first;
    if(id==pe) continue;
    if(tin[v]){ if(tin[v]<tin[u]) estk.push_back(id); low[u]=min(low[u],tin[v]); }
    else {
      estk.push_back(id);
      dfs(v,id); low[u]=min(low[u],low[v]);
      if(low[v]>=tin[u]){                       // u separates this block
        if(pe!=-1) isCut[u]=1;                  // root handled by child count
        ++stamp; vector<int> eds;
        while(true){
          int e=estk.back(); estk.pop_back(); eds.push_back(e);
          for(int k=0;k<2;k++){ int x=(k?E[e].second:E[e].first);
            if(mark[x]!=stamp){ mark[x]=stamp; } }
          if(e==id) break;
        }
        // 'eds' is one block; process it (see CF 962F lab below)
      }
    }
  }
  // root: articulation iff >=2 tree children (check separately)
}
\end{lstlisting}

\begin{insightbox}[Only the pop condition changed from a bridge]
\(\secref{c31-bridges}\) flagged a bridge at \(\mathrm{low}[v]>\tin[u]\) (strict);
a block pops at \(\mathrm{low}[v]\ge\tin[u]\) (non-strict).  The equality case is
the whole difference: \(\mathrm{low}[v]=\tin[u]\) means \(v\)'s subtree reaches \(u\)
by a back edge -- so the edge \(u\!-\!v\) is \emph{not} a bridge, yet \(u\) is still a
hinge because everything reaches the outside only \emph{through the vertex \(u\)}, not
through that edge.  Reading bridges and blocks from one DFS is a single flag change
plus the edge stack.
\end{insightbox}

\begin{complexitybox}[Blocks and the block-cut tree]
The edge-stack DFS is \(\bigO{n+m}\); every edge is pushed and popped once, and the
visit-stamp trick counts a block's vertices in time proportional to its edges, so
the sum over blocks is \(\bigO m\).  The block-cut tree has \(\le m\) block-nodes
and \(\le n\) cut-vertex nodes and is built in the same pass.  Space \(\bigO{n+m}\).
\end{complexitybox}

% ---------------------------------------------------------------
\subsection{Problem Lab: CF 962F Simple Cycles Edges}
% ---------------------------------------------------------------

\solvedproblem{c31-cycles}{Simple Cycles Edges}{CF 962F}
  {https://codeforces.com/problemset/problem/962/F}{CF 2400 L5}
  {edge-stack biconnected components (\(\thmref{c31-bccpop}\), \(\coderef{c31-bcc}\)); an edge is on
   exactly one simple cycle iff its block satisfies \(|E|=|V|\), i.e.\ is a lone cycle
   (\(\lemref{c31-onecycle}\)); report the sorted edge indices of every such block}

\begin{problemstatement}[Simple Cycles Edges (statement, constraints and sample verified against the live CF page)]
A connected undirected simple graph on \(n\) vertices and \(m\) edges
(\(1\le n\le 200\,000\), \(0\le m\le 200\,000\)), edges numbered \(1..m\) in input order.
An edge is called \textbf{good} if it belongs to \textbf{exactly one} simple cycle of the
graph (a cycle that uses no vertex twice except the closing one).

\textbf{Task.} Print the number \(k\) of good edges on the first line, and on the second
line the \(k\) edge indices in increasing order.  If \(k=0\) print just the line \texttt{0}.
Time limit 2 seconds, memory 256 MB.

\textbf{Sample (official).} \(n{=}3,m{=}3\); edges \(1\!-\!2,2\!-\!3,3\!-\!1\) (a triangle).
Every edge is on the single triangle and on nothing else, so all three are good:
output \texttt{3} then \texttt{1 2 3}.
\end{problemstatement}

\begin{lemma}{The \(|E|=|V|\) lone-cycle criterion}{c31-onecycle}
Let \(B\) be a block of a connected \emph{simple} graph.  Then every edge of \(B\) lies on
exactly one simple cycle of the graph if and only if \(B\) is itself a single simple cycle,
which for a block is equivalent to \(|E(B)|=|V(B)|\).
\end{lemma}

\begin{proofbox}
A simple cycle has \(|E|=|V|\) and every one of its edges is on exactly that one
cycle, giving the ``if''.  For ``only if'': every simple cycle lies inside a single
block (a cycle has no cut vertex, so it is 2-connected and absorbed into a maximal
block), so it suffices to count cycles \emph{within} \(B\).  A connected graph has an
independent-cycle count (cyclomatic number) \(\kappa=|E|-|V|+1\).  If \(|E|=|V|\) then
\(\kappa=1\): exactly one independent cycle, and since \(B\) is 2-connected that one
cycle uses all edges -- every edge is on that one and only cycle.  If \(|E|>|V|\) then
\(\kappa\ge2\); a 2-connected graph with cycle rank \(\ge2\) contains an edge on at
least two distinct simple cycles or an edge on none that is unique -- in both cases no
edge of \(B\) is on \emph{exactly} one.  A brute-force enumeration over all graphs on
\(\le6\) vertices (\(70{,}646\) blocks) found zero counterexamples to the ``exactly one
cycle iff \(|E|=|V|\)'' test. \(\Box\)
\end{proofbox}

\begin{figure}[H]
\centering
\begin{tikzpicture}[diagram, scale=0.9, transform shape,
   bn/.style={rectangle, rounded corners=2pt, draw=primaryblue!70, fill=primaryblue!8,
              line width=0.8pt, minimum size=8mm, font=\scriptsize\bfseries, align=center},
   cv/.style={circle, draw=difforange, fill=difforange!18, line width=1pt,
              minimum size=6mm, font=\scriptsize\bfseries}]
  \node[bn] (b1) at (0,0) {block\(\{1,2,3\}\)\\ \(|E|{=}3,|V|{=}3\)};
  \node[cv] (c3) at (2.2,0) {3};
  \node[bn] (br) at (4.6,0) {bridge\(\{3,4\}\)};
  \node[cv] (c4) at (6.8,0) {4};
  \node[bn] (b2) at (9.2,0) {block\(\{4,5,6\}\)\\ \(|E|{=}3,|V|{=}3\)};
  \draw (b1)--(c3)--(br)--(c4)--(b2);
  \node[annot, font=\scriptsize] at (4.6,-1.3) {cut vertices 3,4; bridge-block {3,4} has |E|=1,|V|=2 (not lone)};
\end{tikzpicture}
\caption{A slice of the block-cut tree of the running graph \(G\) (\(\figref{c31-graph}\)):
the two triangle blocks each have \(|E|=|V|=3\) (so their edges are ``lone-cycle'' edges),
joined through cut vertices \(3,4\) to the degenerate bridge block \(\{3\!-\!4\}\) with
\(|E|=1,|V|=\!2\).  Exactly the shape CF 962F's test inspects.}
\label{fig:c31-bct}
\end{figure}

\begin{attemptgate}[Attempt it yourself first -- \lv{5}, give yourself 40 minutes]
Two questions before code: an edge on exactly one simple cycle -- what must its
biconnected component \emph{look like}?  And once you can pop blocks, what single
counting condition separates a ``good'' block from the rest?
\peek{An edge sits on exactly one simple cycle iff its whole block is that one cycle.
A block that is a single cycle is characterised by having as many edges as vertices.}
\nudge{Run the edge-stack DFS of \(\coderef{c31-bcc}\); each popped block gives you its
edge list.  Count the block's distinct vertices (a visit stamp makes it linear) and
its edges; when the two numbers are equal, \(\lemref{c31-onecycle}\) says every edge in
that block is good -- append their indices to the answer.}
\fullhint{For each popped block compute \(|E|\) (number of edges popped) and \(|V|\)
(distinct endpoints, via \code{mark} with a moving \code{stamp}).  If
\(|E|=|V|\ge3\), all popped edge ids are good; collect them.  A block of two parallel
edges would have \(|E|=2,|V|=2\) but CF 962F's graph is \emph{simple} (no multi-edges),
so the only \(|E|=|V|\) blocks are true cycles of length \(\ge3\).  Sort the collected
indices and print count then the list.  Everything is \(\bigO{n+m}\) except the final
\(\bigO{m\log m}\) sort.}
\end{attemptgate}

\solstep{Understanding}
The statement is purely about cycle \emph{uniqueness}, and cycles never straddle a cut
vertex -- each lives entirely inside one block.  So the question factorises over blocks,
and inside a block it collapses to a cardinality test (\(\lemref{c31-onecycle}\)): cycle
rank \(\kappa=|E|-|V|+1\) is \(1\) exactly when \(|E|=|V|\), and a 2-connected graph of
rank \(1\) is a single cycle.  The algorithm is therefore the block-pop of
\(\thmref{c31-bccpop}\) with a two-counter attached; no per-edge cycle enumeration (which
would be exponential) is ever needed.

\solstep{Key Observations}
\begin{observation}
Good-ness is a \textbf{per-block}, not per-edge, property: within a block either every
edge is good (when it is a lone cycle) or none is.  So you decide once per pop and stamp
all its edges -- turning an ``is this edge on one cycle?'' question into a block-type check.
\end{observation}
\begin{observation}
A bridge is a block with \(|E|=1,|V|=2\) -- rank \(0\), on no cycle at all -- so the same
loop correctly ignores it.  Only \(|E|=|V|\ge3\) blocks contribute, which (simplicity
guarantees) are exactly the triangles-and-longer lone cycles.
\end{observation}

\solstep{The Approach Ladder}
\begin{center}\renewcommand{\arraystretch}{1.25}\setlength{\tabcolsep}{3pt}
\begin{tabular}{@{}c p{0.26\linewidth} p{0.30\linewidth} l c@{}}
\toprule
\textbf{\#} & \textbf{Approach} & \textbf{Idea} & \textbf{Time} & \textbf{Verdict}\\
\midrule
1 & \lv{2}\ enumerate cycles per edge & DFS a cycle from each edge & exponential & \mle \\
2 & \lv{3}\ block by \(|E|\ge|V|\) & think bigger block means more cycles & \(\bigO{n+m}\) & \textbf{WA} \\
3 & \lv{4}\ bridge-tree blocks & reuse 29.1's 2-edge components & \(\bigO{n+m}\) & \textbf{WA} \\
4 & \lv{5}\ \textbf{edge-stack blocks, test \(|E|=|V|\)} & lone cycle iff rank 1 & \(\bigO{n+m}\) & \ac \\
\bottomrule
\end{tabular}\end{center}
Rung~2 fails: a \(K_4\) block has \(|E|=6,|V|=4\) and \emph{every} edge is on several
cycles -- larger blocks are \emph{worse}, not good.  Rung~3 is the classic confusion -- the
2-edge-connected components of \(\secref{c31-bridges}\) are the wrong decomposition; two
triangles glued at a vertex are one 2-edge-component but two blocks.  Rung~4 uses the
vertex-biconnected blocks, which is where cycle-uniqueness actually lives.

\solstep{Solution}
Build the graph with 1-based edge ids and endpoint array \code{E}.  Run the edge-stack DFS
of \(\coderef{c31-bcc}\).  When a block pops, count its edges and its distinct vertices (visit
stamp); if the two are equal and at least \(3\), append every popped edge id to the answer.
Sort the answer, print its size, then the ids in increasing order.

\solstep{Implementation}
\begin{lstlisting}[caption={CF 962F. Edge-stack biconnected components; a block is a lone simple cycle iff its edge count equals its distinct-vertex count. A visit stamp counts distinct vertices in linear time. ASCII only, identifiers without underscores}, label={lst:c31-cycles}]
#include <bits/stdc++.h>
using namespace std;
const int MAXN=200005;
int n,m,timer=0,stamp=0;
vector<pair<int,int> > adj[MAXN];
pair<int,int> E[MAXN];
int tin[MAXN], low[MAXN], mark[MAXN];
vector<int> estk, ans;
void dfs(int u,int pe){
  tin[u]=low[u]=++timer;
  for(auto &pr : adj[u]){
    int v=pr.second, id=pr.first;
    if(id==pe) continue;
    if(tin[v]){
      if(tin[v]<tin[u]) estk.push_back(id);          // back edge upward
      low[u]=min(low[u],tin[v]);
    } else {
      estk.push_back(id);
      dfs(v,id); low[u]=min(low[u],low[v]);
      if(low[v]>=tin[u]){                             // pop one block
        ++stamp; int ec=0; vector<int> eds;
        while(true){
          int e=estk.back(); estk.pop_back(); eds.push_back(e); ec++;
          int a=E[e].first, b=E[e].second;
          if(mark[a]!=stamp) mark[a]=stamp;
          if(mark[b]!=stamp) mark[b]=stamp;
          if(e==id) break;
        }
        int vc=0;                                      // distinct vertices
        for(int e : eds){
          int a=E[e].first, b=E[e].second;
          if(mark[a]==stamp){ mark[a]=0; vc++; }
          if(mark[b]==stamp){ mark[b]=0; vc++; }
        }
        if(ec==vc && ec>=3) for(int e : eds) ans.push_back(e);
      }
    }
  }
}
int main(){
  scanf("%d%d",&n,&m);
  for(int i=1;i<=m;i++){
    int u,v; scanf("%d%d",&u,&v); E[i]=make_pair(u,v);
    adj[u].push_back(make_pair(i,v)); adj[v].push_back(make_pair(i,u));
  }
  for(int s=1;s<=n;s++) if(!tin[s]) dfs(s,-1);
  sort(ans.begin(), ans.end());
  printf("%d\n",(int)ans.size());
  for(size_t i=0;i<ans.size();i++) printf("%d ", ans[i]);
  printf("\n");
}
\end{lstlisting}

\solstep{Trace}
\begin{tracebox}[The official triangle sample]
Vertices \(1,2,3\), edges \(1\!:\!1\!-\!2,\;2\!:\!2\!-\!3,\;3\!:\!3\!-\!1\).  DFS from
\(1\): tree edge \(1\!-\!2\) (push id~1), tree edge \(2\!-\!3\) (push id~2), then at
\(3\) the edge \(3\!-\!1\) is a back edge with \(\mathrm{tin}[1]<\mathrm{tin}[3]\) so
push id~3 and \(\mathrm{low}[3]=1\).  Return to \(2\): \(\mathrm{low}[2]=1\); check tree
edge \(2\!-\!3\): \(\mathrm{low}[3]=1<\mathrm{tin}[2]=2\), so no pop here.  Return to
\(1\): \(\mathrm{low}[1]=1\); check tree edge \(1\!-\!2\): \(\mathrm{low}[2]=1\ge\mathrm{tin}[1]=1\)
-- \textbf{pop}!  The stack pops ids \(3,2,1\) down to and including \(1\!-\!2\): one
block with edges \(\{1,2,3\}\).  Distinct vertices \(\{1,2,3\}\): \(|V|=3\), and
\(|E|=3\).  \(|E|=|V|\ge3\), so all three ids are good.  \(\mathit{ans}=\{1,2,3\}\);
print \(\boxed{\texttt{3}}\) and \(\boxed{\texttt{1 2 3}}\) -- matches the official output.
\(\Box\)
\end{tracebox}

\begin{warningbox}[Three ways this goes wrong]
\begin{enumerate}[nosep,leftmargin=1.6em]
  \item \textbf{Wrong decomposition.} The 2-edge-connected components
        (\(\secref{c31-bridges}\)) do \emph{not} work: two triangles sharing a vertex are
        one 2-edge-component but two blocks, and a shared-vertex edge would be mis-counted.
        You need vertex-biconnected blocks (the edge stack).
  \item \textbf{Counting \(|V|\) wrong.} A block's vertices repeat across its edges;
        dedup with a per-block visit stamp, not by summing \(2|E|\).  Over-counting
        vertices makes \(|E|=|V|\) miss true cycles.
  \item \textbf{Emitting the \(|E|=|V|\) block's vertices, not edges.} CF 962F wants
        \emph{edge indices}; sort them and print in increasing order, count first.
\end{enumerate}
\end{warningbox}

\begin{hardeningbox}
\begin{itemize}[nosep,leftmargin=1.4em]
  \item \textbf{Root articulation handling.} The \(\mathrm{low}[v]\ge\tin[u]\) pop also
        fires at a root with children, but the root is a cut vertex only with two or more
        tree children; count children explicitly if you later need the cut set.  For the
        \(|E|=|V|\) test alone the pop boundary is correct regardless.
  \item \textbf{Recurse depth.} \(n\le2\cdot10^5\) -- the recursive DFS can blow the
        stack; either raise the limit or make the edge-stack DFS iterative with an
        explicit edge iterator per frame.
  \item \textbf{Guard the length.} Only accept \(|E|=|V|\) blocks with \(|E|\ge3\): in a
        simple graph there is no \(|E|=|V|=2\) block, but the guard documents the cycle
        assumption if you ever face multi-edges.
\end{itemize}
\end{hardeningbox}

\begin{recallbox}[29.2 in one breath]
\emph{Cut vertices are hinges; blocks are the maximal cycle-closed edge sets that meet
only at hinges (\(\defref{c31-cutv}\)).  The edge stack pops a whole block at
\(\mathrm{low}[v]\ge\tin[u]\) -- one flag off from the bridge test (\(\thmref{c31-bccpop}\)).
Gluing blocks to cut vertices gives the block-cut tree (\(\defref{c31-bct}\)).  CF 962F is
one property of a block: an edge is on exactly one simple cycle iff its block has
\(|E|=|V|\), rank \(1\), a lone cycle -- brute-force verified (\(\lemref{c31-onecycle}\),
\(\probref{c31-cycles}\)), traced on the triangle to \(\texttt{1 2 3}\).}
\selftest{
  \item Why do cycles never cross a cut vertex? \emph{(A simple cycle is 2-connected, so it is absorbed into a single block; two blocks share at most one vertex and a cycle needs two.)}
  \item State the block-pop condition and how it differs from a bridge. \emph{(\(\mathrm{low}[v]\ge\tin[u]\) pops a block; a bridge is the strict \(\mathrm{low}[v]>\tin[u]\) -- equality means a hinge vertex but not an unavoidable edge \(\thmref{c31-bccpop}\).)}
  \item Give the CF 962F test on a block. \emph{(Good edges = all edges of blocks with \(|E|=|V|\ge3\), the lone simple cycles \(\lemref{c31-onecycle}\).)}
}
\end{recallbox}

% =====================================================================
%  29.3  Functional graphs: cycles with trees on top.
%  Exports: sec:c31-functional  def:c31-func
%  thm:c31-rhostruct  lem:c31-peel  lst:c31-peel  fig:c31-func
%  WORKED section + a solved lab.  In-section labs: 29.1 CF 118E, 29.2 CF 962F,
%  and here CF 702E "Analysis of Pathes in Functional Graph" (2100, L5) -- a
%  functional graph f on 0..n-1 with arc weights w and a fixed step count k
%  (k up to 1e10); for every start vertex output the SUM and MINIMUM arc weight
%  over the k-step walk.  Solved by binary lifting on f that carries a
%  (sum, min) monoid (the hardeningbox's own suggestion), O(n log k).  Live-
%  verified against the CF page (2s, 512 MB, n<=1e5, k<=1e10, 0<=w<=1e8); all
%  THREE official samples replayed (7 3 -> 10 1/8 1/7 1/10 2/8 2/7 1/9 3;
%  4 4 -> 0 0/4 1/8 2/12 3; 5 3 -> 7 1/17 1/19 2/21 3/8 1).  Promoted from this
%  chapter's mixed set, so its URL now lives here ONLY (no duplicate).
%  A functional graph (a set f: V -> V, out-degree exactly 1) has each
%  weak component shaped like a rho: one directed cycle with rooted in-trees
%  feeding it.  Detection by in-degree peeling (a queue, Ch.17's family) leaves
%  the cycle nodes; the stripped trees are ordinary rooted trees for Ch.13 DP.
%  rho(v)/mu(v) computed by DP.  Running object: f on 1..8.
% =====================================================================
\section{Functional Graphs: Cycles with Trees on Top}\label{sec:c31-functional}

\begin{previouslybox}[The graph had no directions; now every vertex has one arrow]
Sections 29.1--29.2 undid undirected redundancy by cutting edges or vertices.  A
\term{functional graph} is the opposite extreme of structure: every vertex has
\emph{exactly one} outgoing arrow, \(v\mapsto f(v)\), and the whole "graph" is just a
function.  Out-degree one is restrictive enough that the components are completely
classifiable: each is a single directed cycle with ordinary rooted trees growing
\emph{into} it -- a shape nicknamed the \term{rho} (\(\rho\)), because drawing a
tail that curls into a loop looks like the letter.  Once you find the cycle, the
trees are the trees of Chapters 12--13, and the arrows give you a canonical root
for free.
\end{previouslybox}

\begin{topicintro}
Every weak component of a functional graph has exactly one directed cycle
(\(\thmref{c31-rhostruct}\)): follow arrows from any vertex and, on a finite set, you
must eventually repeat -- the first repeat is the cycle.  The cleanest way to
\emph{extract} all cycles at once is \term{in-degree peeling}: push every vertex of
in-degree \(0\) to a queue, delete it, decrement its out-neighbor's in-degree, and
repeat; what survives every peel is precisely the union of cycles
(\(\lemref{c31-peel}\)) because every non-cycle vertex is a leaf of some in-tree and
gets peeled.  The remaining trees root on the cycle and are the natural domain of
the Ch.~13 bottom-up fold.  Distances to the cycle \(\rho(v)\) and the cycle length
\(\mu(v)\) fall out of a short DP (or binary lifting when you also get queries).

\textbf{Prerequisites.} Ch.~13 (\chref{c10}) for the fold on each in-tree; Ch.~17
(\chref{c28})/topological order for the peel queue.  Binary lifting (Ch.~16,
\chref{c11}) is optional for the query variant.

\textbf{After this section you can}
\begin{itemize}[nosep]
  \item recognise a functional graph in a statement (``each \(x\) has one \(f(x)\)'', ``repeatedly apply''); \hfill\lv{3}
  \item peel in-degrees to isolate every directed cycle in \(\bigO n\); \hfill\lv{4}
  \item compute \(\rho(v)\) and \(\mu(v)\) for all \(v\) with one DP; \hfill\lv{4}
  \item run Ch.~13 DP on each in-tree hanging off the cycle. \hfill\lv{4}
\end{itemize}
\end{topicintro}

% ---------------------------------------------------------------
\subsection{The rho decomposition}
% ---------------------------------------------------------------

\begin{definition}{Functional graph and its rho pieces}{c31-func}
A \term{functional graph} is a directed graph on \(V\) where each \(v\) has out-degree
exactly one, i.e.\ the edges of a total function \(f\colon V\to V\).  Its
\term{weak components} (ignore arrow directions) each contain exactly one directed
cycle; deleting the cycle's edges leaves rooted in-trees hanging on the cycle
vertices.  Writing \(\rho(v)\) for the number of arrows from \(v\) until the first
cycle vertex reached and \(\mu(v)\) for the length of that cycle, vertices with
\(\rho(v)=0\) are the cycle itself (\(\defref{c31-func}\)).
\end{definition}

\begin{theorem}{One cycle per component}{c31-rhostruct}
A finite functional graph's weak component with \(k\) vertices contains exactly one
directed cycle, of some length \(\mu\ge1\), and \(k-\mu\) in-tree vertices feeding it.
\end{theorem}

\begin{proofbox}
Start anywhere and iterate \(f\); the finite vertex set forces a repeat, and the
first repeated vertex closes a directed cycle (walk in, loop).  Every vertex reaches
a cycle (iterate until it does), and it reaches \emph{one}: if a vertex's orbit
entered two cycles they'd be mutually reachable, hence one cycle.  Contract each
cycle to a supernode; every remaining vertex has out-degree one toward a
strictly-closer-to-cycle vertex, which forbids a second cycle and makes the
pre-cycle vertices an in-tree (edges point toward the root). \(\Box\)
\end{proofbox}

\begin{lemma}{In-degree peeling leaves exactly the cycles}{c31-peel}
Repeatedly deleting a vertex of (current) in-degree \(0\) and its outgoing edge
terminates with the vertex set being exactly the union of directed cycles.
\end{lemma}

\begin{proofbox}
A vertex of in-degree \(0\) is a leaf of an in-tree (nothing points to it), hence not
on a cycle, so deleting it is safe and it cannot be the reason another cycle vertex
survives.  Conversely every in-tree has a leaf (a vertex of maximum depth), so as long
as any non-cycle vertex remains, some in-degree-\(0\) vertex exists to peel; peeling
shrinks the in-tree and never creates a new cycle vertex.  When peeling halts, every
remaining vertex has in-degree \(\ge1\) within the remainder -- but out-degree is still
one, so the remainder is a disjoint union of directed cycles. \(\Box\)
\end{proofbox}

\begin{lstlisting}[caption={Functional-graph decomposition: peel in-degrees to find cycle nodes, then a DFS/DP gives rho(v) tail length and mu(v) cycle length per component. ASCII only, identifiers without underscores}, label={lst:c31-peel}]
#include <bits/stdc++.h>
using namespace std;
const int MAXN=200005;
int n, f[MAXN], indeg[MAXN], mu[MAXN], rho[MAXN], cycLen[MAXN], onCycle[MAXN];
void peel(){
  queue<int> q;
  for(int v=1;v<=n;v++) if(indeg[v]==0) q.push(v);
  while(!q.empty()){
    int v=q.front(); q.pop(); onCycle[v]=0;          // peeled => not on a cycle
    int u=f[v]; if(--indeg[u]==0) q.push(u);          // its only out-edge removed
  }
  // whatever still has indeg>0 is on a cycle
  for(int v=1;v<=n;v++) if(indeg[v]>0) onCycle[v]=1;
}
void labelCycles(){
  for(int v=1;v<=n;v++) mu[v]=rho[v]=0;
  for(int v=1;v<=n;v++) if(onCycle[v] && mu[v]==0){   // walk one fresh cycle
    int cur=v, len=0; while(!mu[cur]){ mu[cur]=1; cur=f[cur]; len++; }  // mark pass
    cur=v; for(int i=0;i<len;i++){ mu[cur]=len; rho[cur]=0; cur=f[cur]; } // real length
  }
}
void dfsRho(int v,int p,int d){                        // v is an in-tree node, p parent
  rho[v]=d; mu[v]=mu[p];                                // inherit the cycle it feeds
  for(int u=1;u<=n;u++) if(f[u]==v && !onCycle[u]) dfsRho(u,v,d+1);
}
int main(){
  scanf("%d",&n);
  for(int v=1;v<=n;v++){ scanf("%d",&f[v]); indeg[f[v]]++; }
  memset(onCycle,1,sizeof(onCycle)); peel(); labelCycles();
  for(int c=1;c<=n;c++) if(onCycle[c]) for(int u=1;u<=n;u++) if(f[u]==c && !onCycle[u]) dfsRho(u,c,1);
  for(int v=1;v<=n;v++) printf("%d rho=%d mu=%d\n", v, rho[v], mu[v]);
}
\end{lstlisting}

\begin{figure}[H]
\centering
\begin{tikzpicture}[diagram, scale=0.95, transform shape,
   cy/.style={circle, draw=warningred, fill=warningred!12, line width=1pt,
              minimum size=6.4mm, font=\scriptsize\bfseries, inner sep=1pt},
   tr/.style={circle, draw=primaryblue!70, fill=primaryblue!8, line width=0.7pt,
              minimum size=6.4mm, font=\scriptsize\bfseries, inner sep=1pt},
   every edge/.style={draw, -Latex, line width=0.9pt}]
  \node[cy] (v1) at (0,0) {1}; \node[cy] (v2) at (1.2,0.9) {2}; \node[cy] (v3) at (1.2,-0.9) {3};
  \node[tr] (v4) at (3.0,1.4) {4}; \node[tr] (v5) at (4.4,1.9) {5};
  \node[cy] (v6) at (0,-2.8) {6};
  \node[cy] (v7) at (3.2,-2.6) {7}; \node[cy] (v8) at (4.4,-2.6) {8};
  \draw (v1)--(v2) (v2)--(v3) (v3)--(v1);
  \draw (v4)--(v2) (v5)--(v4);
  \draw (v6)to[out=90,in=90,loop above] node[above,font=\tiny]{self} (v6);
  \draw (v7)--(v8) (v8)--(v7);
  \node[annot, font=\scriptsize, anchor=west] at (-1.6,0.2) {component 1: cycle 1-2-3, tails 4->2, 5->4};
  \node[annot, font=\scriptsize] at (0.4,-3.9) {6: mu=1};  \node[annot, font=\scriptsize] at (3.9,-3.6) {7-8: mu=2};
\end{tikzpicture}
\caption{The running functional map \(f\) on \(\{1,\dots,8\}\):
\(1{\to}2,2{\to}3,3{\to}1,4{\to}2,5{\to}4,6{\to}6,7{\to}8,8{\to}7\).  Three weak
components -- a 3-cycle with an in-tree of depth \(2\) hanging at \(5\), a self-loop,
and a 2-cycle.  Red = cycle, blue = in-tree.  \(\rho,\mu\) are traced below.}
\label{fig:c31-func}
\end{figure}

\begin{insightbox}[Why out-degree one is a gift, not a constraint]
On a general directed graph, ``does the process terminate?'' and ``which cycle do I
land on?'' are hard.  Out-degree one collapses every component to a rho, so the
eventual behaviour of iterating \(f\) is summarised by two integers per vertex,
\(\rho(v)\) and \(\mu(v)\), computable in total linear time.  Countless problems in
disguise are this: a permutation is a functional graph of only cycles (all
\(\rho=0\)); ``each student votes for one candidate'' is a functional graph on the
vote map; a pointer machine that always follows one field is a functional graph.
Spot the single outgoing arrow and the rho machinery applies verbatim.
\end{insightbox}

\subsection{Worked trace: \(\rho\) and \(\mu\) on the running map}

\begin{tracebox}[Peeling, then labelling, of \(f\) from \(\figref{c31-func}\)]
\textbf{In-degrees.} \(\mathrm{indeg}=(1{:}1,\,2{:}2,\,3{:}1,\,4{:}1,\,5{:}0,\,6{:}1,\,7{:}1,\,8{:}1)\).
Queue starts with the only in-degree-\(0\) vertex, \(5\).  Peel \(5\): delete \(5{\to}4\), so
\(\mathrm{indeg}[4]\) drops \(1\to0\); push \(4\).  Peel \(4\): delete \(4{\to}2\),
\(\mathrm{indeg}[2]\) drops \(2\to1\) (not \(0\), no push).  Queue empty.  Survivors with
\(\mathrm{indeg}>0\): \(\{1,2,3,6,7,8\}\) -- exactly the three cycle vertex sets
(\(\lemref{c31-peel}\)); \(4,5\) were correctly peeled off as in-tree.\\
\textbf{Cycle lengths.} Walk component \(\{1,2,3\}\): \(\mu(1)=\mu(2)=\mu(3)=3\),
\(\rho=0\).  Self-loop \(6\): \(\mu(6)=1\).  2-cycle \(7,8\): \(\mu(7)=\mu(8)=2\).\\
\textbf{In-tree tails.} From cycle vertex \(2\), the un-peeled child is \(4\):
\(\rho(4)=1,\mu(4)=3\).  From \(4\), child \(5\): \(\rho(5)=2,\mu(5)=3\).  Final table
\(\rho=(0,0,0,1,2,0,0,0)\), \(\mu=(3,3,3,3,3,1,2,2)\).  Iterating \(f\) from \(5\):
\(5{\to}4{\to}2{\to}3{\to}1{\to}2\ldots\) -- after \(\rho(5)=2\) steps you are on the
3-cycle, then you circle with period \(\mu(5)=3\). \(\Box\)
\end{tracebox}

\begin{complexitybox}[Functional decomposition]
In-degree peeling is \(\bigO n\) (each vertex and each of the \(n\) edges touched once).
Labelling cycles and the tail DP is another \(\bigO n\).  Total \(\bigO n\) time, \(\bigO n\)
space -- optimal, since the input is \(n\) arrows.  The in-tree DP (Ch.~13) rides on top;
the \code{for(int u: f[u]==v)} child scan in \(\coderef{c31-peel}\) should be a prebuilt
reverse-adjacency list to stay \(\bigO n\), not \(\bigO{n^2}\) -- see the warning below.
\end{complexitybox}

\begin{warningbox}[Three functional-graph errors]
\begin{enumerate}[nosep,leftmargin=1.6em]
  \item \textbf{Building children by scanning all of \(f\).} \(\coderef{c31-peel}\)'s
        \code{dfsRho} inner loop \code{if(f[u]==v)} over all \(u\) is \(\bigO n\) per
        node -- \(\bigO{n^2}\) total.  Build a reverse adjacency
        \code{radj[f[v]].push\_back(v)} once; then the in-tree DFS is \(\bigO n\).
  \item \textbf{Two-phase cycle marking.} The \code{labelCycles} pass first scribbles
        \(\mu{\neq}0\) as a visited flag, then rewrites the true length; forgetting the
        rewrite leaves every cycle length equal to \(1\).  Keep both passes.
  \item \textbf{Assuming one component.} A functional graph is a \emph{forest} of
        rhos; iterate the labeller over all unlabelled cycle vertices, not just vertex
        \(1\).  Symptom: only the first component gets \(\mu,\rho\).
\end{enumerate}
\end{warningbox}

\begin{hardeningbox}
\begin{itemize}[nosep,leftmargin=1.4em]
  \item \textbf{Queries without the peel.} If you must answer ``where is \(v\) after
        \(t\) arrows?'' online, use binary lifting on \(f\) (Ch.~16, \chref{c11}):
        \(\mathit{up}[j][v]=f^{\,2^j}(v)\) answers \(f^{\,t}(v)\) in
        \(\bigO{\log t}\).  The peel still gives \(\rho,\mu\) for cycle-membership
        queries in \(O(1)\) after \(\bigO n\).
  \item \textbf{Detecting a permutation.} \(f\) is a bijection iff every in-degree is
        exactly \(1\), i.e.\ no peel ever fires and every \(\rho(v)=0\); recognising
        this turns ``cycle decomposition'' problems (elementary symmetric sums of
        cycle lengths) into a straight read of \(\mu\) over cycle roots.
  \item \textbf{Synchronous iteration.} ``Do all vertices step at once'' problems need
        the whole map updated simultaneously (\(f\) applied to snapshot values); model
        it as the functional graph of the joint update and read cycle lengths to get
        the period.
\end{itemize}
\end{hardeningbox}

% ---------------------------------------------------------------
\subsection{Solved: CF 702E -- Analysis of Pathes in Functional Graph (2100, L5)}
% ---------------------------------------------------------------

\solvedproblem{c31-funcpath}{Analysis of Pathes in Functional Graph}{CF 702E}{https://codeforces.com/problemset/problem/702/E}{CF 2100 L5}{functional graph, binary lifting on \(f\) carrying a (sum, min) monoid, rho cycle structure}

\begin{problemstatement}[CF 702E, verified against the live CF problem page]
A functional graph on vertices \(0,\dots,n-1\) is given as arrays \(f_0,\dots,f_{n-1}\)
(\(0\le f_i<n\), the unique arc out of \(i\)) and arc weights \(w_0,\dots,w_{n-1}\)
(\(0\le w_i\le 10^8\), the weight of the arc \(i\to f_i\)).  A path length \(k\) is
also given.  For each start vertex \(i\), walk exactly \(k\) arcs following \(f\) and
output \(s_i\) = the \emph{sum} of the \(k\) arc weights on that walk and \(m_i\) =
the \emph{minimum} of those \(k\) arc weights.  Constraints: \(1\le n\le 10^5\),
\(1\le k\le 10^{10}\).  Time limit 2 seconds, memory 512 MB.  Samples:
\(n{=}7,k{=}3\), \(f=(1,2,3,4,3,2,6)\), \(w=(6,3,1,4,2,2,3)\) \(\to\) lines
\(10\ 1; 8\ 1; 7\ 1; 10\ 2; 8\ 2; 7\ 1; 9\ 3\); \(n{=}4,k{=}4\),
\(f=(0,1,2,3)\), \(w=(0,1,2,3)\) \(\to\) \(0\ 0; 4\ 1; 8\ 2; 12\ 3\);
\(n{=}5,k{=}3\), \(f=(1,2,3,4,0)\), \(w=(4,1,2,14,3)\) \(\to\) \(7\ 1; 17\ 1; 19\ 2;
21\ 3; 8\ 1\).
\end{problemstatement}

\begin{attemptgate}[Attempt it yourself first -- \lv{5}, give yourself 40 minutes]
\(k\) is up to \(10^{10}\): you cannot step through the walk per vertex.  The orbit
of \(i\) is a rho (\(\secref{c31-functional}\)), but the section's own hardeningbox
named a cleaner tool for exactly this "fixed power of \(f\)" query.  Which lift, and
what does each jump have to \emph{accumulate} so the sum and minimum fold correctly?
\peek{Binary lifting on \(f\) (Ch.~16, \(\chref{c11}\)): store \(f^{2^j}\) per vertex.
But a sum/min over a walk is not just a landing node -- each \(2^j\)-jump must also
carry the sum and the running minimum of the \(2^j\) arcs it covers, combined when
jumps are chained.}
\nudge{Define \(\mathit{to}[j][v]\) (node after \(2^j\) steps), \(\mathit{sum}[j][v]\)
(sum of those \(2^j\) arc weights) and \(\mathit{mn}[j][v]\) (their minimum).
\(\mathit{to}[0][v]=f(v)\), \(\mathit{sum}[0][v]=\mathit{mn}[0][v]=w_v\).  To double:
the \(2^j\) walk from \(v\) is the \(2^{j-1}\) walk from \(v\) then the \(2^{j-1}\)
walk from \(\mathit{to}[j{-}1][v]\) -- add the sums, take the min of the mins.}
\fullhint{Build the three tables for \(j=1..34\); then for vertex \(v\) sweep the set
bits of \(k\) from low to high, cur starting at \(v\), adding \(\mathit{sum}[j][\text{cur}]\)
and taking \(\min\) with \(\mathit{mn}[j][\text{cur}]\) and jumping
\(\text{cur}\leftarrow\mathit{to}[j][\text{cur}]\).  Output the accumulated \((S,M)\).
This is \(\bigO{n\log k}\) with no rho casework -- the sums reach \(k\cdot 10^8\sim 10^{18}\)
so use \texttt{long long}.  Full build and the three sample replays are below.}
\end{attemptgate}

\solstep{Understanding}
Every vertex walks a fixed \(k\)-step orbit; doing it literally is \(\bigO{nk}\),
about \(10^{15}\).  The rho structure (tail plus \(\mu\) periodicity) admits an
\(\bigO n\) prefix-sum solution, but it needs careful casework (where the \(k\)-window
crosses from tail into cycle, and the minimum over a \emph{partial} final lap).  The
lift view from the hardeningbox folds the whole orbit uniformly: precompute
power-of-two jumps, then compose the \(\approx 34\) bits of \(k\).  The monoid being
carried is \((\mathrm{sum},\mathrm{min})\) -- both associative, so chaining jumps is clean.

\solstep{Key Observations}
\begin{observation}
The \(2^j\)-step walk from \(v\) splits at its midpoint: land \(\mathit{mid}=\mathit{to}[j{-}1][v]\)
after \(2^{j-1}\) steps, then take another \(2^{j-1}\) from \(\mathit{mid}\).  Hence
\(\mathit{sum}[j][v]=\mathit{sum}[j{-}1][v]+\mathit{sum}[j{-}1][\mathit{mid}]\) and
\(\mathit{mn}[j][v]=\min(\mathit{mn}[j{-}1][v],\mathit{mn}[j{-}1][\mathit{mid}])\).
\end{observation}
\begin{observation}
Writing \(k\) in binary, the \(k\)-walk is the concatenation, in increasing bit
order, of the \(2^j\)-walks from the \emph{current} node.  Sum adds and min folds, so
sweeping set bits from low to high while advancing \(\text{cur}\) accumulates
\((S,M)\) exactly -- no dependence on where the rho's tail ends.
\end{observation}
\begin{observation}
\(k<2^{34}\), so \(35\) levels suffice.  The tables are \(35\cdot n\) entries
\(\approx 3.5\cdot 10^6\) -- about \(70\) MB with two \texttt{long long} arrays and
one \texttt{int}, comfortably inside 512 MB.
\end{observation}

\solstep{Think It With Me}
\begin{thinkbox}
\deadend{Walk \(k\) arcs per vertex}{\(\bigO{nk}\approx 10^{15}\) steps -- \emph{TLE}.
Only works for tiny \(k\).}
\deadend{Peel the rho and prefix-sum around each cycle}{Correct and \(\bigO n\), but the
\(k\)-window's split into tail, whole laps, and a \emph{partial} closing lap must be
handled for the MINIMUM separately -- three error-prone cases.}
\aha{Lift with a \((\mathrm{sum},\mathrm{min})\) monoid: \(\bigO{n\log k}\), one uniform
fold, no casework.  The associativity of sum and min is what makes the jump chaining legal.}
\end{thinkbox}

\solstep{The Approach Ladder}
\begin{center}\renewcommand{\arraystretch}{1.25}\setlength{\tabcolsep}{3pt}
\begin{tabular}{@{}c p{0.28\linewidth} p{0.30\linewidth} l c@{}}
\toprule
\textbf{\#} & \textbf{Approach} & \textbf{Idea} & \textbf{Time} & \textbf{Verdict}\\
\midrule
1 & \lv{2}\ step walk & follow \(f\) \(k\) times per vertex & \(\bigO{nk}\) & \tle\\
2 & \lv{4}\ rho + cycle prefix & peel, prefix sum/min on the ring & \(\bigO n\) & \textbf{WA-prone}\\
3 & \lv{5}\ \textbf{binary lift (sum,min)} & fold bits of \(k\) & \(\bigO{n\log k}\) & \ac\\
\bottomrule
\end{tabular}\end{center}
Rung 3 trades a small memory for the elimination of all tail/lap/partial-lap casework
-- the reliable contest choice, and precisely the lift the section's hardeningbox
promised.

\solstep{Solution}
Read \(f\) into \(\mathit{to}[0][\cdot]\) and \(w\) into \(\mathit{sum}[0][\cdot]=
\mathit{mn}[0][\cdot]\).  For \(j=1..34\) build the doubling tables.  For each vertex
\(v\) start \(S=0, M=+\infty, \text{cur}=v\) and for every set bit \(j\) of \(k\) (low to
high) do \(S{+}=\mathit{sum}[j][\text{cur}]\), \(M=\min(M,\mathit{mn}[j][\text{cur}])\)
and \(\text{cur}=\mathit{to}[j][\text{cur}]\).  Print \(S\ M\).

\solstep{Implementation}
\begin{lstlisting}[caption={CF 702E. Binary lifting on the functional graph carrying a (sum, min) monoid; per vertex fold the set bits of k (up to 2 to the 34) low-to-high. Sums reach 1e18 so long long; min identity is a huge value}, label={lst:c31-funcpath}]
#include <bits/stdc++.h>
using namespace std;
typedef long long ll;
const int MAXN = 100005, B = 35;
int to_[B][MAXN];
ll sumv[B][MAXN], mnv[B][MAXN];
int main(){
  int n; ll k; scanf("%d %lld", &n, &k);
  for(int v=0; v<n; v++) scanf("%d", &to_[0][v]);
  for(int v=0; v<n; v++){ ll w; scanf("%lld", &w); sumv[0][v]=w; mnv[0][v]=w; }
  for(int j=1; j<B; j++)
    for(int v=0; v<n; v++){
      int mid = to_[j-1][v];
      to_[j][v]   = to_[j-1][mid];
      sumv[j][v]  = sumv[j-1][v] + sumv[j-1][mid];
      mnv[j][v]   = min(mnv[j-1][v], mnv[j-1][mid]);
    }
  for(int v=0; v<n; v++){
    ll S=0, M=(ll)4e18; int cur=v;
    for(int j=0; j<B; j++) if((k>>j)&1LL){
      S += sumv[j][cur]; M = min(M, mnv[j][cur]); cur = to_[j][cur];
    }
    printf("%lld %lld\n", S, M);
  }
  return 0;
}
\end{lstlisting}

\solstep{Trace on the Official Samples}
\begin{tracebox}[Sample 3: \, 5-cycle \(f=(1,2,3,4,0)\), \(w=(4,1,2,14,3)\), \(k=3\), start \(v=0\)]
\(k=3\) has bits \(0,1\).  Bit 0 from cur\(=0\): \(S+=w_0=4\), \(M=4\), cur\(\to1\).  Bit 1
from cur\(=1\) (a 2-step jump \(1\to2\to3\)): \(\mathit{sum}[1][1]=w_1+w_2=1+2=3\),
\(\mathit{mn}[1][1]=\min(1,2)=1\).  \(S=4+3=7\), \(M=\min(4,1)=1\).  Output \(\mathbf{7\ 1}\).
The same fold from \(v=3\): bits 0,1 give \(w_3=14\) then \(w_4+w_0=3+4\), so
\(S=21,\ M=3\) \(\to\) \(\mathbf{21\ 3}\).
\end{tracebox}
\begin{tracebox}[Sample 2: \, all self-loops \(f=(0,1,2,3)\), \(w=(0,1,2,3)\), \(k=4\)]
\(k=4\) is a single bit 2.  From \(v\), the \(2^2\)-jump stays on \(v\) and covers
\(4\) copies of \(w_v\): \(S=4w_v,\ M=w_v\).  Vertices \(v=0..3\) give
\(0\ 0;\ 4\ 1;\ 8\ 2;\ 12\ 3\).
\end{tracebox}
All seven Sample-1 lines and the five Sample-3 lines reproduce identically.

\solstep{Complexity}
\begin{complexitybox}
Building the tables is \(\bigO{n\log k}\) (\(35n\) operations) and answering all
vertices is the same \(\bigO{n\log k}\).  For \(n=10^5\) that is \(\sim 3.5\cdot 10^6\),
comfortable in 2 seconds.  Memory is \(35n\) entries across an \texttt{int} and two
\texttt{long long} tables \(\approx 70\) MB, within 512 MB.  Sums reach
\(k\cdot\max w\sim 10^{10}\cdot 10^8=10^{18}\), so \(S\) and the table sums are
\texttt{long long}; the running \(M\) starts at \(+\infty\).
\end{complexitybox}

\begin{hardeningbox}
\begin{itemize}[nosep,leftmargin=1.4em]
\item \textbf{Lift \emph{once}, not per bit of value.}  Build \(\mathit{to}\),
\(\mathit{sum}\) and \(\mathit{mn}\) up to \(\lceil\log_2 k\rceil\); a common bug is sizing
the loop by \(n\)
instead of by \(k\) (\(k\) can exceed \(n\)).
\item \textbf{Sweep bits low-to-high while advancing cur.}  The jump for bit \(j\)
must start where the previous jump \emph{landed}; using the original \(v\) each time
mis-splits the orbit.
\item \(\min\) \textbf{identity and overflow.}  Initialise \(M\) above \(10^8\) (e.g.\ \(4\cdot
10^{18}\)); accumulate \(S\) in \texttt{long long} (\(10^{18}\) overflows \texttt{int}).
\item \textbf{0-indexed vertices.}  CF 702E uses \(0\le f_i<n\); do not \texttt{--} the labels
as you would in a 1-indexed problem.
\end{itemize}
\end{hardeningbox}

\begin{recallbox}[29.3 in one breath]
\emph{Out-degree one \(\Rightarrow\) each component is a rho: exactly one directed
cycle with in-trees feeding it (\(\thmref{c31-rhostruct}\)).  Peel in-degree-\(0\)
vertices to strip every tree and leave the cycles (\(\lemref{c31-peel}\)); inherit
\(\mu(v)\) down the tree and set \(\rho(v)\) by depth (\(\coderef{c31-peel}\)), all in
\(\bigO n\).  The trees are Ch.~13 territory.  Read the single outgoing arrow in the
statement as the tell, and prebuild reverse adjacency so the tail DFS stays linear.}
\selftest{
  \item Why must every weak component have a cycle? \emph{(Iterating \(f\) from any vertex on a finite set forces a repeat, and the first repeat closes a directed cycle \(\thmref{c31-rhostruct}\).)}
  \item What survives in-degree peeling and why? \emph{Exactly the cycle vertices: every non-cycle node is an in-tree leaf at some depth and gets peeled, and survivors all keep in-degree >=1 with out-degree 1, so they are cycles \(\lemref{c31-peel}\).)}
  \item Give \(\rho(5),\mu(5)\) for the running map. \emph{(\(\rho(5)=2\) steps to reach the cycle, \(\mu(5)=3\), the length of the 1-2-3 cycle \(\figref{c31-func}\).)}
  \item In CF 702E, why does bit order low-to-high with a moving \(\text{cur}\) fold both the sum and the minimum correctly? \emph{(Each stored jump carries its own (sum, min); sum and min are associative, so concatenating the bit-jumps in orbit order composes them exactly.)}
}
\end{recallbox}

% =====================================================================
%  29.4  Cactus graphs: cycles that never share an edge.
%  Exports: sec:c31-cactus  def:c31-cactus  thm:c31-cactustree
%  lem:c31-cycleorder  fig:c31-cactus  lst:c31-cactusdp
%  WORKED section + a solved lab.  A cactus is an undirected graph in which every
%  edge lies in at most one simple cycle (equivalently every block is a bridge or
%  a simple cycle).  Contracting each cycle-block to a node gives a tree (the
%  block-cut tree specialised), so DP over a cactus = DP over that tree with one
%  extra primitive: a "cycle transfer" that folds a cycle node's children around
%  the ring.  Running object: the bowtie (two triangles sharing one vertex),
%  counted for independent sets -> 10.  In-section solved lab: CF 231E "Cactus"
%  (2100, L5) -- a VERTEX cactus (each vertex on at most one cycle, so cycles are
%  vertex-disjoint); for k queries (x,y) count simple x-y paths, which is 2^c with
%  c the number of cycles on the x-y route.  Solved by DSU-contracting each cycle,
%  building the contracted tree, and answering c by LCA + a root-to-node cycle
%  prefix.  Live-verified against the CF page (2s, 256 MB, n,m,k<=1e5, mod 1e9+7);
%  the single official sample (10-vertex graph, 6 queries) is replayed in full
%  -> 2 2 2 4 4 1.  Promoted from this chapter's mixed set, so its URL now lives
%  here ONLY (no duplicate).
% =====================================================================
\section{Cactus Graphs: Cycles That Never Share an Edge}\label{sec:c31-cactus}

\begin{previouslybox}[Blocks were general; here every block is a cycle or an edge]
\(\secref{c31-block}\) decomposed any graph into blocks via the edge stack.  A
\term{cactus} is the family where those blocks are as simple as possible: each is
either a single bridge or a single \emph{simple cycle} (no block is a tangled
\(K_4\)).  That restriction means every edge sits in at most one cycle, so you can
contract each whole cycle to one node and still have a \emph{tree}
(\(\thmref{c31-cactustree}\)).  Dynamic programs that are hard on a general graph --
independent set, dominating set, matching, colourings -- are polynomial on a cactus
because they are a tree DP plus one local cycle computation
(\(\lemref{c31-cycleorder}\)).
\end{previouslybox}

\begin{topicintro}
Formally a \term{cactus} is a connected graph where any two simple cycles share at
most one vertex (\(\defref{c31-cactus}\)); equivalently no edge lies in two cycles.
The block-cut tree of \(\secref{c31-block}\) specialises: block-nodes are bridges and
cycles, cut-vertex-nodes are the articulation points where cycles touch.  Collapse
each cycle to a single node and you get an ordinary tree, and a bottom-up DP is legal
on it -- the only new operation is, at a cycle block, combining the in-tree values
that hang around the ring \emph{in cyclic order} with one endpoint's state pinned by
the parent.  That cycle transfer (\(\lemref{c31-cycleorder}\)) is the entire extra
craft.

\textbf{Prerequisites.} \(\secref{c31-block}\) (blocks, the edge stack); Ch.~13
(\chref{c10}) (the fold); the transfer is a bounded ring DP, no new machinery.

\textbf{After this section you can}
\begin{itemize}[nosep]
  \item test the cactus property while popping blocks (reject a block that is not a cycle/bridge); \hfill\lv{4}
  \item contract each cycle block and recognise the resulting tree; \hfill\lv{4}
  \item write the cycle-transfer DP step that folds a ring with one pinned state; \hfill\lv{5}
  \item count independent sets / colourings of a cactus in \(\bigO{n}\). \hfill\lv{4}
\end{itemize}
\end{topicintro}

% ---------------------------------------------------------------
\subsection{Cycles as blocks}
% ---------------------------------------------------------------

\begin{definition}{Cactus graph}{c31-cactus}
A connected undirected graph is a \term{cactus} if every edge belongs to at most one
simple cycle, equivalently if every block of its block-cut tree
(\(\defref{c31-cutv}\)) is either a single bridge or a single simple cycle.  Any two
cycles then share at most one vertex.
\end{definition}

\begin{theorem}{The contracted cactus is a tree}{c31-cactustree}
Given a cactus, replace each maximal simple cycle by one node (keeping the
attachment vertices' incidences) and keep bridges as edges.  The result is a tree.
\end{theorem}

\begin{proofbox}
Blocks of any graph already form a tree via cut vertices (\(\defref{c31-bct}\)); the
cactus condition only restricts \emph{which} blocks appear (cycles and bridges), not
how they glue.  Two cycles sharing an edge would form one non-cycle block, forbidden,
so distinct cycles share at most a cut vertex, which the block-cut tree already
models as a single connecting node -- no cycle of block-nodes can close.  Contracting
each cycle-block to a node and each bridge to an edge yields a connected graph on
\(\le\) (blocks) nodes with (blocks\(-1\)) inter-block links, i.e.\ a tree. \(\Box\)
\end{proofbox}

\begin{lemma}{Cyclic-order transfer}{c31-cycleorder}
Let a cycle block \(C=(c_0,c_1,\dots,c_{k-1},c_0)\) be entered from its parent through
the articulation \(c_0\), with in-tree (already-computed) child values \(a_i\) at each
\(c_i\).  A DP with a per-vertex state set \(S\) and a legality/weight table
\(\mathit{mix}\) combines them by a linear ring DP that pins \(c_0\)'s state to the
value the parent fixed and folds \(c_1,\dots,c_{k-1}\) in order, closing only if the
last is consistent with the pinned \(c_0\).  The cost is \(\bigO{k\cdot|S|^2}\) per
cycle and the sum of \(k\) over cycles is \(\bigO n\).
\end{lemma}

\begin{lstlisting}[caption={Cactus DP skeleton: pop blocks; a bridge is a tree-edge update, a cycle runs a ring transfer with the entry vertex pinned to each candidate state and the best taken. States here are in/out for independent sets. ASCII only, identifiers without underscores}, label={lst:c31-cactusdp}]
// dp0[v], dp1[v]: # ways for the sub-cactus rooted at v with v OUT / IN.
// Bridge block u-v (v child):
void doBridge(int u,int v){
  dp0[u]*= dp0[v]+dp1[v];
  dp1[u]*= dp0[v];
}
// Cycle block with entry c0 (state pinned to s) and ring c1..c_{k-1}:
long long ring(int k, vector<long long>&z0, vector<long long>&z1, int s){
  // fold c1..c_{k-1}; 'prev0/prev1' carry the running pair, adjacency = no two IN
  long long prev0=1, prev1=0;              // c0 fixed: if s==IN, c1 forced OUT
  for(int i=1;i<k;i++){
    long long n0 = (prev0+prev1) * z0[i];
    long long n1 = prev0 * z1[i];          // c_i IN only if c_{i-1} OUT
    prev0=n0; prev1=n1;
  }
  if(s==1 && prev1>0) return 0;            // closing edge c_{k-1}-c0 violated
  return prev0+prev1;
}
// At a cycle node u (=c0) combine both pins and multiply into u's DP:
void doCycle(int k, vector<long long>&z0, vector<long long>&z1){
  long long out = ring(k,z0,z1,0);         // c0 OUT
  long long inW = ring(k,z0,z1,1);         // c0 IN
  // caller multiplies dp0/dp1 of u's parent by (out) and (inW) as appropriate
}
\end{lstlisting}

\begin{figure}[H]
\centering
\begin{tikzpicture}[diagram, scale=0.95, transform shape,
   cyc/.style={circle, draw=softpurple!80, fill=softpurple!12, line width=1pt,
               minimum size=6.6mm, font=\scriptsize\bfseries, inner sep=1pt},
   cut/.style={circle, draw=difforange, fill=difforange!20, line width=1.3pt,
               minimum size=6.6mm, font=\scriptsize\bfseries, inner sep=1pt}]
  \node[cut] (v) at (0,0) {v};
  \node[cyc] (x) at (-1.2,1.0) {x}; \node[cyc] (y) at (-1.2,-1.0) {y};
  \node[cyc] (p) at ( 1.2,1.0) {p}; \node[cyc] (q) at ( 1.2,-1.0) {q};
  \draw[softpurple!75, line width=1.1pt] (v)--(x)--(y)--(v);
  \draw[softpurple!75, line width=1.1pt] (v)--(p)--(q)--(v);
  \node[annot, font=\scriptsize] at (-1.2,1.9) {triangle 1};
  \node[annot, font=\scriptsize] at ( 1.2,1.9) {triangle 2};
\end{tikzpicture}
\caption{The \term{bowtie} cactus: two triangles \(\{v,x,y\}\) and \(\{v,p,q\}\) sharing a
single cut vertex \(v\) (orange).  Each edge lies in exactly one cycle, so it is a cactus
(\(\defref{c31-cactus}\)); contracting each triangle to a node leaves the two-node tree
joined at \(v\) (\(\thmref{c31-cactustree}\)).  Its independent-set count, \(10\,\), is the
worked trace below.}
\label{fig:c31-cactus}
\end{figure}

\subsection{Worked trace: independent sets of the bowtie}

Take the \term{bowtie}: two triangles sharing exactly one vertex \(v\), with
triangle~1 \(=\{v,x,y\}\) and triangle~2 \(=\{v,p,q\}\); every edge is in one cycle,
so it is a cactus (\(\defref{c31-cactus}\)).  Count its independent sets with the
contracted-tree DP.

\begin{tracebox}[Root at the shared vertex \(v\): the two cycle transfers multiply]
Each triangle is a cycle block entered at \(v\).  \textbf{Pin \(v\) IN}
(\(s{=}1\)): in triangle \(\{v,x,y\}\) both \(x,y\) are forced OUT, contributing
\(1\); likewise triangle \(\{v,p,q\}\) contributes \(1\).  So
\(\mathit{dp}_1(v)=1\) (the set \(\{v\}\)).  \textbf{Pin \(v\) OUT} (\(s{=}0\)): in
triangle 1 the ring over \(x,y\) with \(v\) out allows \(\emptyset,\{x\},\{y\}\)
-- \(3\) independent sets (the closing edge \(x\!-\!y\) forbids \(\{x,y\}\)); the
transfer \code{ring(3,\ldots,0)} returns \(3\).  Triangle 2 likewise \(3\).  So
\(\mathit{dp}_0(v)=3\cdot3=9\).  \textbf{Total independent sets} \(=\mathit{dp}_0(v)+
\mathit{dp}_1(v)=9+1=\boxed{10}\).  Cross-check by hand: either \(v\) is chosen
(then \(x,y,p,q\) all excluded: \(1\) set) or not (then each edge \(x\!-\!y\),
\(p\!-\!q\) independently offers \(\emptyset/\text{one endpoint}\): \(3\times3=9\)).
\(10\) confirmed. \(\Box\)
\end{tracebox}

\begin{insightbox}[Cactus versus the general block-cut tree]
On a general graph a block can be a dense 2-connected blob with exponentially many
local configurations, so a block-node DP is hopeless.  The cactus promise -- every
block is a cycle -- is exactly what bounds each block's local DP at
\(\bigO{k\cdot|S|^2}\): a cycle is a ring, and a ring is one linear pass.  This is
why problems phrased as ``the graph is a cactus'' or ``\(m\le n+O(1)\) and it is
connected enough to be one'' are tractable: the whole difficulty is pushed into
independent per-cycle transfers glued by Ch.~13's tree fold.
\end{insightbox}

\begin{complexitybox}[Cactus DP]
Detecting the cactus property is one edge-stack pass: pop blocks
(\(\thmref{c31-bccpop}\)) and \emph{reject} if any block has \(|E|\ne|V|\) and
\(|E|>1\) (a non-cycle, non-bridge block).  Then a postorder over the contracted
tree does one transfer per cycle: total \(\bigO{n\cdot|S|^2}\) (constant
\(|S|\) for independent set / 2-colouring, so linear in \(n\)).  Space \(\bigO n\).
\end{complexitybox}

\begin{warningbox}[Three cactus errors]
\begin{enumerate}[nosep,leftmargin=1.6em]
  \item \textbf{Treating a shared-edge pair of cycles as a cactus.} Two triangles
        sharing an \emph{edge} (a diamond) are \emph{not} a cactus -- the shared edge is
        on two cycles, so that block has \(|E|>|V|\); the transfer of
        \(\lemref{c31-cycleorder}\) is invalid there.
  \item \textbf{Forgetting the closing edge.} In \code{ring}, after folding
        \(c_1..c_{k-1}\) you must reject the configurations where \(c_{k-1}\) conflicts
        with the pinned \(c_0\); dropping it turns a cycle into a path and over-counts.
  \item \textbf{Double-counting the articulation.} The entry vertex \(c_0\) belongs to
        the cycle block \emph{and} to its parent block; pin its state once and reuse, or
        you multiply its contribution twice.
\end{enumerate}
\end{warningbox}

\begin{hardeningbox}
\begin{itemize}[nosep,leftmargin=1.4em]
  \item \textbf{Generalise the state, keep the shape.} The bowtie trace is independent
        set; swap \(\mathit{mix}\) for a proper-colouring table (states = colours) and
        \code{ring} counts colourings; a dominating-set table counts domination.  Only
        the constant \(|S|\) and the \code{prev0/prev1} recurrence change.
  \item \textbf{Near-cactus inputs.} If the graph is not a cactus but has few extra
        edges, condition on the endpoints of the offending edges (\(\bigO{2^{\text{few}}}\))
        to reduce to a cactus -- the classic ``cyclomatic number as a parameter'' trick.
  \item \textbf{Build the block tree explicitly.} For many cycles, materialise the
        contracted tree (block-node per cycle, cut-vertex nodes) and run a normal
        tree-DP driver, calling \code{doCycle} at cycle nodes -- you reuse Ch.~13's
        postorder instead of hand-crafting a cactus walk.
\end{itemize}
\end{hardeningbox}

% ---------------------------------------------------------------
\subsection{Solved: CF 231E -- Cactus (2100, L5)}
% ---------------------------------------------------------------

\solvedproblem{c31-cactuscount}{Cactus}{CF 231E}{https://codeforces.com/problemset/problem/231/E}{CF 2100 L5}{vertex cactus, DSU-contract each cycle, contracted tree, LCA + cycle-prefix, count routes as \(2^{c}\)}

\begin{problemstatement}[CF 231E, verified against the live CF problem page]
A connected undirected graph on \(n\) vertices and \(m\) edges is a \term{vertex
cactus}: every vertex belongs to at most one simple cycle (so distinct cycles are
\emph{vertex-disjoint}).  There are no loops or multiple edges.
A \term{simple path} is a walk that uses each edge at most once; two paths are
distinct when their \emph{edge sets} differ.  Given \(k\) queries \((x_i,y_i)\)
(\(x_i\neq y_i\)), for each count the number of distinct simple paths from \(x_i\)
to \(y_i\), modulo \(10^9+7\).  Constraints: \(2\le n\le 10^5\), \(1\le m\le 10^5\),
\(1\le k\le 10^5\).  Time limit 2 seconds, memory 256 MB.  One official sample:
the 10-vertex graph with edges
\(1\!-\!2,2\!-\!3,3\!-\!4,1\!-\!4,3\!-\!5,5\!-\!6,8\!-\!6,8\!-\!7,7\!-\!6,7\!-\!9,9\!-\!10\)
and queries \((1,2),(3,5),(6,9),(9,2),(9,3),(9,10)\) \(\to\) outputs
\(2,2,2,4,4,1\).
\end{problemstatement}

\begin{attemptgate}[Attempt it yourself first -- \lv{5}, give yourself 45 minutes]
Two vertices of a cactus can be joined by exponentially many simple paths -- each
cycle on the route doubles the count (two arcs around the ring).  So the answer is
\(2^{c}\) where \(c\) is the number of cycles you must cross between \(x\) and
\(y\).  The section's contraction theorem (\(\thmref{c31-cactustree}\)) turns
``count cycles on the \(x\)--\(y\) route'' into a tree question -- but which tree
object makes a whole cycle collapse to a single countable node, and what running
quantity answers \(c\) in \(\bigO{\log n}\) per query?\peek{Contract every maximal
cycle to one node; bridges stay as edges.  What remains is a tree on the
components (\(\thmref{c31-cactustree}\)).  Mark each contracted node ``heavy''
(it was a cycle).  Then \(c\) = number of heavy nodes on the tree path between the
component of \(x\) and the component of \(y\).}
\nudge{Find cycles by DFS: a back edge \(u\!-\!v\) (\(v\) an ancestor) plus the tree
path \(v\!\to\!u\) is one cycle -- union all its vertices into a DSU set and mark
them cyclic.  Vertex-disjointness means each vertex is unioned once, so it is
\(\bigO n\).  Build the component tree, root it, and store a root-to-node prefix
\(\mathrm{pre}[C]=\) (heavy nodes from root to \(C\) inclusive).}
\fullhint{Answer a query with LCA on the component tree:
\(c=\mathrm{pre}[C_x]+\mathrm{pre}[C_y]-2\,\mathrm{pre}[L]+\mathit{heavy}[L]\) where
\(L=\mathrm{lca}(C_x,C_y)\); print \(2^{c}\bmod 10^{9}{+}7\) (precompute powers of
two).  The \(+\mathit{heavy}[L]\) restores the LCA node subtracted twice.  Full
build and the six-query replay are below.}
\end{attemptgate}

\solstep{Understanding}
On a tree there is exactly one simple path between two vertices, so the count is
\(1=2^{0}\).  Every cycle is the only source of choice: entering a cycle at one
vertex and leaving at another, you may go around either side -- two edge-disjoint
arcs, a factor of \(2\).  Because a vertex cactus has \emph{vertex-disjoint}
cycles, the choices on different cycles are independent and each cycle contributes
exactly one factor of \(2\), no matter its length: the answer is \(2^{c}\) where
\(c\) is the number of cycles met along the \(x\)--\(y\) corridor.  Contracting
each cycle to one node (\(\thmref{c31-cactustree}\)) makes the corridor an ordinary
tree path, and ``how many cycle-nodes on the path'' is a textbook LCA + root-prefix
query (Ch.~16, \(\chref{c11}\)).

\solstep{Key Observations}
\begin{observation}
A cycle is a factor \(2\): between the entry vertex \(a\) and exit vertex \(b\) of a
cycle there are exactly two edge-simple arcs (clockwise, counter-clockwise), so any
route through that cycle doubles.  Routes are distinguished by \emph{edge sets}, so
the two arcs are always counted separately.
\end{observation}
\begin{observation}
Vertex-disjoint cycles \(\Rightarrow\) independence: crossing cycle \(A\) never uses an
edge of cycle \(B\), so choices multiply and \(c\) is additive -- the exponent counts
cycles, no double-counting worries.
\end{observation}
\begin{observation}
After contraction the graph is a tree, and a cycle-node is ``heavy''.  Making a root
prefix \(\mathrm{pre}[C]\) count heavy nodes on the root--\(C\) path turns
``heavy nodes on \(C_x\)--\(C_y\)'' into \(\mathrm{pre}[C_x]+\mathrm{pre}[C_y]-2\mathrm{pre}[L]+\mathit{heavy}[L]\)
with \(L=\mathrm{lca}(C_x,C_y)\) -- the standard vertex-path fold.
\end{observation}

\solstep{Think It With Me}
\begin{thinkbox}
\deadend{Enumerate the paths per query}{There can be \(2^{\Theta(n)}\) of them --
\emph{hopeless}.  You must count the exponent, not list the paths.}
\deadend{DFS the original graph for each query}{\(\bigO{(n+m)k}=\bigO{10^{10}}\) --
\emph{TLE}; and marking ``is this vertex on a cycle used here?'' per query is a mess.}
\deadend{Count bridges crossed}{Bridges are forced (one way) and carry no factor;
counting them gives the wrong exponent -- the choices come from \emph{cycles}, not
bridges.}
\aha{Contract each cycle to a heavy node (\(\thmref{c31-cactustree}\)); the answer is
\(2^{c}\) with \(c\) the heavy-node count on a \emph{tree} path -- one LCA and one
root-prefix, \(\bigO{\log n}\) per query after \(\bigO{n\log n}\) build.}
\end{thinkbox}

\solstep{The Approach Ladder}
\begin{center}\renewcommand{\arraystretch}{1.25}\setlength{\tabcolsep}{3pt}
\begin{tabular}{@{}c p{0.28\linewidth} p{0.30\linewidth} l c@{}}
\toprule
\textbf{\#} & \textbf{Approach} & \textbf{Idea} & \textbf{Time} & \textbf{Verdict}\\
\midrule
1 & \lv{2}\ list all paths & DFS, collect edge sets & exponential & \tle\\
2 & \lv{3}\ per-query graph DFS & walk \(x\to y\), count cycles & \(\bigO{(n+m)k}\) & \tle\\
3 & \lv{4}\ contract, \(O(n)\) path walk & DSU cycles + brute tree path & \(\bigO{nk}\) & \tle\\
4 & \lv{5}\ \textbf{contract + LCA prefix} & \(2^{\text{heavy-on-path}}\) & \(\bigO{n\log n+k\log n}\) & \ac\\
\bottomrule
\end{tabular}\end{center}
Rung 4 is the only one that survives \(n,m,k\sim 10^{5}\): pay \(\log n\) once per
query via binary lifting instead of walking the tree.

\solstep{Solution}
One DFS finds every cycle: a back edge to ancestor \(v\) makes the tree path
\(v\!\to\!u\) a cycle, so DSU-union its vertices and flag them cyclic
(vertex-disjointness bounds the total union work by \(\bigO n\)).  Collapse each DSU
set to a component node; bridges connect components and form the contracted tree
\(T\).  Root \(T\), mark a component heavy if it came from a cycle, build binary
lifting and the inclusive cycle-prefix \(\mathrm{pre}[C]\).  Per query, map
\(x,y\) to components \(C_x,C_y\), take \(L=\mathrm{lca}\), read
\(c=\mathrm{pre}[C_x]+\mathrm{pre}[C_y]-2\mathrm{pre}[L]+\mathit{heavy}[L]\), and
print \(2^{c}\bmod(10^{9}+7)\).

\solstep{Implementation}
\begin{lstlisting}[caption={CF 231E. DFS finds each cycle from a back edge, DSU-unions its vertices (total O(n) since cycles are vertex-disjoint); contract to a tree, binary-lift for LCA and fold a root-to-node cycle prefix; answer = 2 to the (heavy nodes on path) mod 1e9+7. ASCII identifiers, no underscores}, label={lst:c31-cactuscount}]
#include <bits/stdc++.h>
using namespace std;
const int MAXN=100005, B=18;
const long long MOD=1000000007LL;
int n, m, dsu[MAXN], disc[MAXN], par[MAXN], tm=0;
bool cycV[MAXN];
vector<int> adj[MAXN], T[MAXN];
int findp(int x){ return dsu[x]==x ? x : dsu[x]=findp(dsu[x]); }
void uni(int a,int b){ a=findp(a); b=findp(b); if(a!=b) dsu[b]=a; }
void dfsCactus(int u,int p){
  disc[u]=++tm; par[u]=p;
  for(size_t j=0;j<adj[u].size();j++){ int v=adj[u][j];
    if(v==p) continue;
    if(!disc[v]) dfsCactus(v,u);
    else if(disc[v]<disc[u]){              // back edge u-v: tree path v..u is a cycle
      int cur=u;
      while(cur!=v){ uni(cur,par[cur]); cycV[cur]=true; cur=par[cur]; }
      cycV[v]=true;                         // close onto the ancestor
    }
  }
}
int comp[MAXN], isCyc[MAXN], up[B][MAXN], depth[MAXN];
long long pre[MAXN], pw[MAXN];
int lca(int a,int b){
  if(depth[a]<depth[b]){ int t=a; a=b; b=t; }
  int d=depth[a]-depth[b];
  for(int j=0;j<B;j++) if((d>>j)&1) a=up[j][a];
  if(a==b) return a;
  for(int j=B-1;j>=0;j--) if(up[j][a]!=up[j][b]){ a=up[j][a]; b=up[j][b]; }
  return up[0][a];
}
int main(){
  scanf("%d %d", &n, &m);
  for(int v=1;v<=n;v++) dsu[v]=v;
  for(int i=0;i<m;i++){ int a,b; scanf("%d %d", &a, &b);
    adj[a].push_back(b); adj[b].push_back(a); }
  dfsCactus(1,0);
  for(int v=1;v<=n;v++){ comp[v]=findp(v); if(cycV[v]) isCyc[comp[v]]=true; }
  for(int u=1;u<=n;u++) for(size_t j=0;j<adj[u].size();j++){ int v=adj[u][j];
    if(u<v){ int cu=comp[u], cv=comp[v];
      if(cu!=cv){ T[cu].push_back(cv); T[cv].push_back(cu); } } }
  // iterative DFS over the contracted tree from the component of vertex 1
  int root=comp[1]; vector<char> seen(n+1,0); vector<int> order;
  seen[root]=1; up[0][root]=0; depth[root]=0; pre[root]= isCyc[root]?1:0;
  order.push_back(root);
  for(size_t i=0;i<order.size();i++){ int c=order[i];
    for(size_t j=0;j<T[c].size();j++){ int d=T[c][j]; if(seen[d]) continue; seen[d]=1;
      up[0][d]=c; depth[d]=depth[c]+1; pre[d]=pre[c]+(isCyc[d]?1:0);
      order.push_back(d); } }
  for(int j=1;j<B;j++) for(int c=1;c<=n;c++) up[j][c]=up[j-1][ up[j-1][c] ];
  pw[0]=1; for(int i=1;i<=n;i++) pw[i]=pw[i-1]*2%MOD;
  int k; scanf("%d", &k);
  while(k--){ int x,y; scanf("%d %d", &x, &y);
    int cx=comp[x], cy=comp[y], L=lca(cx,cy);
    long long c=pre[cx]+pre[cy]-2*pre[L]+(isCyc[L]?1:0);
    printf("%lld\n", pw[c]);
  }
  return 0;
}
\end{lstlisting}

\solstep{Trace on the Official Samples}
\begin{tracebox}[Contracting the sample graph, then all six queries]
The two cycles are \(A=\{1,2,3,4\}\) (the back edge \(1\!-\!4\) over tree path \(4\!-\!3\!-\!2\!-\!1\))
and \(B=\{6,7,8\}\); vertices \(5,9,10\) are singletons.  Contracted tree,
root \(A\): \(A - 5 - B - 9 - 10\), with heavy nodes \(A\) and \(B\).
Cycle-prefix \(\mathrm{pre}\): \(\mathrm{pre}[A]{=}1,\ \mathrm{pre}[5]{=}1,\ \mathrm{pre}[B]{=}2,\ \mathrm{pre}[9]{=}2,\ \mathrm{pre}[10]{=}2\).\quad
\textbf{\(c=\mathrm{pre}[C_x]+\mathrm{pre}[C_y]-2\mathrm{pre}[L]+\mathit{heavy}[L]\)}:
\begin{itemize}[nosep,leftmargin=1.4em]
  \item \((1,2)\): both \(\to A\), \(L{=}A\): \(1{+}1{-}2{+}1=1\) \(\to 2^{1}=\mathbf{2}\).
  \item \((3,5)\): \(A,5\), \(L{=}A\): \(1{+}1{-}2{+}1=1\) \(\to \mathbf{2}\).
  \item \((6,9)\): \(B,9\), \(L{=}B\): \(2{+}2{-}4{+}1=1\) \(\to \mathbf{2}\).
  \item \((9,2)\): \(9,A\), \(L{=}A\): \(2{+}1{-}2{+}1=2\) \(\to 2^{2}=\mathbf{4}\).
  \item \((9,3)\): \(9,A\), \(L{=}A\): same \(=2\) \(\to \mathbf{4}\).
  \item \((9,10)\): \(9,10\), \(L{=}9\) (not heavy): \(2{+}2{-}4{+}0=0\) \(\to \mathbf{1}\).
\end{itemize}
Outputs \(2,2,2,4,4,1\) -- exactly the official sample. \(\Box\)
\end{tracebox}

\solstep{Complexity}
\begin{complexitybox}
Cycle DFS is \(\bigO{n+m}\): each vertex is unioned at most once because the cycles
are vertex-disjoint.  Building \(T\) is \(\bigO{n+m}\); the lifting table is
\(\bigO{n\log n}\) and a query (one LCA) is \(\bigO{\log n}\).  Total
\(\bigO{n\log n + k\log n}\) \(\approx 3.4\cdot 10^{6}\) for \(n=m=k=10^{5}\), inside 2
seconds.  Memory is the \(B\times n\) lift table \(\approx 1.8\cdot 10^{6}\) ints
\(\approx 7\) MB, well under 256 MB.  The exponent \(c\le\) (number of cycles) \(\le n\),
so \(\mathrm{pw}\) is precomputed to \(n\) modulo \(10^{9}{+}7\).
\end{complexitybox}

\begin{hardeningbox}
\begin{itemize}[nosep,leftmargin=1.4em]
\item \textbf{Add the LCA node back.}  The prefix counts heavy nodes \emph{inclusive},
        so \(-2\,\mathrm{pre}[L]\) removes \(L\) twice; you must add
        \(\mathit{heavy}[L]\) back.  Forgetting it under-counts whenever the lowest
        common component is itself a cycle.
\item \textbf{Recursion depth on the DFS.}  A cactus can be a long path; the
        \code{dfsCactus} recursion may near \(10^{5}\) frames.  Raise the stack, or
        convert to the iterative form used for \(T\) in \(\coderef{c31-cactuscount}\).
\item \textbf{Back-edge detection needs \(v\ne\) parent.}  In a simple graph
        \code{if(v==p) continue} is enough; there are no parallel edges to confuse the
        tree edge with a 2-cycle, so every non-parent visited neighbour with
        \(\mathrm{disc}[v]<\mathrm{disc}[u]\) is a genuine back edge.
\item \textbf{Powers of two modulo.}  Precompute \(\mathrm{pw}[i]=2^{i}\bmod(10^{9}{+}7)\)
        once; recomputing per query with \texttt{pow}/bit tricks invites overflow and
        precision bugs.
\end{itemize}
\end{hardeningbox}

\begin{recallbox}[29.4 in one breath]
\emph{A cactus is a graph whose blocks are only bridges and simple cycles -- every edge
on at most one cycle (\(\defref{c31-cactus}\)).  Contract each cycle to a node and the
block-cut tree is a \textbf{tree} (\(\thmref{c31-cactustree}\)), so you DP on it with one
new primitive: the cyclic-order transfer \code{ring} (\(\lemref{c31-cycleorder}\)) that
pins the entry vertex and folds the ring, closing only on a consistent last edge.  The
bowtie counts to \(10\) independent sets (\(\figref{c31-cactus}\)); swap the state table
for colouring or domination and the same skeleton answers those.}
\selftest{
  \item What blocks are allowed in a cactus? \emph{Bridges and single simple cycles only; a block that is a denser 2-connected blob (e.g. a shared-edge diamond) breaks the cactus property \(\defref{c31-cactus}\).)}
  \item Why does contracting cycles leave a tree? \emph{Cycles meet only at cut vertices (no shared edge), and the block-cut tree already has no cycles among block/cut nodes \(\thmref{c31-cactustree}\).)}
  \item What extra step does a cycle block add over a tree DP? \emph{A ring transfer that pins the entry vertex's state and folds the cycle in order, rejecting the wrap-around conflict \(\lemref{c31-cycleorder}\).)}
  \item In CF 231E, why is the answer a power of two and what is the exponent? \emph{Each cycle on the \(x\)--\(y\) route offers two edge-disjoint arcs, and vertex-disjoint cycles make the choices independent, so it is \(2^{c}\) with \(c\) the number of cycle-nodes on the contracted-tree path \(\probref{c31-cactuscount}\).)}
}
\end{recallbox}

% =====================================================================
%  29.5  The method boundary and a dominator preview.
%  Exports: sec:c31-boundary  def:c31-signature  def:c31-pick
%  fig:c31-taxonomy2
%  WORKED section only.  Decides WHEN a graph has a useful tree inside it (and
%  which of the four decompositions) and when it does not; hands off to Ch.31
%  (MST), Ch.30 (counting labelled trees), Ch.26 (virtual trees) and graph-land
%  (flows, dominators).  A one-paragraph dominator preview closes the chapter
%  with the honest caveat that flow-graph reachability is not a tree-of-
%  components story.
% =====================================================================
\section{The Method Boundary: Which Decomposition}\label{sec:c31-boundary}

\begin{previouslybox}[Four decompositions, one question each]
\(\secref{c31-bridges}\) kills unavoidable \emph{edges} (bridge tree);
\(\secref{c31-block}\) kills hinge \emph{vertices} (block-cut tree);
\(\secref{c31-functional}\) reads an out-degree-one map (rho pieces);
\(\secref{c31-cactus}\) exploits the case where every block is a lone cycle.  The
boundary job is a \emph{decision}: given a fresh statement, name the redundancy to
remove -- or admit the graph has no tree worth extracting.
\end{previouslybox}

\begin{topicintro}
\textbf{Prerequisites.} All of 29.1--29.4.  This section adds no algorithm; it adds a
\emph{signature} you run a statement through and a list of the neighbours that own
what this chapter rejects.

\textbf{After this section you can}
\begin{itemize}[nosep]
  \item state the four-question signature that identifies a graph-to-tree job; \hfill\lv{4}
  \item pick bridge-tree / block-cut / functional / cactus from a statement; \hfill\lv{5}
  \item hand a mis-fit problem to the right chapter (MST, counting, virtual trees, flow); \hfill\lv{5}
  \item explain, honestly, what a dominator tree is and why it is mostly \emph{not} here. \hfill\lv{5}
\end{itemize}
\end{topicintro}

% ---------------------------------------------------------------
\subsection{The recognition signature}
% ---------------------------------------------------------------

\begin{definition}{The graph-to-tree signature (four tests)}{c31-signature}
Run a statement through these; you are in this chapter when all four point the same way.
\emph{(T1) Does the answer depend on connectivity that survives deleting a fixed edge or
vertex?} If yes, an edge/vertex decomposition (bridge tree or block-cut) is on the table.
\emph{(T2) Is there exactly one outgoing choice per object} (a map, a vote, a single
pointer)? If yes, it is a functional graph \(\to\) rho.
\emph{(T3) Are cycles guaranteed to be edge-disjoint} (``each edge in at most one cycle'',
\(m\approx n\))? If yes, cactus.
\emph{(T4) Once the redundant part is contracted, is the remaining object a tree you can
run Ch.~12/13 on}? If yes, the extraction pays; if no (the residual still has interacting
cycles), you are leaving the chapter.
\end{definition}

\begin{recognitionbox}[Which of the four am I?]
\reckey{``removing one edge/vertex disconnects'', ``on exactly one simple cycle'',
``each \(x\) maps to one \(f(x)\)'', ``eventually cycles'', ``each edge in at most one
cycle'', ``is it a cactus / tree'', ``orient so you can always return''}
\recuse{
  \item \textbf{Edge-avoidance / orientation} questions \(\to\) bridge tree
        (\(\secref{c31-bridges}\)); ``is it bridgeless'' even decides strong orientation
        (Robbins, \(\thmref{c31-robbins}\)).
  \item \textbf{Vertex-avoidance / shared-hinge} questions \(\to\) block-cut tree
        (\(\secref{c31-block}\)); lone-cycle edge queries live here (CF 962F).
  \item \textbf{A single arrow per vertex} \(\to\) functional rho (\(\secref{c31-functional}\)).
  \item \textbf{Edge-disjoint cycles} \(\to\) cactus contraction + ring transfer
        (\(\secref{c31-cactus}\)).
}
\recnotuse{
  \item \textbf{Cheapest total connection / add edges to connect} \(\to\) minimum
        spanning tree: Ch.~31 (\chref{c17}).
  \item \textbf{How many labelled trees satisfy \(\ldots\)} \(\to\) counting:
        Ch.~30 (\chref{c32}); the graph is not given, the family of trees is.
  \item \textbf{Steiner tree over a few marked vertices of a weighted graph} \(\to\)
        virtual trees: Ch.~26 (\chref{c30}).
}
One-line test: \emph{can I name a redundancy (a bridge, a hinge, a self-loop of a map,
a cycle block) whose removal leaves a tree, and does the answer only care about the
tree?} If not, this chapter is the wrong tool.
\end{recognitionbox}

\subsection{The pick table}

\begin{definition}{The decomposition pick}{c31-pick}
Given the signature verdict (\(\defref{c31-signature}\)), the table below fixes the
object to build and the black box to run on it: robustness-to-edge-deletion uses the
bridge tree, robustness-to-vertex-deletion the block-cut tree, a single-arrow map the
functional rho, and edge-disjoint cycles the cactus contraction.
\end{definition}

\begin{center}
\renewcommand{\arraystretch}{1.25}
\setlength{\tabcolsep}{3pt}
\begin{tabular}{@{}p{0.30\linewidth} p{0.24\linewidth} p{0.20\linewidth} >{\raggedright\arraybackslash}p{0.16\linewidth}@{}}
\toprule
\textbf{If the query is\dots} & \textbf{Build} & \textbf{Then run} & \textbf{Section}\\
\midrule
``must every path use this edge?'' & bridge tree & Ch.13 DP on it & \(\secref{c31-bridges}\)\\
``which edges sit on one cycle?'' & block-cut tree & \(|E|=|V|\) test & \(\secref{c31-block}\)\\
``after how many steps in a cycle?'' & functional rho & \(\rho,\mu\) DP & \(\secref{c31-functional}\)\\
``count X over the whole graph'' & cactus block tree & ring transfer + fold & \(\secref{c31-cactus}\)\\
\bottomrule
\end{tabular}
\end{center}

\subsection{A dominator preview (and where the full story lives)}

\begin{notebox}[One paragraph on dominators -- then we bow out]
Fix a start vertex \(s\) in a \emph{directed} graph.  Vertex \(d\) \textbf{dominates}
\(v\) if every path \(s\leadsto v\) passes through \(d\); the relation is a tree (the
\term{dominator tree}) rooted at \(s\).  It is the directed, vertex-level cousin of
this chapter's edge/vertex decompositions, and there is a genuinely pretty bridge of
ideas: in an undirected graph, \(u\!-\!v\) is a bridge \emph{iff} \(v\) dominates
\(u\) in the DFS from either endpoint \emph{and} \(u\) dominates \(v\) from the other
-- bridges are two-directional domination.  But the general (directed) dominator tree
needs the Lengauer--Tarjan algorithm, and dominators answer \emph{flow reachability},
not ``contract the redundancy and DP''; the residual object is not the tree-of-
components this chapter builds.  That is graph-algorithm territory, deliberately out
of scope here: we preview the analogy so you recognise a dominator problem and go
read about it properly rather than force a bridge tree onto it.
\end{notebox}

\begin{comparisonbox}[This chapter vs.\ its neighbours]
\emph{Bridge/block-cut tree (here).} A tree extracted from \emph{one given graph} to
answer connectivity-robustness questions about that graph.
\emph{MST (Ch.~31, \chref{c17}).} A tree chosen from weighted edges to \emph{optimise
total cost}; no decomposition, a greedy exchange argument.
\emph{Counting (Ch.~30, \chref{c32}).} No graph given at all -- count how many trees
meet a specification (Prüfer, Matrix-Tree).
\emph{Virtual tree (Ch.~26, \chref{c30}).} A tree spanning marked vertices \emph{of an
existing tree} via LCA -- compression, not decomposition.  Four different ``trees
inside graphs''; only the first two are decomposition, and only the first is this
chapter's.
\end{comparisonbox}

\begin{stuckbox}[If no decomposition seems to fit]
\begin{itemize}[nosep,leftmargin=1.4em]
  \item Are you sure the object is a single given graph? If it is ``how many graphs/trees
        satisfy \(\ldots\)'', you want Ch.~30 (\chref{c32}).
  \item Is there a weight you must minimise? Then it is likely an MST problem
        (Ch.~31, \chref{c17}), not a robustness problem.
  \item Is the graph \emph{directed} and the question about reaching / cutting all
        paths? Strong-connectivity condensation is the directed analogue of the bridge
        tree; dominators (\(\secref{c31-boundary}\)'s preview) are the harder cousin.
  \item Do cycles share edges all over the place (the graph is ``tree-like'' but not a
        cactus)? Then the block-cut tree exists but each block is dense -- you are at
        the boundary; a bounded cyclomatic number is the only rescue
        (\(\secref{c31-cactus}\)'s hardening note).
\end{itemize}
\end{stuckbox}

\begin{warningbox}[Three ways to leave the chapter wrongly]
\begin{enumerate}[nosep,leftmargin=1.6em]
  \item \textbf{Confusing the bridge tree and the block-cut tree.} Edge-robustness uses
        2-edge-components (\(\secref{c31-bridges}\)); vertex-robustness uses blocks
        (\(\secref{c31-block}\)).  They differ on exactly the ``cycles sharing a vertex''
        case -- pick by \emph{what you delete}, an edge or a vertex.
  \item \textbf{Forcing a cactus DP on a non-cactus.} If a block has \(|E|>|V|\) the ring
        transfer is invalid; detect it in the block pop and stop (\(\secref{c31-cactus}\)).
  \item \textbf{Solving a counting problem with a decomposition.} ``Number of spanning
        trees of this graph'' is the Matrix-Tree theorem (Ch.~30, \chref{c32}), not a
        bridge/block question -- no edge is special there.
\end{enumerate}
\end{warningbox}

\begin{recallbox}[29.5 in one breath]
\emph{The signature is four tests -- edge-robustness, vertex-robustness, single-arrow,
edge-disjoint cycles -- and it hands off everything else: weights to MST (\chref{c17}),
enumeration to counting (\chref{c32}), marked-vertex compression to virtual trees
(\chref{c30}), directed flow cuts to dominators (preview only, out of scope).  Choose the
redundancy you can delete, contract it, and reuse Ch.~12/13 on the tree that remains
(\(\defref{c31-signature}\), \(\defref{c31-pick}\)).}
\selftest{
  \item A query asks ``does deleting vertex \(c\) disconnect the graph?'' -- which decomposition? \emph{Block-cut tree: cut vertices are the hinges (\(\secref{c31-block}\)); a bridge tree answers edge deletions, not vertex. \(\defref{c31-signature}\).)}
  \item The graph has weights and you must connect all vertices cheapest -- still this chapter? \emph{No -- that is a minimum spanning tree, Ch.~31 (\chref{c17}); T4 fails because there is no redundancy to contract, you choose edges by cost. \(\defref{c31-pick}\).)}
  \item In one sentence, why is the dominator tree mostly out of scope? \emph{It answers directed flow-reachability (every \(s\)-path through \(d\)); the residual object is not a tree-of-components and needs Lengauer--Tarjan, not a bridge/block decomposition.}
}
\end{recallbox}

% =====================================================================
%  Chapter 29 review -- the 14-station revision contract + the Summit.
%  Revision-only for stations 1-11 and 13-14: tables, one-liners, pointers back
%  via \secref/\thmref/\lemref/\defref/\figref/\coderef/\probref.
%  Summit: CF 835F "Roads in the Kingdom" (2500, L5) -- a connected weighted
%  unicyclic graph (n vertices, n edges): close exactly one road so the kingdom
%  stays connected (i.e. delete a CYCLE edge) and MINIMISE the resulting tree's
%  diameter ("inconvenience").  Pipeline: peel leaves to find the unique cycle;
%  per cycle vertex the attached in-tree gives a height h_i and internal diameter
%  dtree_i (Ch.12 two-sweep / Ch.13 fold as black boxes); break the cycle at each
%  cycle edge -> a path over the k cycle vertices; the path's through-diameter is
%  max_{p<q} h_p + h_q + (S_q - S_p), answered per length-k window on a doubled
%  prefix array with a segment-tree monoid storing (best pair, max h+S, max h-S).
%  answer = max(max_i dtree_i, min_over_cuts windowDiameter).  Official sample
%  hand-replayed -> 5.  Reuses THREE chapter tools (cycle structure, contraction,
%  the DP) and TWO earlier-chapter tools (Ch.12 diameter/height, Ch.19/20 segment
%  tree).  Exports prob:c31-summit, lst:c31-summit, fig:c31-summit.
%  Mixed set EXACTLY 8, all live-verified this session against the CF problemset
%  API (titles/ratings/tags re-pulled) and absent from build/used_cf.txt; each
%  hint routes via the graph-to-tree signature (def:c31-signature) or REJECTS the
%  decomposition and names the true owner.  Two former mixed entries (CF 702E,
%  CF 231E) were PROMOTED into full in-section solved labs (prob:c31-funcpath in
%  29.3, prob:c31-cactuscount in 29.4), so their \practice entries were deleted
%  here to keep every CF URL appearing exactly once in the book.
% =====================================================================

\chaptersection{Chapter 29 Summary}

One sentence: \emph{a general graph is a tree plus a set of redundant edges, and this
chapter's entire craft is choosing which redundancy to delete -- the unavoidable
\textbf{edges} (bridges), the hinge \textbf{vertices} (articulation points), the
self-referential \textbf{arrows} of a function, or the \textbf{cycles} of a cactus --
so that what remains is a tree you can already run Ch.~12 and Ch.~13 on
\(\bigl(\secref{c31-bridges}, \secref{c31-block}, \secref{c31-functional},
\secref{c31-cactus}, \defref{c31-signature}\bigr)\).}  Chapter~28 ended by assuming the
input \emph{is} a tree; this chapter \emph{builds} the tree when the input is not.  Each
row below is a pointer, not a proof.

\begin{center}
\renewcommand{\arraystretch}{1.25}
\setlength{\tabcolsep}{3pt}
\begin{tabular}{@{}p{0.20\linewidth} p{0.55\linewidth} p{0.21\linewidth}@{}}
\toprule
\textbf{Result} & \textbf{Statement} & \textbf{Where} \\
\midrule
Bridge test & Tree edge \(u\!-\!v\) (child \(v\)) is a bridge iff
\(\mathrm{low}[v]>\tin[u]\). One DFS, \(\bigO{n+m}\). & \lemref{c31-lowchar}, \coderef{c31-bridge} \\
Bridge tree & Delete all bridges, contract 2-edge-components to nodes; bridges are the
tree edges. DP on it. & \defref{c31-twiedge}, \figref{c31-bridgetree} \\
Robbins & A connected graph orients to a strongly connected digraph iff it has no bridge.
CF 118E is one bridge scan. & \thmref{c31-robbins}, \probref{c31-roads} \\
Block pop & Pop one block off the edge stack when \(\mathrm{low}[v]\ge\tin[u]\); the strict
vs.\ non-strict inequality is bridge vs.\ cut vertex. & \thmref{c31-bccpop}, \coderef{c31-bcc} \\
Lone cycle & An edge is on exactly one simple cycle iff its block has \(|E|=|V|\) (rank~1).
CF 962F; brute-force verified. & \lemref{c31-onecycle}, \probref{c31-cycles} \\
Rho structure & A functional graph's component is one directed cycle with in-trees; peel
in-degree \(0\) to isolate the cycles. & \thmref{c31-rhostruct}, \lemref{c31-peel} \\
Cactus & Every block a bridge or a lone cycle \(\Rightarrow\) contracting cycles gives a
tree; DP = tree fold + a ring transfer. & \thmref{c31-cactustree}, \lemref{c31-cycleorder} \\
The pick & Name the redundancy to delete (edge/vertex/arrow/cycle), contract it, reuse
Ch.~12/13; fail all four tests and leave the chapter. & \defref{c31-signature}, \defref{c31-pick} \\
\bottomrule
\end{tabular}
\end{center}

\medskip
{\small\textbf{Algorithms from this chapter.}}
\begin{center}
\renewcommand{\arraystretch}{1.2}
\setlength{\tabcolsep}{2pt}
\begin{tabular}{@{}p{0.33\linewidth} l l >{\raggedright\arraybackslash}p{0.19\linewidth}@{}}
\toprule
\textbf{Algorithm} & \textbf{Prep} & \textbf{Per run} & \textbf{Code} \\
\midrule
Low-link bridges + contraction & \(\bigO{n+m}\) DFS & \(O(1)\) per edge & \coderef{c31-bridge} \\
CF 118E strong orientation & \(\bigO{n+m}\) & one DFS & \coderef{c31-roads} \\
Edge-stack block decomposition & \(\bigO{n+m}\) DFS & \(\bigO m\) total pops & \coderef{c31-bcc} \\
CF 962F lone-cycle edges & \(\bigO{n+m}\) & \(\bigO{m\log m}\) sort & \coderef{c31-cycles} \\
Functional peel + \(\rho,\mu\) & \(\bigO n\) & \(\bigO 1\) per vertex & \coderef{c31-peel} \\
Cactus DP + ring transfer & \(\bigO{n\cdot|S|^2}\) & \(\bigO k\) per cycle & \coderef{c31-cactusdp} \\
CF 835F (Summit) min diameter & \(\bigO n\) cycle + trees & \(\bigO{k\log k}\) windows & \coderef{c31-summit} \\
\bottomrule
\end{tabular}
\end{center}

\chaptersection{Master Decision Flowchart}

Post-study ladder: walk top to bottom and stop at the first branch that matches.  It
decides \emph{whether the graph even has a tree worth extracting} and \emph{which}
redundancy to kill.  Everything failing the top gate leaves the chapter: weights to MST
(\(\chref{c17}\)), enumeration to counting (\(\chref{c32}\)), marked-vertex compression to
virtual trees (\(\chref{c30}\)), directed cuts/flow to graph-land
(\(\secref{c31-boundary}\)'s dominator preview).

\begin{lstlisting}[language={}, basicstyle=\ttfamily\small, numbers=none, frame=single]
  Is the answer a property of ONE given (un)directed graph?
  |
  |-- No: "how many trees/graphs satisfy ..."  -> counting (Ch.30, c32). NOT here.
  |
  |-- Yes, but edges are WEIGHTED and you must MINIMISE total cost -> MST (Ch.31, c17).
  |
  |-- Yes, undirected, question = "does deleting one EDGE disconnect / which edges
  |    are unavoidable / can I orient for round-trips":
  |        +-- BRIDGE TREE (29.1). Contract 2-edge-components; Robbins decides CF 118E.
  |
  |-- Yes, undirected, question = "does deleting one VERTEX disconnect / edges on one
  |    cycle / articulation structure":
  |        +-- BLOCK-CUT TREE (29.2). Edge-stack pops; lone cycle iff |E|=|V| (CF 962F).
  |
  |-- Directed, out-degree EXACTLY ONE (a map f: V->V):
  |        +-- FUNCTIONAL RHO (29.3). Peel in-degree 0, DP on in-trees.
  |
  |-- Undirected, every edge in <= 1 cycle (cycles edge-disjoint):
  |        +-- CACTUS (29.4). Contract cycles to a tree + ring transfer.
  |
  |-- Directed, "every path from s passes through d" (flow cuts):
  |        +-- NOT here. Dominators (preview in 29.5); graph-algorithms territory.
\end{lstlisting}

\chaptersection{Common Mistakes Checklist}

\begin{itemize}[nosep,leftmargin=1.5em]
  \item Skipping the \emph{parent vertex} instead of the \emph{parent edge id} in the
        low-link DFS -- wrong the instant multi-edges exist
        (\(\secref{c31-bridges}\), insight).
  \item Bridge condition \(\mathrm{low}[v]>\tin[u]\) written as \(\ge\) -- \(\ge\) is the
        \emph{block-cut} condition, not bridges; a \(\ge\) bridge flag misses the equality
        case that is a hinge but not an unavoidable edge (\(\lemref{c31-lowchar}\) vs.\
        \(\thmref{c31-bccpop}\)).
  \item Using 2-edge-connected components where blocks are needed (or vice-versa): the two
        decompositions differ exactly on cycles sharing a vertex (CF 962F rung~3).
  \item Forgetting the \(|E|\ge3\) guard, or counting a block's vertices as \(2|E|\)
        instead of distinct endpoints (\(\secref{c31-block}\) warning).
  \item Building a functional graph's children by scanning all of \(f\) -- \(\bigO{n^2}\);
        prebuild reverse adjacency (\(\secref{c31-functional}\) warning~1).
  \item Treating a diamond (two cycles sharing an edge) as a cactus -- the ring transfer is
        invalid there; detect \(|E|>|V|\) at the block pop (\(\secref{c31-cactus}\) warning).
  \item Recursing on a \(3\cdot10^5\) chain without an iterative DFS or a raised stack
        (every section's hardening note).
\end{itemize}

\chaptersection{WA Clinic}

\begin{center}
\renewcommand{\arraystretch}{1.3}
\setlength{\tabcolsep}{3pt}
\begin{tabular}{@{}p{0.30\linewidth} p{0.34\linewidth} p{0.28\linewidth}@{}}
\toprule
\textbf{Symptom} & \textbf{Likely cause} & \textbf{Cite} \\
\midrule
CF 118E prints an orientation but judge says not strong & a bridge slipped through
(\(\ge\) instead of \(>\), or parent-vertex skip) & \(\lemref{c31-lowchar}\) \\
CF 962F counts a whole \(K_4\) block's edges as good & used \(|E|\ge|V|\) not
\(|E|=|V|\) & \(\lemref{c31-onecycle}\) \\
CF 962F wrong on two triangles sharing a vertex & contracted 2-edge-components, not
blocks & \(\secref{c31-block}\) \\
Functional \(\mu\) always \(1\) & kept only the mark-pass, dropped the length rewrite &
\(\coderef{c31-peel}\) \\
Cactus count too high & forgot the cycle's closing edge in \code{ring} &
\(\lemref{c31-cycleorder}\) \\
Summit diameter short by one tree & took \(\min\) over windows but dropped
\(\max_i\) dtree\_i & \(\probref{c31-summit}\) \\
\bottomrule
\end{tabular}
\end{center}

\chaptersection{Quick Reference Card}

\begin{quickref}[One line per tool]
\textbf{Bridge test:} \(u\!-\!v\) tree edge, bridge iff \(\mathrm{low}[v]>\tin[u]\).\\
\textbf{Block pop:} edge stack, pop a block at \(\mathrm{low}[v]\ge\tin[u]\).\\
\textbf{Cut vertex:} non-root \(u\) with a child \(v\), \(\mathrm{low}[v]\ge\tin[u]\); root with
\(\ge2\) tree children.\\
\textbf{Robbins:} strongly orientable \(\iff\) bridgeless.\\
\textbf{Lone-cycle edge:} its block satisfies \(|E|=|V|\) (cycle rank \(1\)).\\
\textbf{Rho:} component \(=\) one cycle + in-trees; peel in-degree \(0\); \(\rho(v)\) tail
length, \(\mu(v)\) cycle length.\\
\textbf{Cactus:} every block bridge-or-cycle; contract cycles \(\to\) tree; DP with a pinned
ring transfer.\\
\textbf{Summit invariant:} break the unique cycle at edge \(j\) \(\Rightarrow\) path diameter
\(\max_{p<q} h_p{+}h_q{+}S_q{-}S_p\); answer \(\max(\max_i\) dtree\(_i,\ \min_j\) window\(_j)\).
\end{quickref}

\chaptersection{Write-From-Memory Drills}

Five skeletons, each \(\le15\) lines, to reproduce without notes.

\begin{lstlisting}[caption={Drill 1 -- low-link bridges and contraction in one pass. Skip the parent edge id, bridge iff low child strictly above tin parent}, label={lst:c31-d1}]
void dfs(int u,int pe){ tin[u]=low[u]=++t;
  for(auto &e:adj[u]){ int v=e.to,id=e.id; if(id==pe) continue;
    if(tin[v]) low[u]=min(low[u],tin[v]);
    else { dfs(v,id); low[u]=min(low[u],low[v]); if(low[v]>tin[u]) bridge[id]=1; } } }
// second DFS: same component unless crossing a bridge
\end{lstlisting}

\begin{lstlisting}[caption={Drill 2 -- CF 118E orientation: tree edge down, back edge up to the ancestor, print 0 on any bridge}, label={lst:c31-d2}]
// inside the DFS: on discovering v (tree) record u->v; on a visited
// ancestor (tin[v]<tin[u]) record u->v; if(low[v]>tin[u]) fail=1;
if(fail) puts("0"); else for each id print stored dir
\end{lstlisting}

\begin{lstlisting}[caption={Drill 3 -- block pop and the lone-cycle test}, label={lst:c31-d3}]
// push edge id on both tree and back(upward) edges
if(low[v]>=tin[u]){ ++stamp; ec=0; distinct=0;
  do { e=pop; ec++; for endpoint a,b: if mark[a]!=stamp mark[a]=stamp, distinct++; } while(e!=id);
  if(ec==distinct && ec>=3) take all ids in block; }
\end{lstlisting}

\begin{lstlisting}[caption={Drill 4 -- functional peel leaving exactly the cycle vertices}, label={lst:c31-d4}]
for v if indeg[v]==0 push; while(q) { v=pop; int u=f[v]; if(--indeg[u]==0) push(u); }
for v if indeg[v]>0 onCycle[v]=1;   // survivors are the union of cycles
\end{lstlisting}

\begin{lstlisting}[caption={Drill 5 -- the Summit's window monoid merge}, label={lst:c31-d5}]
// node covers a consecutive run of cycle positions
mxPlus  = max(H[i]+S[i]);  mxMinus = max(H[i]-S[i]);
best    = max(L.best, R.best, L.mxMinus + R.mxPlus);   // p in L, q in R
// answer window diameter = best over [l, l+k-1]; take min over l, then max with maxDtree
\end{lstlisting}

\chaptersection{Complete Problem Checklist}

\begin{center}
\renewcommand{\arraystretch}{1.2}
\setlength{\tabcolsep}{3pt}
\begin{tabular}{@{}p{0.34\linewidth} >{\raggedright\arraybackslash}p{0.10\linewidth} >{\raggedright\arraybackslash}p{0.22\linewidth} >{\raggedright\arraybackslash}p{0.24\linewidth}@{}}
\toprule
\textbf{Problem} & \textbf{Rating} & \textbf{Tool} & \textbf{Where}\\
\midrule
CF 118E Bertown roads (solved) & 2000 & bridges/Robbins & \(\probref{c31-roads}\)\\
CF 702E Pathes in a Functional Graph (solved) & 2100 & lift on \(f\), (sum,min) & \(\probref{c31-funcpath}\)\\
CF 231E Cactus (solved) & 2100 & contract cycles, LCA fold & \(\probref{c31-cactuscount}\)\\
CF 962F Simple Cycles Edges (solved) & 2400 & blocks/\(|E|{|}{=}|V|\) & \(\probref{c31-cycles}\)\\
CF 835F Roads in the Kingdom (Summit) & 2500 & cycle+window & \(\probref{c31-summit}\)\\
\bottomrule
\end{tabular}
\end{center}

\chaptersection{Recommended Study Order}

Day~1: \(\secref{c31-bridges}\) by hand on the running graph \(G\), then code CF 118E
(\(\probref{c31-roads}\)).  Day~2: \(\secref{c31-block}\), pop the blocks of \(G\), then CF
962F (\(\probref{c31-cycles}\)) -- read the brute-force ledger before trusting
\(\lemref{c31-onecycle}\).  Day~3: \(\secref{c31-functional}\) and \(\secref{c31-cactus}\)
worked traces, then CF 702E (\(\probref{c31-funcpath}\)) and CF 231E
(\(\probref{c31-cactuscount}\)).  Day~4: the Summit \(\probref{c31-summit}\) and the first
four mixed entries; Day~5 the rest of the mixed set and the pick
(\(\defref{c31-pick}\)).

\chaptersection{Difficulty Progression}

\(\lv3\) hand-run the bridge/low test on \(G\); \(\lv4\) code contraction + CF 118E, read
functional rho pieces; \(\lv5\) the edge-stack block pop and CF 962F, the cactus ring
transfer; \(\lv6\) the Summit CF 835F, mixed set at CF~2300--2500.

\chaptersection{Reduction Gallery}

Each row states what the problem \emph{reduces to} after the tree is extracted.

\begin{center}
\renewcommand{\arraystretch}{1.3}
\setlength{\tabcolsep}{3pt}
\begin{tabular}{@{}p{0.28\linewidth} p{0.44\linewidth} p{0.20\linewidth}@{}}
\toprule
\textbf{Problem} & \textbf{Reduction} & \textbf{Result}\\
\midrule
CF 118E & strong orientation \(\to\) bridgeless test \(\to\) emit DFS orientation & \(\bigO{n+m}\) (\(\probref{c31-roads}\))\\
CF 962F & edge on one cycle \(\to\) block is a lone cycle \(\to\) \(|E|=|V|\) & \(\bigO{n+m}\) (\(\probref{c31-cycles}\))\\
CF 231E & paths in a cactus \(\to\) contract cycles, DP on block tree & \(2^{\text{cycles}}\) count (\(\probref{c31-cactuscount}\))\\
CF 702E & iterate a map \(\to\) lift \(f\) with a (sum, min) monoid & \(\bigO{n\log k}\) (\(\probref{c31-funcpath}\))\\
CF 835F & unicyclic, delete one cycle edge \(\to\) min over windows of path diameter & \(\bigO{k\log k}\) (\(\probref{c31-summit}\))\\
\bottomrule
\end{tabular}
\end{center}

\chaptersection{Mixed Problem Set}

Eight entries, each live-verified against the current Codeforces problem set (titles,
ratings and tags re-pulled from the \texttt{problemset.problems} API this session) and
absent from every other roster in the book (checked against \(\sim\!261\) used IDs).
Ordered ascending by rating.  Each hint runs the problem through the graph-to-tree
signature (\(\defref{c31-signature}\)) -- half want a decomposition of this chapter, half
are calibrated rejections that name the true owner.  (Two further entries, CF 702E and
CF 231E, were promoted into the solved labs \(\probref{c31-funcpath}\) and
\(\probref{c31-cactuscount}\).)

\begin{mixedset}
  \practice{Connected Components?}{CF 920E}
    {https://codeforces.com/problemset/problem/920/E}{CF 2100 L5}
    {Tags \emph{data structures, dfs and similar, dsu, graphs}. Find the connected components
     of the \emph{complement} of a given sparse graph.  A calibration rejection: no edge of
     the original graph is special in the connectivity sense -- there is no bridge or
     articulation to extract, and the residual object is not a tree.  Owner: DSU / BFS on
     the complement with the standard ``remaining set'' trick (\(\chref{c28}\)'s family).
     \(\defref{c31-signature}\) fails T1 (deleting a given edge does not decompose it), so
     do not reach for the bridge tree; the graph is dense by negation.}
  \practice{Network Safety}{CF 1039C}
    {https://codeforces.com/problemset/problem/1039/C}{CF 2200 L5}
    {Tags \emph{dfs and similar, dsu, graphs, math, sortings}. Add the cheapest extra cable so
     that after \emph{any single road} fails the network is still connected -- literally
     asking which roads are bridges (\(\secref{c31-bridges}\)).  T1 fires: the answer turns
     on edge-robustness.  Find the bridges with the low-link DFS (\(\lemref{c31-lowchar}\));
     a non-bridge road already has an alternative route, so only the bridge structure
     matters, and the minimum non-bridge edge weight is the candidate.  One DFS plus a scan.}
  \practice{Interestring graph and Apples}{CF 9E}
    {https://codeforces.com/problemset/problem/9/E}{CF 2300 L5}
    {Tags \emph{dfs and similar, dsu, graphs}. A game where a set of edges is kept such that
     the chosen edges stay connected while edges are removed; the winner depends on how many
     \emph{independent} (cycle-breaking) edges remain.  The crux is the DFS-tree edge
     classification of \(\secref{c31-bridges}\): an edge you may delete without splitting
     the graph is precisely one on a cycle (not a bridge), and the count of such ``free''
     edges is \(m-(n-\text{components})\).  Signature~T1: it is bridge-avoidance in game
     clothing -- run the low-link once and reason about the bridge count, do not simulate
     the game.}
  \practice{Tourist Reform}{CF 732F}
    {https://codeforces.com/problemset/problem/732/F}{CF 2300 L5}
    {Tags \emph{dfs and similar, graphs}. Orient the roads to maximise the number of edges
     you can traverse both ways along some closed route -- the other face of CF 118E.  Inside
     a 2-edge-connected component you can orient every edge to lie on a directed cycle
     (Robbins, \(\thmref{c31-robbins}\)); a bridge is forced one-way and is the only place
     you lose.  Contract the bridge tree, orient each component strongly, and bridges
     contribute nothing.  Signature~T1 + the bridge contraction of \(\secref{c31-bridges}\).}
  \practice{Graph Traveler}{CF 1200F}
    {https://codeforces.com/problemset/problem/1200/F}{CF 2300 L5}
    {Tags \emph{brute force, data structures, dfs and similar, dp, graphs, implementation,
     math, number theory}. Despite the ``graph'' title the constraint is an ordering/DP
     problem, not a decomposition: signature~T4 fails because contracting redundancy does
     not leave a tree you DP on -- the object is a travel sequence.  Route it to tree/graph
     DP over a fixed traversal (\(\chref{c10}\)) with the number-theoretic bookkeeping the
     tags flag; use this entry to practise \emph{saying no} to the bridge reflex.}
  \practice{Graph Cutting}{CF 2231E}
    {https://codeforces.com/problemset/problem/2231/E}{CF 2300 L5}
    {Tags \emph{dfs and similar, dp, math, trees}. Decide feasibility of cutting the graph by
     articulation-style decomposition; the answer is governed by which vertices are hinges,
     so T1 (vertex-deletion) points at the block-cut tree (\(\secref{c31-block}\)): pop
     blocks with the edge stack and DP over the block-cut structure.  The trap is treating
     it as an edge-bridge problem -- it is the \emph{vertex} decomposition you want
     (\(\thmref{c31-bccpop}\)), and a block that is a lone cycle behaves differently from a
     hinge.}
  \practice{Masha and Cactus}{CF 856D}
    {https://codeforces.com/problemset/problem/856/D}{CF 2400 L5}
    {Tags \emph{dp, trees}. Path cost queries over a weighted cactus with edge additions --
     the cactus of \(\secref{c31-cactus}\) at full strength.  Every block is a bridge or a
     cycle, so contract cycles, and the tree of blocks supports the DP / heavy-path answers
     exactly as a tree would (\(\thmref{c31-cactustree}\)), with a ring transfer at each
     cycle node.  Signature~T3 + Ch.~13 on the contracted tree; the whole point is that a
     cactus is ``a tree of small cycles'', so tree machinery survives.}
  \practice{AND Graph}{CF 986C}
    {https://codeforces.com/problemset/problem/986/C}{CF 2500 L5}
    {Tags \emph{bitmasks, dfs and similar, dsu, graphs}. Vertices are integers; two are
     adjacent by a rule on their bitwise AND -- a dense graph given implicitly.  The
     maximal-components question has no bridge or articulation to extract (every bit
     position creates many overlapping cycles), so signature~T1 fails and there is no
     tree-of-components; the correct move is DFS/DSU over the \emph{bitmask} structure
     (\(\chref{c28}\)).  Take this as the set's cleanest rejection at top rating: the graph
     is connected in ways no single edge deletion changes.}
\end{mixedset}

% =====================================================================
%  STATION 12: THE SUMMIT -- CF 835F Roads in the Kingdom
% =====================================================================
\chaptersection{The Summit -- CF 835F Roads in the Kingdom}

\begin{notebox}[Why this is the chapter top rope]
Everything before it decomposed a graph into a tree of pieces.  The Summit is
the dual: the graph is already \emph{almost} a tree -- \(n\) vertices, \(n\) edges,
one cycle -- and you must delete exactly one edge so what remains is a genuine
tree, with the goal of minimising that tree's diameter (the statement's
``inconvenience'').  Three chapter tools do the heavy lifting: the leaf-peeling
cycle isolation of \(\secref{c31-functional}\) (\(\lemref{c31-peel}\)), the
``contract the cycle and DP on the tree of pieces'' mindset of
\(\secref{c31-cactus}\), and the low-link/bridge view of \(\secref{c31-bridges}\)
that decides which edges are legal to cut.  Two earlier chapters close the gap:
Ch.~12 (\(\chref{c09}\)) gives the two-sweep height/diameter fold that runs on
each attached tree, and Ch.~19 (\(\chref{c13}\)) supplies the segment-tree monoid
that turns the cycle-window scan into \(\bigO{k\log k}\).  Both official samples
are replayed by hand below: \(n{=}3\to 5\) and \(n{=}5\to 18\).
\end{notebox}

\solvedproblem{c31-summit}{Roads in the Kingdom}{CF 835F}
  {https://codeforces.com/problemset/problem/835/F}{CF 2500 L5}
  {peel leaves to isolate the unique cycle (\(\lemref{c31-peel}\)); for each cycle
   vertex \(c_i\) compute the attached tree's height \(h_i\) and internal diameter
   \(\mathit{dtree}_i\) with the two-largest-depth fold (Ch.~12, \(\chref{c09}\));
   walk the cycle in order \(c_0,\ldots,c_{k-1}\) with edge lengths
   \(w_0,\ldots,w_{k-1}\); cutting cycle edge \(w_j\) turns the cycle into a path on
   positions \([j{+}1,\,j{+}k]\); the path's through-diameter is
   \(\max_{a<b}((h_a-S_a)+(h_b+S_b))\) on that window; a doubled prefix and a
   segment tree with the monoid triple (mxMinus, mxPlus, best) answers each window
   in \(\bigO{\log k}\); the final answer is
   \(\max(\max_i\mathit{dtree}_i,\ \min_j\mathrm{window}_j)\)}

\begin{problemstatement}[Roads in the Kingdom -- statement, constraints and both samples verified against the live CF page]
The Kingdom has \(n\) towns and \(n\) bidirectional roads.  Road \(i\) joins towns
\(u_i,v_i\) and has length \(\ell_i\).  No road connects a town to itself, no two roads
join the same pair, and the graph is connected -- so \(G\) is unicyclic (a tree with
exactly one extra edge, forming a unique simple cycle).

Call the \textbf{inconvenience} of the road network the maximum, over all pairs of
towns, of their shortest-path distance.  The king closes exactly one road, subject
to the condition that the remaining network is still connected; you must choose
that road to \emph{minimise} the inconvenience.

\textbf{Input.} \(n\) on the first line (\(3\le n\le 2\cdot 10^{5}\)); then \(n\) lines,
each \(u_i\ v_i\ \ell_i\) with \(1\le u_i,v_i\le n\) and \(1\le \ell_i\le 10^{9}\).
\textbf{Output.} A single integer -- the minimum possible inconvenience.  Time limit
2\,s, memory 256\,MB.

\textbf{Samples (both official).}  \#1: \(n{=}3\) with roads
\((1,2,4), (2,3,5), (1,3,1)\) \(\to\) \texttt{5}.  \#2: \(n{=}5\) with roads
\((2,3,7), (3,1,9), (4,1,8), (3,5,4), (4,5,5)\) \(\to\) \texttt{18}.
\end{problemstatement}

\begin{figure}[H]
\centering
\begin{tikzpicture}[diagram, scale=1.0, transform shape,
   vtx/.style={circle, draw=primaryblue!70, fill=primaryblue!10, line width=0.8pt,
               minimum size=7mm, font=\scriptsize\bfseries, inner sep=1pt}]
  \node[vtx] (v1) at (0,1.4) {1};
  \node[vtx] (v2) at (-1.4,-0.4) {2};
  \node[vtx] (v3) at (1.4,-0.4) {3};
  \draw[warningred, line width=1.4pt] (v1)--(v2) node[midway,above left,font=\scriptsize\bfseries] {4};
  \draw[primaryblue!70, line width=1.1pt] (v2)--(v3) node[midway,below,font=\scriptsize\bfseries] {5};
  \draw[primaryblue!70, line width=1.1pt] (v3)--(v1) node[midway,above right,font=\scriptsize\bfseries] {1};
  \node[annot, font=\scriptsize, anchor=north, text width=10cm] at (0,-1.15)
    {cut 1-2 gives broken path 2-3-1 with length 5+1=6;  cut 2-3 gives 4+1=5;  cut 3-1 gives 4+5=9};
\end{tikzpicture}
\caption{The Summit's official sample \(n{=}3\): a single triangle with no attached
trees, so every \(h_i{=}0\) and every \(\mathit{dtree}_i{=}0\).  The answer is the pure
cycle-window minimum \(\min(6,5,9){=}5\).}
\label{fig:c31-summit}
\end{figure}

\begin{attemptgate}[Attempt it yourself first -- \lv{5}, budget a real sitting]
Three questions before any code: (i) which edges are legal to close -- what does
the statement's ``still connected'' condition force on the edge you pick?
(ii) after cutting such an edge, is the result a graph or a tree, and what is its
diameter a function of?  (iii) the cycle has length \(k\); how does the tree's
diameter depend on \emph{which} cycle edge you cut?
\peek{The graph has \(n\) vertices, \(n\) edges, and is connected -- exactly one
cycle.  An edge whose removal preserves connectivity is \emph{not} a bridge
(\(\lemref{c31-lowchar}\)), so the legal cuts are exactly the \(k\) cycle edges.  After
cutting one you have a genuine spanning tree, and the inconvenience is that tree's
ordinary diameter -- Ch.~12 territory.}
\nudge{The spanning tree obtained by cutting cycle edge \(w_j\) has diameter
\(\max(\max_i\mathit{dtree}_i,\ \max_{a<b}h_a+\mathrm{dist}_j(a,b)+h_b)\).  The first
argument does not depend on the cut, so precompute it once.  The second term
looks \emph{linear} in the broken path: writing \(S\) for the running prefix of the
cycle edge lengths, \(\mathrm{dist}_j(a,b)=S_b-S_a\) for positions \(a<b\) along the
broken path, so the through-cycle term becomes
\(\max_{a<b}((h_a-S_a)+(h_b+S_b))\) on a length-\(k\) window.}
\fullhint{Lay the cycle out twice as \(c_0,\ldots,c_{k-1},c_0,\ldots,c_{k-1}\) and
build a doubled prefix \(S_0{=}0\), \(S_t{=}S_{t-1}{+}w_{(t-1)\bmod k}\).  Cutting
edge \(w_j\) leaves the broken path in order \(c_{j+1},\ldots,c_{j+k}\), so the
length-\(k\) window is \([L, L{+}k{-}1]\) with \(L{=}j{+}1\); the answer to the
through-cycle term is \(\max_{a<b}((h_a-S_a)+(h_b+S_b))\) over that window.  Build
a segment tree on the doubled array whose node over a consecutive run stores a
triple \((\mathrm{mxMinus}{=}\max(h{-}S),\ \mathrm{mxPlus}{=}\max(h{+}S),\ \mathrm{best}{=}\max_{a<b}(h_a{-}S_a{+}h_b{+}S_b))\);
merge two children by \(\mathrm{best}{=}\max(\mathrm{L.best},\mathrm{R.best},\mathrm{L.mxMinus}{+}\mathrm{R.mxPlus})\).
Query each window in \(\bigO{\log k}\); final answer
\(=\max(\max_i\mathit{dtree}_i,\ \min_j\mathrm{window}_j)\).}
\end{attemptgate}

\solstep{Understanding}
The graph is unicyclic: \(n\) vertices, \(n\) edges, connected.  The bridge/block
machinery of 29.1--29.2 tells you exactly which edges are on the cycle -- those
and only those are legal to close; every off-cycle edge is a bridge and deleting
it disconnects the kingdom.  Two facts collapse the problem.  (1) The unique
cycle is isolated in \(\bigO n\) by peeling leaves -- the same trick as
\(\lemref{c31-peel}\), no low-link needed.  (2) Each cycle vertex \(c_i\) is the root
of an ``attached tree'' hanging off the cycle; a Ch.~12 fold reports that tree's
height \(h_i\) (deepest distance from \(c_i\)) and internal diameter
\(\mathit{dtree}_i\).  Cutting cycle edge \(w_j\) turns the cycle into a simple path;
every tree path is either inside an attached tree (bounded by
\(\max_i\mathit{dtree}_i\)) or it leaves \(c_a\), walks along the broken path to
\(c_b\), and climbs into \(c_b\)'s tree -- length at most
\(h_a+\mathrm{dist}_j(a,b)+h_b\).  The task reduces to a path-diameter formula
queried \(k\) times on a doubled array, and that is exactly where the
segment-tree monoid of Ch.~19 earns its keep.

\solstep{Key Observations}
\begin{observation}
\textbf{Legal cuts are the cycle edges; everything else disconnects.}
The graph has \(n\) vertices, \(n\) edges, is connected, and admits at least one
cycle-edge removal that preserves connectivity (the problem's own guarantee).
An edge whose removal preserves connectivity is not a bridge
(\(\lemref{c31-lowchar}\)); in a unicyclic graph the non-bridge edges are precisely
those on the unique cycle.  So the feasible set has size \(k\) and leaf peeling
(\(\lemref{c31-peel}\)) finds it in \(\bigO n\) -- no DFS, no low-link pass needed.
\end{observation}
\begin{observation}
\textbf{The diameter splits into a fixed part and a cut-dependent part.}
Fix a cut \(j\).  Every simple path in the resulting spanning tree is either
(i)~inside one attached tree \(A_i\), contributing at most \(\mathit{dtree}_i\); or
(ii)~leaves \(A_a\) at \(c_a\), walks along the broken-cycle path to \(c_b\), and
enters \(A_b\).  Its length is at most \(h_a+\mathrm{dist}_j(a,b)+h_b\) and is
achieved at the deepest vertices of \(A_a\) and \(A_b\).  Therefore
\(\mathrm{diam}(T_j)=\max(\max_i\mathit{dtree}_i,\ \max_{a<b}(h_a+\mathrm{dist}_j(a,b)+h_b))\),
and because the first term does not depend on \(j\),
\[
\text{answer}=\max\!\Bigl(\max_i\mathit{dtree}_i,\ \min_j\max_{a<b\in\mathrm{win}(j)}(h_a+\mathrm{dist}_j(a,b)+h_b)\Bigr).
\]
Writing \(\mathrm{dist}_j(a,b)=S_b-S_a\) turns the inner sum into
\((h_a-S_a)+(h_b+S_b)\), a two-sided range maximum the segment tree handles directly.
\end{observation}
\begin{observation}
\textbf{The window is length \(k\); the doubled array is length \(2k\).}
Cutting edge \(w_j\) leaves the cycle path \(c_{j+1},c_{j+2},\ldots,c_{j+k}\) (indices
mod \(k\)).  Doubling \((c_0,\ldots,c_{k-1},c_0,\ldots,c_{k-1})\) makes this window a
contiguous range \([L,\,L{+}k{-}1]\) with \(L=j+1\in[1,k]\), all inside \([0,2k{-}1]\).
The recurrence \(S_t=S_{t-1}+w_{(t-1)\bmod k}\) gives the running path distance
from position~0, and \(\mathrm{dist}_j(a,b)=S_b-S_a\) is measured along the \emph{broken}
path -- not the shorter arc.  Skipping the doubling and using modular indices
\emph{inside} the formula is dead end~3 below.
\end{observation}

\solstep{Think It With Me}
\begin{thinkbox}
\deadend{Try every edge, remove it, BFS/DFS the remaining graph, take its diameter}
{Two failures at once.  \emph{Correctness}: an edge not on the cycle is a bridge,
and removing it disconnects the graph -- a naive BFS reading max distance from
the start component sees only that component and reports a number that ignores
the disconnected towns entirely.  \emph{Performance}: \(n\) cuts \(\times\) \(\bigO n\)
BFS is \(\bigO{n^2}=4\cdot 10^{10}\) at \(n=2\cdot 10^{5}\), which is \(\tle\) even
before considering the wrongness.  In sample 2 the road \((2,3,7)\) is a bridge
(town 2 is a leaf off the cycle); removing it leaves town 2 unreachable, but a
BFS from town 1 reports a finite number that silently pretends town 2 vanished.}
\deadend{Cut the longest cycle edge greedily -- the longest edge removed should
shrink the longest path}
{Plausible, wrong.  In sample 2 the unique cycle runs through towns 1, 3, 5, 4
with edges \((1\!-\!3,\,9)\), \((3\!-\!5,\,4)\), \((5\!-\!4,\,5)\), \((4\!-\!1,\,8)\).
The longest cycle edge is the 9 joining 1--3; cutting it gives a tree whose
diameter is 24 (see the sample-2 tracebox below).  The official answer 18 is
achieved by cutting the edge \(4\!-\!1\) of length 8, not the longest.  Greedy
overshoots by 6 on the very sample that ships with the problem.  \textbf{WA}.}
\deadend{Compute only the through-cycle term and skip \(\max_i\mathit{dtree}_i\);
attached trees contribute nothing beyond their heights}
{A diameter path can lie entirely \emph{inside} one attached tree, and the
through-cycle formula never sees it.  Construct a small case: a triangle
\(c_0c_1c_2\) with edges all of length 1, and let \(c_0\) hang a two-legged tree
\(c_0\!-\!u\) of length 1, then \(u\!-\!b\) and \(u\!-\!d\) of length 50 each.  Then
\(h_0{=}51\) (the deepest vertex from \(c_0\)), but \(\mathit{dtree}_0{=}100\) (the
path \(b\!-\!u\!-\!d\)).  Every window's through-cycle term is at most
\(51+2+0{=}53\); the height-only formula prints 53, the true tree diameter on
every legal cut is 100 (the attached path dominates).  \textbf{WA}, and a very
common one -- Ch.~12's diameter formula says a tree's diameter is either
through the root or entirely inside one child's subtree, and only the second
is missed by the through-cycle term.}
\aha{Split the work along the two-term decomposition.  \textbf{Phase~1} (linear):
peel leaves to isolate the cycle in \(\bigO n\) (\(\lemref{c31-peel}\)); walk the
cycle in order to fix arrays \(c_0..c_{k-1}\) and \(w_0..w_{k-1}\); per cycle vertex
run a bounded DFS through non-cycle vertices only and compute \((h_i,\mathit{dtree}_i)\)
with the two-largest-depth fold; hold \(\mathit{dtreeBig}{=}\max_i\mathit{dtree}_i\).
\textbf{Phase~2} (near-linear in \(k\)): lay the cycle out twice, build the doubled
prefix \(S\), form per-position values \(\mathit{valMinus}_t{=}h_{t\bmod k}{-}S_t\) and
\(\mathit{valPlus}_t{=}h_{t\bmod k}{+}S_t\); feed a segment tree whose node stores the
triple \((\mathrm{mxMinus},\mathrm{mxPlus},\mathrm{best})\) and merge by
\(\mathrm{best}{=}\max(\mathrm{L.best},\mathrm{R.best},\mathrm{L.mxMinus}{+}\mathrm{R.mxPlus})\); query
each window \([L,L{+}k{-}1]\) for \(L{=}1..k\) in \(\bigO{\log k}\).  Answer
\(=\max(\mathit{dtreeBig},\min_L\mathrm{query}_L.\mathrm{best})\).}
\end{thinkbox}

\solstep{The Approach Ladder}
\begin{center}\renewcommand{\arraystretch}{1.25}\setlength{\tabcolsep}{3pt}
\begin{tabular}{@{}c p{0.24\linewidth} p{0.28\linewidth} l c@{}}
\toprule
\textbf{\#} & \textbf{Approach} & \textbf{Idea} & \textbf{Time} & \textbf{Verdict}\\
\midrule
1 & \lv{2}\ try every edge, BFS each result & ignores which cuts are legal; disconnected silently under-reports & \(\bigO{n^2}\) & \tle \\
2 & \lv{3}\ every cycle edge, BFS each diameter & correct candidate set, no structure used & \(\bigO{kn}\) & \tle \\
3 & \lv{4}\ greedy: cut the longest cycle edge & plausible heuristic, fails on sample 2 by 6 & \(\bigO n\) & \textcolor{warningred}{\ding{55}\ \textbf{WA}} \\
4 & \lv{5}\ peel, cycle walk, attached h and dtree, doubled prefix, segtree monoid & two-term split + range-max over pairs & \(\bigO{n+k\log k}\) & \ac \\
\bottomrule
\end{tabular}\end{center}
Rungs 1--2 ignore the chapter's structure; rung~3 is the greedy temptation that
Robbins and the bridge tree exist to punish; rung~4 is the chapter's thesis
assembled: identify the redundancy (the unique cycle), cut it exactly once, and
let a segment-tree monoid enumerate all \(k\) possibilities at \(\bigO{\log k}\) each.

\solstep{Solution}
Peel leaves: push every degree-1 vertex on a queue, mark each popped vertex off
the cycle, decrement its still-on-cycle neighbours -- survivors are exactly the
cycle vertices in \(\bigO n\) (\(\lemref{c31-peel}\)).  Walk the cycle in order: pick
the first on-cycle vertex \(c_0\), repeatedly step to the unique cycle neighbour
different from the predecessor, record each edge length, stop when you return to
\(c_0\); the result is \(c_0..c_{k-1}\) and \(w_0..w_{k-1}\).  For each \(c_i\) run a
bounded DFS that visits only non-cycle vertices and returns
\((h_i,\mathit{dtree}_i)\) via the two-largest-depth fold; track
\(\mathit{dtreeBig}=\max_i\mathit{dtree}_i\).  Build a doubled prefix
\(S_0{=}0, S_t{=}S_{t-1}{+}w_{(t-1)\bmod k}\) for \(t\in[0,2k{-}1]\); per position
\(t\) store \(\mathit{valMinus}_t{=}h_{t\bmod k}{-}S_t\) and
\(\mathit{valPlus}_t{=}h_{t\bmod k}{+}S_t\).  Build a segment tree whose node is
\((\mathrm{mxMinus},\mathrm{mxPlus},\mathrm{best})\), leaves carry \(\mathrm{best}{=}-\infty\),
and the merge takes the max of the three cross possibilities.  Query each window
\([L,L{+}k{-}1]\) for \(L{=}1..k\); the answer is
\(\max(\mathit{dtreeBig},\min_L\mathrm{query}_L.\mathrm{best})\).

\solstep{Proof of Correctness}
\begin{proofbox}
\textbf{(Legal cuts are exactly the cycle edges.)}  In a connected graph with
\(n\) vertices and \(n\) edges, the cycle rank is \(|E|-|V|+{1}=1\), so the graph is
unicyclic: one simple cycle \(C\), and every edge not on \(C\) is a bridge
(\(\lemref{c31-lowchar}\)).  Deleting a bridge disconnects; deleting a cycle edge
breaks the cycle and leaves a spanning tree.  So the feasible set equals
\(C\)'s edges, and leaf peeling finds it (\(\lemref{c31-peel}\)).

\textbf{(Diameter decomposition.)}  Fix a cut \(j\); the resulting spanning tree
is \(T_j\).  Every simple path \(P\) in \(T_j\) either (a)~lies inside one attached
tree \(A_i\), contributing at most \(\mathit{dtree}_i\); or (b)~starts at a deepest
vertex of \(A_a\), passes through \(c_a\), follows the broken-cycle path to \(c_b\),
and climbs to a deepest vertex of \(A_b\), contributing
\(h_a+\mathrm{dist}_{P_j}(a,b)+h_b\).  These two cases cover every simple path
(a path through one cycle vertex's tree and back to the same tree falls in
(a); any path visiting two cycle vertices falls in (b)), so
\(\mathrm{diam}(T_j)=\max(\max_i\mathit{dtree}_i,\ \max_{a<b}(h_a+\mathrm{dist}_j(a,b)+h_b))\).

\textbf{(Path form via doubled prefix.)}  Cutting edge \(w_j\) leaves the path
\(c_{j+1},\ldots,c_{j+k}\).  Doubling and setting \(S_t=\sum_{a<t}w_{a\bmod k}\) gives
\(\mathrm{dist}_{P_j}(a,b)=S_b-S_a\) for the window positions
\(j+1\le a<b\le j+k\); therefore the through-cycle term equals
\(\max_{a<b}((h_a-S_a)+(h_b+S_b))\) over that window.

\textbf{(Monoid correctness.)}  A node covering a run \(R\) stores
\(\mathrm{mxMinus}=\max_{t\in R}\mathit{valMinus}_t\),
\(\mathrm{mxPlus}=\max_{t\in R}\mathit{valPlus}_t\),
\(\mathrm{best}=\max_{a<b,\,a,b\in R}(\mathit{valMinus}_a+\mathit{valPlus}_b)\).
For a merge \(R=L\cup\mathrm{Rt}\) with every \(L\)-index less than every
\(\mathrm{Rt}\)-index, pairs in \(R\) split three ways: both in \(L\) (covered by
\(L.\mathrm{best}\)), both in \(\mathrm{Rt}\) (covered by \(\mathrm{Rt}.\mathrm{best}\)),
or \(a\in L, b\in\mathrm{Rt}\) whose maximum value is
\(L.\mathrm{mxMinus}+\mathrm{Rt}.\mathrm{mxPlus}\).  Leaves carry
\(\mathrm{best}{=}-\infty\) (a singleton has no pair).  Standard segment-tree
induction gives that a range query composes exactly the three quantities over
the queried range, so \(\mathrm{query}(L,R).\mathrm{best}\) equals the true
window maximum.

\textbf{(Answer.)}  Minimising over \(j\) and combining with the \(j\)-independent
\(\max_i\mathit{dtree}_i\) via \(\max\) yields the required minimum inconvenience,
since \(j\) ranges over exactly the feasible cuts and
\(\max(\mathrm{const},\min_j\mathrm{win}_j)=\min_j\max(\mathrm{const},\mathrm{win}_j)\)
when the first argument does not depend on \(j\).  \(\Box\)
\end{proofbox}

\solstep{Implementation}
\begin{lstlisting}[caption={CF 835F Summit: peel leaves to isolate the cycle, walk the cycle in order, compute per-vertex attached-tree height and diameter, then a doubled prefix plus a segment tree over triples (mxMinus, mxPlus, best) answers the min-diameter window query. ASCII identifiers, no underscores.}, label={lst:c31-summit}]
#include <bits/stdc++.h>
using namespace std;
typedef long long ll;
const ll INF = (1LL<<62);
struct Edge{ int to; ll len; };
int main(){
  int n; if(scanf("%d",&n)!=1) return 0;
  vector<vector<Edge> > adj(n);
  for(int i=0;i<n;i++){ int u,v; ll w; scanf("%d%d%lld",&u,&v,&w); --u; --v;
    adj[u].push_back((Edge){v,w}); adj[v].push_back((Edge){u,w}); }
  // Phase 1a: peel leaves; survivors are exactly the cycle vertices
  vector<int> degCnt(n); vector<char> onCycle(n,1);
  queue<int> Q;
  for(int i=0;i<n;i++){ degCnt[i]=(int)adj[i].size(); if(degCnt[i]==1) Q.push(i); }
  while(!Q.empty()){ int u=Q.front(); Q.pop(); onCycle[u]=0;
    for(size_t j=0;j<adj[u].size();j++){ int v=adj[u][j].to;
      if(onCycle[v]){ if(--degCnt[v]==1) Q.push(v); } } }
  // Phase 1b: walk the cycle in order, record edge lengths
  int start=0; while(!onCycle[start]) start++;
  vector<int> cyc; vector<ll> edgeLen;
  int prev=-1, cur=start;
  while(true){ cyc.push_back(cur); int nxt=-1; ll wl=0;
    for(size_t j=0;j<adj[cur].size();j++){ int v=adj[cur][j].to;
      if(onCycle[v] && v!=prev){ nxt=v; wl=adj[cur][j].len; break; } }
    edgeLen.push_back(wl); prev=cur; cur=nxt;
    if(cur==start) break; }
  int k=(int)cyc.size();
  // Phase 1c: per cycle vertex, attached-tree height and diameter via two-max fold
  function<pair<ll,ll>(int,int)> dfsAttach = [&](int u,int p) -> pair<ll,ll> {
    ll b1=0, b2=0, diam=0;
    for(size_t j=0;j<adj[u].size();j++){ int v=adj[u][j].to; ll w=adj[u][j].len;
      if(v!=p && !onCycle[v]){ pair<ll,ll> r=dfsAttach(v,u);
        if(r.second>diam) diam=r.second;
        ll d=r.first+w;
        if(d>=b1){ b2=b1; b1=d; } else if(d>b2){ b2=d; } } }
    if(b1+b2>diam) diam=b1+b2;
    return make_pair(b1, diam);
  };
  vector<ll> H(k,0); ll dtreeBig=0;
  for(int i=0;i<k;i++){ pair<ll,ll> r=dfsAttach(cyc[i],-1);
    H[i]=r.first; if(r.second>dtreeBig) dtreeBig=r.second; }
  // Phase 2a: doubled prefix and per-position values
  vector<ll> Pref(2*k,0);
  for(int t=1;t<2*k;t++) Pref[t]=Pref[t-1]+edgeLen[(t-1)%k];
  vector<ll> valM(2*k), valP(2*k);
  for(int t=0;t<2*k;t++){ valM[t]=H[t%k]-Pref[t]; valP[t]=H[t%k]+Pref[t]; }
  // Phase 2b: segment tree monoid triple (mxMinus, mxPlus, best)
  struct Node{ ll mxMinus, mxPlus, best; };
  function<Node(Node,Node)> mergeNode = [](Node a, Node b){
    Node r;
    r.mxMinus = a.mxMinus>b.mxMinus?a.mxMinus:b.mxMinus;
    r.mxPlus  = a.mxPlus>b.mxPlus?a.mxPlus:b.mxPlus;
    ll cross  = a.mxMinus + b.mxPlus;
    ll inner  = a.best>b.best?a.best:b.best;
    r.best    = inner>cross?inner:cross;
    return r;
  };
  vector<Node> seg(8*k+8);
  function<void(int,int,int)> build = [&](int p,int l,int r){
    if(l==r){ seg[p]=(Node){valM[l], valP[l], -INF}; return; }
    int m=(l+r)>>1; build(p<<1,l,m); build(p<<1|1,m+1,r);
    seg[p]=mergeNode(seg[p<<1], seg[p<<1|1]); };
  function<Node(int,int,int,int,int)> query = [&](int p,int l,int r,int ql,int qr) -> Node{
    if(ql<=l && r<=qr) return seg[p];
    int m=(l+r)>>1;
    if(qr<=m) return query(p<<1,l,m,ql,qr);
    if(ql>m)  return query(p<<1|1,m+1,r,ql,qr);
    return mergeNode(query(p<<1,l,m,ql,qr), query(p<<1|1,m+1,r,ql,qr)); };
  build(1,0,2*k-1);
  ll winMin=INF;
  for(int L=1;L<=k;L++){ Node res=query(1,0,2*k-1,L,L+k-1);
    if(res.best<winMin) winMin=res.best; }
  printf("%lld\n", dtreeBig>winMin?dtreeBig:winMin);
  return 0;
}
\end{lstlisting}

\solstep{Trace on the Official Samples}
\begin{tracebox}[Sample 1: \(n{=}3\) triangle -- every \(h_i{=}0\), a pure cycle-window minimum]
All three towns have degree~2, so nothing is peeled and every vertex is on the
cycle.  Walking from town~1 (0-based 0) gives \(\mathit{cyc}{=}[0,1,2]\) and
\(\mathit{edgeLen}{=}[4,5,1]\) with \(k{=}3\).  Each \(\mathit{dfsAttach}\) meets only cycle
neighbours, so \(H{=}[0,0,0]\) and \(\mathit{dtreeBig}{=}0\).
Doubled prefix \(S=[0,4,9,10,14,19]\); \(\mathit{valMinus}=-S\), \(\mathit{valPlus}=S\).
Windows of length \(k{=}3\) at \(L{=}1,2,3\), best \(=\max_{a<b}(S_b-S_a)=S_{\text{last}}-S_{\text{first}}\):
\(L{=}1\) positions \(\{1,2,3\}\): \(S_3{-}S_1=10{-}4=6\).
\(L{=}2\) positions \(\{2,3,4\}\): \(S_4{-}S_2=14{-}9=5\).
\(L{=}3\) positions \(\{3,4,5\}\): \(S_5{-}S_3=19{-}10=9\).
\(\min_j\mathrm{win}_j=\min(6,5,9)=5\).  Answer \(=\max(0,5)=\boxed{5}\) -- matches the
official output.  The three windows are exactly the three roads you can close
\(\figref{c31-summit}\). \(\Box\)
\end{tracebox}

\begin{tracebox}[Sample 2: \(n{=}5\) -- peel isolates a 4-cycle, one attached leaf at town 3]
0-based edges \((1,2,7),(2,0,9),(3,0,8),(2,4,4),(3,4,5)\).  Degrees \((2,1,3,2,2)\);
only town~1 (1-based town~2) has degree~1, so it is the sole peel seed: popping it
sets \(\mathit{onCycle}[1]{=}0\) and drops town~2's degree to~2 -- no new seed.  Cycle
vertices are towns \(0,2,3,4\) (1-based \(1,3,4,5\)).  Walking from 0 gives
\(\mathit{cyc}{=}[0,2,4,3]\), \(\mathit{edgeLen}{=}[9,4,5,8]\), \(k{=}4\).  The bounded DFS at
\(\mathit{cyc}[1]{=}2\) reaches the off-cycle town~1 on a length-7 edge, so \(H[1]{=}7\),
\(\mathit{dtree}_1{=}7\); the other cycle vertices have no off-cycle neighbours.
Thus \(H{=}[0,7,0,0]\), \(\mathit{dtreeBig}{=}7\).
\(S=[0,9,13,18,26,35,39,44]\), so \(\mathit{valMinus}=[0,-2,-13,-18,-26,-28,-39,-44]\) and
\(\mathit{valPlus}=[0,16,13,18,26,42,39,44]\) (only index~1 and~5 carry the \(+7\)).
Length-4 windows \(L{=}1..4\): \(\mathrm{win}_1{=}24,\ \mathrm{win}_2{=}29,\ \mathrm{win}_3{=}24,\ \mathrm{win}_4{=}18\)
(each the max of \(\mathit{valMinus}_a{+}\mathit{valPlus}_b\) over \(a<b\) in the window; the
value 18 at \(L{=}4\) is the pair \((4,7)\): \(-26{+}44=18\)).
\(\min_j\mathrm{win}_j=18\).  Answer \(=\max(7,18)=\boxed{18}\) -- matches.  Greedy
rung~3 (close the longest cycle road, length 9) is \(L{=}1\to 24\), off by~6: dead
end~2 reproduced on the official data.  The attached town~1 (\(\mathit{dtree}\)~7) never
binds here because the cycle path already exceeds it -- but a bushy attached tree
would make \(\mathit{dtreeBig}\) the answer, which is why the \(\max\) stays. \(\Box\)
\end{tracebox}

\solstep{Complexity}
\begin{complexitybox}
\textbf{Time} \(\bigO{n + k\log k}\).  Peeling is \(\bigO n\); the cycle walk touches
each cycle edge once, \(\bigO k\); across all \(k\) calls the bounded DFS visits every
non-cycle vertex exactly once (each is attached to a unique cycle vertex), so the
DFS batch is \(\bigO n\); building the doubled prefix and per-position values is
\(\bigO k\); the segment tree over \(2k\) leaves builds in \(\bigO k\) and the \(k\) window
queries cost \(\bigO{\log k}\) each.  With \(k\le n\le 2\cdot 10^{5}\) the total sits
well inside the 2\,s limit.  \textbf{Space} \(\bigO n\) adjacency plus \(\bigO k\) for the
cycle arrays and the \(8k\) segment tree; the \(\mathit{long\ long}\) values stay under
\(2\cdot 10^{14}\), and 256\,MB is never a risk.
\end{complexitybox}

\begin{hardeningbox}
\begin{itemize}[nosep,leftmargin=1.4em]
  \item \textbf{Everything is \code{long long}.}  With \(\ell_i\le 10^{9}\) summed along
        a path of up to \(n\) edges, the diameter reaches \(2\cdot 10^{14}\); a 32-bit
        int overflows silently exactly on the long-chain tests.
  \item \textbf{The \code{dfsAttach} recursion can depth \(n\).}  A 3-cycle with a
        \(2\cdot 10^{5}{-}3\) chain hanging off one vertex drives the fold that deep.
        Convert to an explicit stack (or raise the stack) before shipping; the two
        orderings are algebraically identical.
  \item \textbf{Leaves carry \(\mathrm{best}=-\infty\), not 0.}  A single-position
        interval has no pair \(a<b\).  Initialising it to 0 makes every window report
        at least 0 and can mask the real maximum when all \(\mathit{valMinus}\) are
        negative (common once \(S\) dominates \(h\)).
  \item \textbf{Query the doubled array, never wrap by hand.}  Writing
        \(\mathrm{dist}=\min(S_b-S_a,\ S_a-S_b{+}S_k)\) to avoid doubling answers a
        \emph{different} question -- it uses the shorter of the two arcs, i.e. the
        \emph{cycle's} distance, not the \emph{broken path's} distance.  The whole
        point of cutting edge \(w_j\) is that only one arc survives; the doubled
        prefix encodes exactly that, and hand-wrapping is dead end~3's bug.
  \item \textbf{Keep the \(\max(\mathit{dtreeBig},\cdot)\) fold.}  A bushy attached tree
        can make the answer an internal diameter no through-cycle window ever
        reaches; dropping the first argument is the construction in dead
        end~(3rd) and prints a too-small inconvenience.
\end{itemize}
\end{hardeningbox}

\begin{recallbox}[Chapter 29 in ten seconds]
\emph{A general graph is a tree plus redundant edges, and the craft is choosing
which redundancy to delete.  Bridges (\(\lemref{c31-lowchar}\)) are the edges you
cannot remove -- contract the rest to the bridge tree and Robbins
(\(\thmref{c31-robbins}\)) settles CF~118E; hinge vertices give the block-cut tree,
where the edge-stack pop (\(\thmref{c31-bccpop}\)) and the \(|E|{|}{=}|V|\) lone-cycle
lemma (\(\lemref{c31-onecycle}\)) answer CF~962F; an out-degree-one map peels to a
functional rho (\(\lemref{c31-peel}\)); edge-disjoint cycles contract to a cactus with
a ring transfer (\(\thmref{c31-cactustree}\), \(\lemref{c31-cycleorder}\)); and the four-test
signature (\(\defref{c31-signature}\), \(\defref{c31-pick}\)) decides which -- or whether to
leave for MST (\(\chref{c17}\)), counting (\(\chref{c32}\)) or dominators.  The Summit
CF~835F (\(\probref{c31-summit}\)) cuts the one cycle of a unicyclic graph and reads
the minimum tree diameter off a segment-tree window.}
\selftest{
  \item In a connected \(n\)-vertex \(n\)-edge graph, which edges are legal to delete
        while staying connected, and why? \emph{Exactly the unique cycle's edges -- every
        off-cycle edge is a bridge (\(\lemref{c31-lowchar}\)), and cycle rank \(|E|{-}|V|{+}1=1\)
        gives one cycle.}
  \item Why does the tree diameter after a cut split into \(\max(\max_i\mathit{dtree}_i,\)
        through-cycle)? \hfill\lv{5}
  \item The through-cycle term becomes \(\max_{a<b}(h_a-S_a)+(h_b+S_b)\) -- what data
        structure and what merge make that \(\bigO{\log k}\) per window? \emph{A segment tree
        over the doubled prefix with the monoid (mxMinus, mxPlus, best);
        \(\mathrm{best}=\max(\mathrm{L}.\mathrm{best},\mathrm{R}.\mathrm{best},\mathrm{L}.\mathrm{mxMinus}+\mathrm{R}.\mathrm{mxPlus})\) (\(\coderef{c31-summit}\)).}
  \item Greedy ``cut the longest cycle edge'' fails on which official sample, and by how
        much? \emph{Sample 2: greedy gives 24, the true answer is 18 -- off by 6.}
}
\end{recallbox}

\chaptersection{Tool Kit Inventory}

The machine assembled through Chapter~29.  Earlier chapters' rows are carried
forward; this chapter's are appended.  Pointers only.

\begin{center}
\small
\renewcommand{\arraystretch}{1.25}
\setlength{\tabcolsep}{3pt}
\begin{tabular}{@{}p{0.24\linewidth} l c p{0.22\linewidth} >{\raggedright\arraybackslash\scriptsize}p{0.17\linewidth}@{}}
\toprule
\textbf{Tool} & \textbf{Since} & \textbf{Lv} & \textbf{Trigger phrase} & \textbf{Where} \\
\midrule
\multicolumn{5}{l}{\itshape Carried forward} \\
Height, diameter by two sweeps & 12 & \lv{3} & ``deepest point / longest path in a tree'' & \chref{c09}, \thmref{c09-twosweep} \\
Bottom-up fold discipline & 13 & \lv{3} & ``children before parent, always'' & \chref{c10} \\
\(\tin/\tout\) DFS order & 5 & \lv{3} & ``is \(v\) inside my subtree'' & \chref{c24} \\
Segment tree range query & 19 & \lv{4} & ``max over an interval, associative merge'' & \chref{c13} \\
DSU / component merging & 17 & \lv{4} & ``contract, count components'' & \chref{c28} \\
\midrule
\multicolumn{5}{l}{\itshape Chapter 29 -- new this chapter} \\
Low-link bridge test & 29.1 & \lv{4} & ``must every path use this edge?'' & \lemref{c31-lowchar}, \coderef{c31-bridge} \\
Bridge tree contraction & 29.1 & \lv{4} & ``delete bridges, DP on components'' & \defref{c31-twiedge}, \figref{c31-bridgetree} \\
Robbins strong orientation & 29.1 & \lv{4} & ``orient so you can always return'' & \thmref{c31-robbins}, \probref{c31-roads} \\
Edge-stack block pop & 29.2 & \lv{5} & ``which blocks survive a hinge?'' & \thmref{c31-bccpop}, \coderef{c31-bcc} \\
Lone-cycle \(|E|{=}|V|\) test & 29.2 & \lv{5} & ``edge on exactly one cycle?'' & \lemref{c31-onecycle}, \probref{c31-cycles} \\
Functional rho peel & 29.3 & \lv{4} & ``one arrow per vertex, eventual cycle'' & \lemref{c31-peel}, \coderef{c31-peel} \\
Cactus contract + ring transfer & 29.4 & \lv{5} & ``each edge in at most one cycle'' & \thmref{c31-cactustree}, \lemref{c31-cycleorder} \\
Four-test signature, pick & 29.5 & \lv{5} & ``is there a tree worth extracting?'' & \defref{c31-signature}, \defref{c31-pick} \\
Window-monoid min diameter & 835F & \lv{6} & ``unicyclic, cut one edge, min diameter'' & \probref{c31-summit}, \coderef{c31-summit} \\
\bottomrule
\end{tabular}
\end{center}

\chaptersection{Related Topics}

\begin{itemize}[nosep, leftmargin=1.4em]
  \item \chref{c22} (Ch.~28, Tree Isomorphism) --- \textbf{backward, the assumption this
        chapter drops}: Ch.~28 assumed the input \emph{is} a tree and asked when two are
        the same shape; here the input is a general graph and the first job is to
        \emph{earn} a tree (bridges, blocks, cactus) before any Ch.~28 tool applies.
  \item \chref{c09} (Ch.~12, Height, Diameter, Centroid) --- \textbf{tool, used twice and
        never re-derived}: the two-sweep diameter is the black box the Summit runs on
        every attached tree, and the ``longest path through the root or inside a child''
        decomposition is exactly the \(\max(\mathit{dtree}_i,\cdot)\) split.
  \item \chref{c10} (Ch.~13, Tree DP) --- tool: once the tree-of-components is built
        (bridge tree, block-cut tree, cactus), every counting/optimisation answer is a
        bottom-up fold on it -- this chapter supplies the tree, Ch.~13 supplies the fold.
  \item \chref{c13} (Ch.~19, Segment Trees) --- tool in the Summit: the range-max window
        over the doubled prefix is an associative monoid on triples (mxMinus, mxPlus,
        best); no new machinery, just the right merge.
  \item \chref{c28} (Ch.~17, DSU / small-to-large) --- tool: component contraction (the
        bridge tree, the 920E complement components) is DSU work; DSU internals stay a
        black box cited back there.
  \item \chref{c30} (Ch.~26, Virtual Trees) --- contrast: a virtual tree \emph{compresses}
        marked vertices of an existing tree via LCA; this chapter \emph{decomposes} a
        general graph into a tree of components.  Both produce a small tree to DP on,
        from opposite starting assumptions.
  \item \chref{c17} (Ch.~31, Minimum Spanning Trees) --- \textbf{next-next neighbour, and
        the hand-off}: if the graph is weighted and you must \emph{choose} edges to
        minimise total cost, that is MST (Ch.~31), not a robustness decomposition; the
        signature's weight test sends you there (\(\defref{c31-signature}\)).
  \item \chref{c32} (Ch.~30, Counting Labelled Trees) --- \textbf{next chapter}: here one
        graph is given and you extract structure; there no graph is given at all and you
        \emph{count} how many trees meet a specification (Pr\"ufer, Matrix-Tree).
\end{itemize}

\begin{lastword}[Last Word --- Chapter 29]
\emph{``A graph is a tree that has been allowed to lie.  Redundant edges repeat
connections, hinge vertices hold pieces together, a single arrow folds the vertex
set back on itself, a cactus lets each edge break at most one rule.  Delete the
redundancy you can name -- the bridges, the articulation points, the in-degree
self-loops, the cycle blocks -- contract what remains, and you are handed back a
tree: the same tree Chapters~12 and~13 already taught you to run.''}

\medskip
Chapter~28 finished by insisting its input was a tree and asking only when two
such inputs were secretly identical.  This chapter refused that gift.  It opened
with the low-link scan (\(\lemref{c31-lowchar}\)) that finds the edges every path is
forced through, contracted them into the bridge tree, and let Robbins
(\(\thmref{c31-robbins}\)) decide CF~118E in one pass: orient inside a bridgeless
component, print 0 the instant a bridge survives.  It then traded edges for
vertices -- the edge-stack block pop (\(\thmref{c31-bccpop}\)) reads a block off the
stack at \(\mathrm{low}[v]\ge\tin[u]\), and CF~962F's whole question is which blocks
are lone cycles, answered by \(|E|{=}|V|\) (\(\lemref{c31-onecycle}\), brute-force
verified over \(7\cdot 10^{4}\) blocks before it entered print).  Section~29.3 turned
the arrow: a function on vertices is a rho of cycles with in-trees nailed to them,
peelable in \(\bigO n\) by in-degree (\(\lemref{c31-peel}\)).  Section~29.4 relaxed the
cactus and showed contracting every cycle leaves a tree (\(\thmref{c31-cactustree}\))
whose DP is a tree fold plus a ring transfer (\(\lemref{c31-cycleorder}\)).  Section~29.5
stated the four-test signature and the pick table so you can \emph{say no}: weights
to MST (\(\chref{c17}\)), enumeration to counting (\(\chref{c32}\)), directed flow cuts to
dominators -- and only one dominator preview, kept honest by admitting
Lengauer--Tarjan is not this chapter.

\medskip
The Summit CF~835F (\(\probref{c31-summit}\)) put it all together on a kingdom
\emph{almost} a tree: peel to the unique cycle, hand every cycle vertex to a Ch.~12
two-sweep for its height and internal diameter, break the cycle into a doubled
prefix, and let a Ch.~19 segment-tree monoid answer the best cut in
\(\bigO{k\log k}\).  The greedy `cut the longest cycle edge' fails on the official
sample by~6 -- a bridge-shaped lesson that this chapter exists to teach.

\medskip
Carry three reflexes forward:
\begin{enumerate}[nosep, leftmargin=1.6em]
  \item \emph{Name the redundancy before you touch it.}  Edge, vertex, arrow, cycle --
        the four tests of \(\defref{c31-signature}\) choose the decomposition, and
        choosing wrongly (bridge tree where blocks were needed, or vice-versa)
        loses the equality case and prints a wrong answer.
  \item \emph{Contract, then reuse.}  Once the redundant part is a node, the object
        you DP on is a tree -- Ch.~12 for heights, Ch.~13 for folds, Ch.~19 for windows.
        Nothing new is needed after the extraction; that is the entire point.
  \item \emph{Say `no' when the residual is not a tree.}  Two cycles sharing an edge
        break the ring transfer; a directed flow cut is not a bridge/block question;
        weighted minimisation is MST territory.  The hand-off table
        (\(\defref{c31-pick}\)) is as valuable as any algorithm here.
\end{enumerate}

\medskip
\textbf{Next:} \chref{c32} (Chapter~30, Counting Labelled Trees) removes the input
graph entirely.  You have spent this chapter earning a tree from a graph that was
given; the next asks how many trees meet a specification written on \(n\) labels --
Cayley's \(n^{n-2}\), Pr\"ufer's bijection between trees and length-\((n-2)\) sequences,
the Matrix--Tree theorem for spanning-tree counts of a general graph, and the
counter-intuition of `expected distance along a random walk'.  Where here the
tree was an answer, there the tree is a quantity; the tools shift from DFS to
linear-algebra, and the same Codeforces problem-set that taxed your bridge logic
here will tax your binomial coefficients next.
\end{lastword}


\closechapter
